\documentclass[11pt]{article}

\usepackage[margin=1.1in]{geometry}

\usepackage{amssymb,mathtools,natbib}

% Anchored formal environments for causalsmith papers.
% First mandatory argument = obj_id (machine anchor joining the paper to the
% Lean crosswalk; invisible in the rendered PDF). Optional argument = name.
% Each environment also defines \label{obj:<obj_id>} so prose can use
% \cref{obj:T-1}; cleveref derives the printed kind and number from the target.
\usepackage{amsmath}
\usepackage{mathrsfs}
\usepackage{amsthm}
\usepackage{booktabs}
% Long displays may break across pages; let TeX stretch a little harder before an
% overfull \hbox so wide formulas/tables are less likely to run off the page.
\allowdisplaybreaks
\setlength{\emergencystretch}{3em}
\newtheorem{theorem}{Theorem}
\newtheorem{lemma}{Lemma}
\newtheorem{assumption}{Assumption}
\newtheorem{definition}{Definition}
\newtheorem{proposition}{Proposition}
\newtheorem{remark}{Remark}
\newtheorem{citedresult}{Cited result}
% The amsthm counter is `algorithmbox` (not `algorithm`) so a later `\usepackage{algorithm}` cannot clash.
\newcounter{algorithmbox}
\renewcommand{\thealgorithmbox}{\arabic{algorithmbox}}
\NewDocumentEnvironment{theoremv}{m o s}{\begin{theorem}\label{obj:#1}{\normalfont[\texttt{\detokenize{#1}}]\IfValueT{#2}{\ (#2)}\IfBooleanT{#3}{\verificationfootnotemark}\ }}{\end{theorem}}
\NewDocumentEnvironment{lemmav}{m o s}{\begin{lemma}\label{obj:#1}{\normalfont[\texttt{\detokenize{#1}}]\IfValueT{#2}{\ (#2)}\IfBooleanT{#3}{\verificationfootnotemark}\ }}{\end{lemma}}
\NewDocumentEnvironment{assumptionv}{m o s}{\begin{assumption}\label{obj:#1}{\normalfont[\texttt{\detokenize{#1}}]\IfValueT{#2}{\ (#2)}\IfBooleanT{#3}{\verificationfootnotemark}\ }}{\end{assumption}}
\NewDocumentEnvironment{definitionv}{m o s}{\begin{definition}\label{obj:#1}{\normalfont[\texttt{\detokenize{#1}}]\IfValueT{#2}{\ (#2)}\IfBooleanT{#3}{\verificationfootnotemark}\ }}{\end{definition}}
% A `propositionv` is a proved result tiered as a Proposition (theorem-grade, machine-verified).
\NewDocumentEnvironment{propositionv}{m o s}{\begin{proposition}\label{obj:#1}{\normalfont[\texttt{\detokenize{#1}}]\IfValueT{#2}{\ (#2)}\IfBooleanT{#3}{\verificationfootnotemark}\ }}{\end{proposition}}
% A `remarkv` holds NON-load-bearing interpretive / open-question content (e.g. a causalsmith `oeq:`
% open question). It is neither proved nor assumed; no theorem may depend on it.
\NewDocumentEnvironment{remarkv}{m o s}{\begin{remark}\label{obj:#1}{\normalfont[\texttt{\detokenize{#1}}]\IfValueT{#2}{\ (#2)}\IfBooleanT{#3}{\verificationfootnotemark}\ }}{\end{remark}}
% An `algorithmv` is a DEFINITION-kind object presented as a boxed, numbered-step procedure (an
% estimator / selection rule built in stages). Same anchor contract as `definitionv`; its body is
% ordinary LaTeX whose steps are `\begin{enumerate}` items, so tex2html needs no extra machinery.
% Presented in the CONVENTIONAL algorithm style readers expect from `algorithm`+`algorithmic`:
% a rule above, an "Algorithm N (Title)" caption line, a rule under the caption, the numbered
% steps, and a closing rule. It is NOT a float — the anchor contract requires the object to stay
% where the outline places it, and floats drift. The obj-id tag is suppressed here (it is
% bookkeeping, and a caption line carrying it reads as debug output in a finished paper).
\NewDocumentEnvironment{algorithmv}{m o s}%
  {\par\addvspace{\medskipamount}\noindent\rule{\linewidth}{0.8pt}\par\nobreak\vspace{2pt}%
   \refstepcounter{algorithmbox}\label{obj:#1}%
   \noindent\textbf{Algorithm \thealgorithmbox}\IfValueT{#2}{ \textnormal{(#2)}}%
   \IfBooleanT{#3}{\verificationfootnotemark}\par\nobreak\vspace{2pt}%
   \noindent\rule{\linewidth}{0.4pt}\par\nobreak\vspace{2pt}\normalfont}%
  {\par\nobreak\vspace{2pt}\noindent\rule{\linewidth}{0.8pt}\par\addvspace{\medskipamount}}
% A `citedv` is an IMPORTED external result the paper relies on but does NOT prove
% (a causalsmith `gate_class:"cited"` node). It is machine-checked against the cited
% source at F2.5, never proved here; the fixed provenance note makes that explicit
% so the environment can never be read as a contribution of this paper. The \citep
% attribution is supplied by the surrounding prose (P2), keyed to references.bib.
\NewDocumentEnvironment{citedv}{m o s}{\begin{citedresult}\label{obj:#1}{\normalfont[\texttt{\detokenize{#1}}]\IfValueT{#2}{\ (#2)}\IfBooleanT{#3}{\verificationfootnotemark}\ \textit{(imported result; machine-checked against the cited source, not proved here.)}\ }}{\end{citedresult}}
% \leanref{obj_id}{display text}: an inline first-mention link to a from-note object's
% verified Lean code. In the PDF it renders the display text plainly (the Lean link is a
% web-only affordance); tex2html turns it into a drawer-opening span.
\newcommand{\leanref}[2]{#2}
% For a theorem/lemma whose Lean declaration accepts a source-matched published proposition as an
% input, the starred anchored environment puts the mark in its heading; the generated text remains
% outside the frozen mathematical environment. Keep the combined macro for older paper bundles.
\newcommand{\verificationfootnotemark}{\footnotemark}
\newcommand{\verificationfootnotetext}[1]{\footnotetext{#1}}
\newcommand{\verificationfootnote}[1]{\verificationfootnotemark\verificationfootnotetext{#1}}
% xcolor before hyperref (hyperref would otherwise pull in plain `color` for colorlinks, and a
% later xcolor load clashes). The page-1 identity stamp at the end of this file needs a grey.
\usepackage{xcolor}
\usepackage{graphicx}
\usepackage{eso-pic}
% hyperref follows natbib and the environment macros; cleveref must follow hyperref.
\usepackage[colorlinks=true,linkcolor=blue,citecolor=blue,urlcolor=blue]{hyperref}
\usepackage[capitalise,nameinlink,noabbrev]{cleveref}
% Pin the names explicitly: labels are placed inside the underlying amsthm environments, and these
% declarations make the PDF contract independent of cleveref's theorem-name inference. House style:
% reference names always print with a capital first letter ("Theorem 1", "Definition 2"), in any
% sentence position — hence `capitalise` above and capitalized \crefname forms below.
\crefname{theorem}{Theorem}{Theorems}
\Crefname{theorem}{Theorem}{Theorems}
\crefname{lemma}{Lemma}{Lemmas}
\Crefname{lemma}{Lemma}{Lemmas}
\crefname{assumption}{Assumption}{Assumptions}
\Crefname{assumption}{Assumption}{Assumptions}
\crefname{definition}{Definition}{Definitions}
\Crefname{definition}{Definition}{Definitions}
\crefname{proposition}{Proposition}{Propositions}
\Crefname{proposition}{Proposition}{Propositions}
\crefname{remark}{Remark}{Remarks}
\Crefname{remark}{Remark}{Remarks}
\crefname{citedresult}{Cited result}{Cited results}
\Crefname{citedresult}{Cited result}{Cited results}
\crefname{algorithmbox}{Algorithm}{Algorithms}
\Crefname{algorithmbox}{Algorithm}{Algorithms}
% ---------------------------------------------------------------------------
% Page-1 identity stamp (arXiv style) and the pinned title-block date.
%
% Split of responsibility: THIS file owns the FORMAT (placement, font, colour, punctuation),
% the generated `paper_stamp.tex` owns only the VALUES (number+version, area, revised date,
% DOI, link target). P4 refreshes this template on every emit, so a layout fix reaches every
% bundle from a recompile; and because `paper_stamp.tex` is not one of the P2 assembly
% digest's authored inputs, a DOI that arrives after publication needs only that one small
% file rewritten plus a recompile — never a redraft, never a reassembly.
%
% Defaults first, so a bundle WITHOUT a stamp file compiles exactly as it always did:
% `\cspaperdate` falls back to `\today` and no stamp is drawn.
\providecommand*{\cspaperdate}{\today}  % the \date{} target
\providecommand*{\csstampid}{}          % e.g. CSWP-2026-008v3 (empty ⇒ no stamp at all)
\providecommand*{\csstamparea}{}        % e.g. Stat
\providecommand*{\csstampdate}{}        % e.g. 27 Aug 2026
\providecommand*{\csstampdoi}{}         % e.g. 10.5281/zenodo.123 (empty ⇒ DOI part omitted)
\providecommand*{\csstamplink}{}        % href target (DOI url, else the paper's page; may be empty)
\IfFileExists{paper_stamp.tex}{% GENERATED by CausalSmith P4 — paper identity stamp values. Do not hand-edit.
%
% Field order on the stamp (owned by paper_macros.tex):
%   <number>v<version> · doi:<doi> · [<Area>] <date>   — the DOI is never the last token, so a
%   text extractor cannot run it into whatever follows. Without a DOI that part is omitted.
% Format lives in paper_macros.tex; only the values are here, and this file is NOT one of
% the P2 assembly digest's authored inputs. A DOI arriving after publication therefore needs
% this file rewritten plus a `latexmk` recompile — never a redraft or a reassembly.
\renewcommand*{\csstampid}{CSWP-2026-027v3}
\renewcommand*{\csstamparea}{Stat}
\renewcommand*{\csstampdate}{1 Oct 2026}
\renewcommand*{\csstampdoi}{}
\renewcommand*{\csstamplink}{https://causalsmith.org/papers/stat_mar_nearcomplete_frontier_v1}
\renewcommand*{\cspaperdate}{1 October 2026}
}{}
\makeatletter
% Field order is deliberate: the DOI is NEVER the last token on the line. A trailing identifier
% runs into whatever a text extractor emits next, which makes it unsafe to match mechanically
% (the Zenodo stamp matcher reads exactly this line); the date is a safe terminator.
%   <number>v<version> · doi:<doi> · [<Area>] <date>      — with a DOI
%   <number>v<version> · [<Area>] <date>                  — without
\newcommand{\cs@stampsep}{\,\textperiodcentered\,}
\newcommand{\cs@stamptext}{%
  \csstampid
  \ifx\csstampdoi\@empty\else\cs@stampsep doi:\csstampdoi\fi
  \ifx\csstamparea\@empty\else\cs@stampsep[\csstamparea]\fi
  \ifx\csstampdate\@empty\else\quad\csstampdate\fi
}
\ifx\csstampid\@empty\else
  \definecolor{csstampgrey}{gray}{0.45}
  % Background shipout material: it occupies no space, so the text block and the \thanks
  % footnote are byte-identical to an unstamped compile. The starred form applies to the
  % NEXT page shipped only — i.e. page 1. Placed in the left margin (the text block starts
  % at 1.1in), reading bottom-to-top, in a Times-like face at 8pt.
  \AddToShipoutPictureBG*{%
    \AtPageLowerLeft{%
      \put(\LenToUnit{0.42in},\LenToUnit{1.1in}){%
        \rotatebox{90}{%
          \fontfamily{ptm}\fontsize{8}{10}\selectfont
          \ifx\csstamplink\@empty
            \textcolor{csstampgrey}{\cs@stamptext}%
          \else
            \expandafter\href\expandafter{\csstamplink}{\textcolor{csstampgrey}{\cs@stamptext}}%
          \fi
        }%
      }%
    }%
  }
\fi
\makeatother


\title{Minimax estimation and honest inference in randomized experiments with missing outcomes}

\author{CausalSmith\thanks{Machine-generated by the CausalSmith research pipeline. Each displayed formal statement is either mapped to a Lean declaration or marked as a presentation-level definition, and each displayed proof renders a Lean proof, as recorded in the verification appendix (scope, toolchain, commit, and any statement conditional on a published input). Cited work otherwise supplies framing and positioning.}}

\date{\cspaperdate}

\begin{document}

\maketitle

\begin{abstract}
We characterize optimal estimation and honest inference for the population average treatment effect in a balanced randomized experiment with incomplete primary outcomes. Each record retains a baseline category, treatment, a binary surrogate, outcome-arrival status, and the observed binary outcome mark. Outcome arrival is missing at random conditional on baseline, treatment, and surrogate, with a known cellwise arrival floor \(q\in[1/2,1]\). For sample size \(n\ge1\) and at most \(d\ge1\) baseline categories, the minimax squared-error risk is, up to universal constants,
\[
\frac1n+(1-q)^2\min\left\{1,\left[\frac{d}{n\log(e+n)}\right]^2\right\}.
\]
For each fixed noncoverage level \(\alpha\in(0,1/2)\), the minimax worst-law expected length of honest connected confidence intervals has the square-root order of this benchmark. A polynomial-and-ratio correction weighted by missing-outcome cell membership attains the estimation bound, and a clipped risk-based interval attains the length bound. Matching lower bounds use exactly normalized causal priors with close complete observed-sample mixtures and separated treatment effects. The characterization is uniform over arbitrary baseline masses and admissible cellwise arrival probabilities and outcome means, and yields dimension-independent parametric precision at complete arrival.
\end{abstract}

\section{Introduction}

This paper characterizes how outcome arrival and baseline dimension jointly determine optimal precision for a population average treatment effect in a balanced randomized experiment. The experiment records baseline category, treatment assignment, a binary surrogate measurement, and outcome-arrival status for every unit, together with a binary primary outcome when it arrives. Arrival is missing at random conditional on baseline category, treatment, and surrogate. Under a known cellwise arrival floor between one half and one, we establish matching finite-sample bounds for minimax squared-error risk and for the worst-law expected length of connected confidence intervals with uniform coverage.

The model permits arbitrary baseline masses, including many sparsely occupied categories, and freely varying cellwise arrival probabilities and outcome means subject to its causal and positivity restrictions. Treatment is an independent fair coin, and realized surrogate measurements and outcomes agree with their assigned potential measurements. The population target is the mean difference between the two potential primary outcomes. These conditions define the randomized-MAR law class in \cref{obj:def:law-class}; the corresponding observed-law identification formula is established in \cref{obj:lem:identification-infrastructure}. The resulting uniform estimation problem concerns recovery of an aggregate treatment effect across cells whose individual outcome means may be sparsely observed.

Let \(n\ge1\) denote the sample size, \(d\ge1\) the baseline-support bound, and \(q\in[1/2,1]\) the known outcome-arrival floor. Our principal result, \cref{obj:thm:point-frontier}, establishes that minimax squared-error risk over this class is comparable to
\[
\frac1n+(1-q)^2
\min\left\{1,\left[\frac{d}{n\log(e+n)}\right]^2\right\},
\]
with universal comparison constants uniform over all such \(n,d,q\). For each fixed noncoverage level \(\alpha\in(0,1/2)\), \cref{obj:thm:interval-frontier} establishes that the minimax worst-law expected length among honest connected confidence intervals is comparable to the square root of this expression, with constants depending only on \(\alpha\). Honesty here requires coverage at least \(1-\alpha\) under every law in the specified class at each sample size.

The rate separates sampling uncertainty from the cost of recovering outcome means across many baseline categories. The additional root mean squared error is proportional to the arrival deficit \(1-q\), multiplied by the smaller of one and \(d/[n\log(e+n)]\). Thus improving arrival permits larger baseline supports at sampling-order precision. In particular, an arrival deficit at most \(n^{-1/2}\) yields root mean squared error and optimal honest expected length of order \(n^{-1/2}\), uniformly over the support bound. At complete arrival, \cref{obj:prop:complete-arrival-reduction} gives dimension-independent parametric rates and identifies the estimator as a projected outcome contrast using the known treatment-allocation weights.

The constructive argument starts from a simple decomposition: arrived outcomes estimate the observed part of each cell contribution, while records with missing outcomes still reveal their cell membership and therefore determine the mass requiring correction. The estimator in \cref{obj:def:mm-estimator} combines these two pieces. Separate auxiliary streams select between polynomial correction in lightly populated cells and an empirical ratio in heavily populated cells. Conditional averaging over the auxiliary randomization yields a procedure determined by the observed sample, and \cref{obj:lem:upper-risk} supplies its uniform squared-risk bound. A clipped interval based on this bound satisfies finite-sample uniform coverage and attains the expected-length order, as established in \cref{obj:lem:honest-interval-construction}. The exhibited polynomial tuning begins when \(\log(e+n)\ge128\).

The converse has two branches. When the squared arrival-weighted dimension scale is at least \(n^{-1}\), oppositely oriented pairs of baseline labels and a fixed filler label give exactly normalized priors whose complete observed-sample mixtures are close while their treatment effects concentrate on separated sides of a deterministic center; this is the comparison in \cref{obj:thm:normalized-converse}. When the squared arrival-weighted dimension scale is below \(n^{-1}\), the parametric floors supply the required lower bounds. Testing reductions combine these existential prior comparisons with parametric sampling floors to obtain both minimax lower bounds while preserving the factor \(1-q\).

The closest identification and efficiency benchmark is \citet[Lemma~1.1 under Assumptions~1--3, Theorem~2.1, and Sections~3--4]{KallusMao2025}, which develops surrogate-assisted ATE identification, efficiency, and estimation under its nuisance-convergence conditions. Our analysis characterizes uniform precision over unrestricted finite-cell nuisance functions under balanced randomization. The dimension-dependent complete-record analysis of \citet[Theorems~1--2 and Appendix~C.5]{ZengBalakrishnanHanKennedy2024} provides a related benchmark with unknown categorywise treatment propensities. Polynomial estimation for discrete treatment effects and an independent annotation experiment are developed in \citet{CausalSmith2026,CausalSmith2026Annotation}. The polynomial and moment-matching methods also build on \citet[Sections~III--IV]{JiaoVenkatHanWeissman2015} and \citet[Proposition~3 and Appendix~B.2]{WuYang2016}. Against these precedents, the present contribution is the joint arrival-and-dimension characterization for estimation and honest expected length in the specified randomized-MAR experiment.

The exposition proceeds through related work, the statistical setup and criteria, the minimax estimation and honest-inference results, and their statistical implications, followed by limitations and future work. The appendices develop identification and estimation bounds, the normalized-prior converse, and the testing reductions completing the main comparisons, together with the verification note and proofs of the main results.


\section{Related work}

The closest identification and efficiency analysis is that of \citet[Lemma~1.1 under Assumptions~1--3, Theorem~2.1, and Sections~3--4]{KallusMao2025}. Their experiment observes baseline covariates, treatment, and surrogate outcomes for all units, together with primary outcomes for a labeled subset. Treatment unconfoundedness, missingness at random conditional on covariates, treatment, and surrogates, and the corresponding overlap conditions identify the population average treatment effect. Theorem~2.1 characterizes the semiparametric efficiency bound; Sections~3--4 develop estimation and asymptotic inference under nuisance convergence and moment conditions, including a regime with a dominant unlabeled sample. Surrogates supply predictive information about the primary outcome within this identification framework. The present analysis specializes treatment assignment to balanced randomization and studies uniform performance over finite baseline supports with arbitrary cell masses and unrestricted cellwise arrival probabilities and outcome means, subject to the restrictions in \cref{obj:def:law-class}. This comparison separates identification and regular efficiency from the uniform estimation difficulty generated by many sparsely observed cells.
% [Kallus and Mao](https://academic.oup.com/jrsssb/article/87/2/480/7829031)

A closely related dimension effect appears in \citet[Theorems~1--2 and Appendix~C.5]{ZengBalakrishnanHanKennedy2024}. Their experiment supplies complete covariate--treatment--outcome records and imposes treatment positivity with an unknown categorywise propensity score. Theorem~1 bounds squared error by a sampling term and a dimension term proportional to the squared ratio of alphabet size to sample size. For alphabet sizes at most a constant multiple of sample size times its logarithm, Theorem~2 establishes a lower bound with a squared logarithmic improvement in the dimension term, with constants depending on the fixed positivity level. Appendix~C.5 develops the moment-matching comparison through approximately normalized masses in an auxiliary Poisson experiment and a reduction to probability laws. Here, treatment probabilities are known by design, while outcome arrival varies across observed cells. The arrival-and-dimension characterization in \cref{obj:thm:point-frontier} therefore concerns a different source of incomplete information.
% [Zeng, Balakrishnan, Han, and Kennedy](https://arxiv.org/html/2405.00118v3)

For the complete-record discrete-confounding experiment, \citet[Sharp fixed-interior minimax theorem]{CausalSmith2026} obtains, for each fixed overlap level \(\epsilon\in(0,1/2)\),
\[
\frac1n+\frac{d^2}{n^2(\log n)^2}
\]
up to \(\epsilon\)-dependent constants whenever \(n\ge N_\epsilon\) and \(d\le\rho_\epsilon n\log n\). The same work gives the parametric endpoint rate \(n^{-1}\) under balanced treatment for every positive \(n,d\). These results supply direct precedents for polynomial correction in sparse cells and dimension-independent estimation under complete observation and balanced assignment. In the present experiment, \cref{obj:prop:complete-arrival-reduction} identifies the corresponding complete-arrival reduction, while the main characterization tracks the additional difficulty associated with outcome arrival. In the independent annotation experiment, \citet[Sharp annotation frontier theorem]{CausalSmith2026Annotation} combines \(n\) complete outcome-labeled records with \(m\) independent treatment--covariate records and obtains
\[
\min\left\{1,\frac1n+\frac{d^2}{(n+m)^2\log^2(en)}\right\}
\]
for fixed interior overlap, all \(n\ge1\), \(m\ge0\), and \(d\ge2\), with constants depending on the overlap level. The current experiment instead observes arrival indicators and surrogate measurements within one randomized sample, with cellwise arrival probabilities bounded below by \(q\); its extra term is therefore weighted by \((1-q)^2\) and uses the within-sample arrival information.
% [CausalSmith: fixed-overlap estimation](https://causalsmith.org/papers/stat_discrete_ate_minimax_loggap_polynomial_upper_match/), [CausalSmith: independent annotation](https://causalsmith.org/papers/stat_semisupervised_discrete_ate_annotation_frontier_v1/slides/)

Structured decaying-labeling theory provides a complementary benchmark. \citet[Theorem~2.2 and equation~(2.4)]{ZhangChakraborttyBradic2025} establishes asymptotic normality at an effective sample size determined by the inverse moment of the treatment-and-labeling propensity. Its conditions include correct specification of both nuisance limits, a product-rate restriction on nuisance estimation errors, diverging effective sample size, and the stated moment and tail conditions. Their Theorems~4.1--4.2 develop bias-reducing procedures under offset labeling models, design and noise regularity, and alternative sparsity restrictions. The decoupled procedure additionally controls estimation of the treatment propensity; centered asymptotic normality retains the required control of the product of nuisance approximation errors. These results accommodate declining labeling probabilities through structured nuisance estimation. The present uniform analysis concerns arrival floors between one half and one and freely varying finite-cell nuisance functions. Its dimension parameter measures the number of baseline categories, whereas the sparsity conditions in those structured procedures regulate the complexity of nuisance representations.
% [Zhang, Chakrabortty, and Bradic](https://arxiv.org/html/2305.12789)

Uniform inference with high-dimensional missing data provides an earlier point of contact. \citet[pp.~285--319]{Robins1997} studies how selection probabilities and growing nuisance complexity can obstruct uniform inference in missing-data experiments. Higher-order influence-function analyses subsequently derive minimax conclusions under smoothness, density, and related regularity restrictions \citep[pp.~335--421]{Robins2008}; \citet[pp.~1305--1321]{Robins2009} and \citet[pp.~1951--1987]{Robins2017} develop related minimax results. The present experiment replaces those smooth-function restrictions with a finite but unrestricted collection of cell masses, arrival probabilities, and outcome means. Its contribution is an exact finite-sample arrival-and-dimension rate and a normalized-prior comparison on the complete observed-record space, rather than a higher-order expansion over a smooth nuisance class.

The estimation and converse methods have substantial precedents in discrete functional estimation. \citet[Sections~III--IV]{JiaoVenkatHanWeissman2015} combines polynomial approximation with unbiased estimation of polynomial terms and develops lower bounds through prior comparisons. \citet[Proposition~3 and Appendix~B.2]{WuYang2016} uses moment matching and comparison of Poisson mixtures to establish the dimension-dependent entropy lower bound. These methods explain how aggregate functional estimation can exploit sparse categories even when individual cell parameters are difficult to estimate. Their role here is to guide the polynomial correction for missing outcomes and the construction of statistically close alternatives with separated treatment effects.
% [Jiao, Venkat, Han, and Weissman](https://arxiv.org/pdf/1406.6956), [Wu and Yang](https://arxiv.org/pdf/1407.0381)

The expected-length criterion places inference beside the theory of honest confidence intervals. \citet[Theorem~1 and Section~5]{CaiLow2004} relates uniform coverage and expected length to statistical distinguishability for linear functionals in Gaussian experiments, including nonconvex parameter spaces formed from finite unions of convex classes. In causal inference, \citet{Armstrong2021} studies average treatment effects conditional on realized treatments and covariates under smoothness or shape restrictions: finite-sample optimality uses Gaussian errors with known variance, and feasible procedures with unknown error distributions receive asymptotic guarantees. For randomized trials with missing outcomes, \citet[Theorem~2 under Conditions~2--3]{DiazVanDerLaan2017} establishes doubly robust asymptotic normality under a Donsker condition and nuisance convergence faster than the inverse fourth root of sample size, with at least one correct nuisance limit. The criterion in \cref{obj:thm:interval-frontier} is finite-sample uniform coverage together with worst-law expected length over the unrestricted finite-cell class.
% [Cai and Low](https://arxiv.org/pdf/math/0503662), [Armstrong and Kolesár](https://onlinelibrary.wiley.com/doi/abs/10.3982/ECTA16907), [Díaz and van der Laan](https://arxiv.org/html/1704.01538)

Against these benchmarks, the contribution is the joint characterization of outcome arrival and baseline dimension for minimax squared error and honest expected length in \cref{obj:thm:point-frontier,obj:thm:interval-frontier}. The converse in \cref{obj:thm:normalized-converse} constructs exactly normalized causal laws and compares the induced mixtures of the complete observed record, including baseline category, treatment, surrogate, arrival indicator, and observed outcome mark. This experiment-specific comparison supports both statistical criteria while retaining the information carried by every observed coordinate.


\section{Setup and statistical criteria}

Throughout, \leanref{sym:n}{\ensuremath{n}} denotes the positive integer sample size, \leanref{sym:d}{\ensuremath{d}} the positive integer bound on the baseline alphabet, and \leanref{sym:q}{\ensuremath{q\in[1/2,1]}} the known outcome-arrival floor. Generic constants \(c,C>0\) may change across displays; subscripts indicate permitted dependence on fixed parameters. For nonnegative quantities, \(u\lesssim v\) means \(u\le Cv\) with the stated constant dependence, and \(u\asymp v\) means both \(u\lesssim v\) and \(v\lesssim u\). We use natural logarithms and write \(\|f\|_{p,Q}=(\mathbb E_Q|f|^p)^{1/p}\) for \(1\le p<\infty\). Baseline indices range over \(x=0,\ldots,d-1\), binary indices over \(a,s\in\{0,1\}\), and sample indices over \(i=1,\ldots,n\).

The statistical experiment observes \((X,A,S,R,RY)\): \(X\) is the baseline category, \(A\) treatment, \(S\) the realized surrogate, \(R\) outcome-arrival status, and \(RY\) the observed outcome mark. The argument below uses this observed experiment, the law class, the treatment-effect target, and the two minimax criteria. The next three definitions record the finite coordinate representation and observation map needed to connect those statistical objects to the potential-outcome assumptions.

\begin{definitionv}{synth_3}[Baseline and potential coordinates]
A potential-data atom consists of a baseline category \(X\), binary potential surrogate measurements \(S(0),S(1)\), and binary potential outcomes \(Y(0),Y(1)\). For a nonnegative integer \(d\), its coordinates range over
\[
X\in\{0,\ldots,d-1\},
\qquad
(S(0),S(1),Y(0),Y(1))\in\{0,1\}^{4},
\]
where the baseline domain is empty when \(d=0\). The collection of all such atoms is finite and has decidable equality.
\end{definitionv}

\begin{definitionv}{synth_4}[Full-data atom coordinates]
A full-data atom, parameterized by a nonnegative integer \(d\), consists of the baseline category \(X\), the potential surrogates \(S(0),S(1)\), the potential outcomes \(Y(0),Y(1)\), the treatment indicator \(A\), the realized surrogate \(S\), the realized outcome \(Y\), and the outcome-arrival indicator \(R\). Its coordinates range over
\[
\bigl(X,S(0),S(1),Y(0),Y(1),A,S,Y,R\bigr)
\in \{0,\ldots,d-1\}\times\{0,1\}^{8},
\]
where the baseline set is empty when \(d=0\). The realized coordinates \(S,Y\) are separate from the potential coordinates; every tuple in the displayed Cartesian product is a full-data atom.
\end{definitionv}

\begin{definitionv}{synth_2}[Observed records and space]
The observed-record space \(O\) consists of tuples
\[
O=(X,A,S,R,RY)\in \{0,\ldots,d-1\}\times\{0,1\}^{4},
\]
where \(d\) is a nonnegative integer and the baseline-label set is empty when \(d=0\). The coordinates are the baseline category \(X\), treatment indicator \(A\), realized surrogate \(S\), outcome-arrival indicator \(R\), and observed outcome mark \(RY\). Each of the last four coordinates is binary. The record space is finite and is equipped with the discrete \(\sigma\)-algebra, so every subset is measurable.
\end{definitionv}

The potential coordinates describe measurements under control and treatment \citep{Rubin1974}; consistency restricts realized measurements to the assigned arm. The observed record \leanref{sym:O}{\ensuremath{O}} is obtained from a full-data atom by the observation map
\[
\bigl(X,S(0),S(1),Y(0),Y(1),A,S,Y,R\bigr)
\longmapsto (X,A,S,R,RY).
\]
Given the full-data probability law \leanref{sym:P}{\ensuremath{P}}, specified in \cref{obj:synth_1}, let \(P_O\) denote its image under this map. When \(R=0\), the baseline category, treatment assignment, surrogate, and arrival indicator remain observed, while the outcome mark \(RY\) equals zero. Retaining \(R\) distinguishes a missing outcome from an arrived outcome equal to zero. When \(R=1\), the outcome mark equals \(Y\). Thus the induced law is supported on records whose fifth coordinate is zero whenever their fourth coordinate is zero, within the Cartesian record space of \cref{obj:synth_2}. We use \(O^n\) for the Cartesian space of \(n\) records and \(P_O^{\otimes n}\) for their independent product law.

The statistical model combines independent sampling with restrictions on treatment assignment, realized measurements, and outcome arrival. The sampling condition is stated directly for the observed records.

\begin{assumptionv}{ass:iid}[Independent observed sampling]
The observed records \(O_1,\ldots,O_n\), with sample size \(n\), are independent and identically distributed under the observed-data law \(P_O\):
\[
O_1,\ldots,O_n\stackrel{\mathrm{iid}}{\sim}P_O.
\]
\end{assumptionv}

Independent superpopulation sampling makes \(P_O^{\otimes n}\) the statistical experiment and supplies the common population law underlying the risk and coverage criteria \citep{vanDerVaart1998}. Treatment assignment obeys two complementary design restrictions: independence from pretreatment and potential measurements, and a known probability of assignment to each arm.

\begin{assumptionv}{ass:randomized-independence}[Randomized treatment independence]
The treatment indicator \(A\) is independent of the baseline category \(X\), the potential surrogates \(S(0),S(1)\), and the potential outcomes \(Y(0),Y(1)\):
\[
A\perp \bigl(X,S(0),S(1),Y(0),Y(1)\bigr).
\]
\end{assumptionv}

Randomized treatment assignment gives the treatment and control arms the same population distribution of baseline and potential measurements; assignment by an independent random draw satisfies this condition by design \citep{ImbensRubin2015}. We specialize its assignment probability as follows.

\begin{assumptionv}{ass:balanced-randomization}[Balanced treatment assignment]
Under the full-data law \(P\), the treatment indicator \(A\) satisfies
\[
P(A=1)=\frac12.
\]
\end{assumptionv}

Balanced Bernoulli randomization fixes both arm probabilities at one half, while allowing the realized treatment counts to vary across samples \citep{ImbensRubin2015}. To connect the assignment mechanism to the measured surrogate and primary outcome, we impose consistency.

\begin{assumptionv}{ass:consistency}[Consistency of realized measurements]
The realized surrogate \(S\) and outcome \(Y\) equal the potential measurements selected by the treatment indicator \(A\):
\[
(S,Y)=(S(A),Y(A)).
\]
\end{assumptionv}

Consistency of observed and potential measurements ensures that each arm reveals the measurements corresponding to its assigned treatment \citep{KallusMao2025}. Outcome arrival may depend on that treatment and on the realized surrogate. Its relationship to the primary outcome is governed by conditional independence.

\begin{assumptionv}{ass:mar}[Conditionally independent outcome arrival]
The outcome-arrival indicator \(R\) is conditionally independent of the realized outcome \(Y\), given the baseline category \(X\), treatment indicator \(A\), and realized surrogate \(S\):
\[
R\perp Y\mid(X,A,S).
\]
\end{assumptionv}

Surrogate-assisted missing at random makes arrived outcomes representative of their baseline--treatment--surrogate cell \citep{KallusMao2025}. The condition is appropriate, for example, when primary-outcome collection is randomized using information already recorded in those coordinates. It permits arrival probabilities to differ across cells and uses the surrogate as an observed conditioning variable.

We now specify the full-data law and the cellwise quantities that connect these conditions to the population target.

\begin{definitionv}{synth_1}[Full-data probability law]
The full-data law \(P\) is a probability mass function on full-data atoms of covariate dimension \(d\), where \(d\) is a nonnegative integer. Each atom consists of a potential-data record and separate binary coordinates \(A,S,Y,R\), representing treatment, the realized surrogate, the realized outcome, and outcome arrival, respectively.
\end{definitionv}

For this law, the population average treatment effect \leanref{sym:tau}{\ensuremath{\tau(P)}} is
\[
\tau(P)=\mathbb E_P\!\left[Y(1)-Y(0)\right].
\]
Binary potential outcomes place the target in \([-1,1]\). The observed law \(P_O\) is the image of \(P\) under the observation map above. Expectations of procedures based on the observed sample are consequently taken under \(P_O^{\otimes n}\), even when written as \(\mathbb E_P\).

Identification also requires arrived outcomes in every occupied conditioning cell. The next definition fixes the conditional quantities and their values on null cells before we impose an arrival floor.

\begin{definitionv}{synth_5}[Cellwise arrival probability $\rho_P$ and outcome mean $\mu_P$]
In the setting of \cref{obj:ass:iid,obj:ass:randomized-independence,obj:ass:balanced-randomization,obj:ass:consistency,obj:ass:mar,obj:ass:arrival-floor,obj:def:law-class,obj:def:identified-functional,obj:def:rate,obj:def:point-risk,obj:def:interval-class,obj:def:length-risk,obj:def:mm-estimator,obj:def:mm-interval,obj:def:paired-prior-handle,obj:lem:identification-infrastructure,obj:lem:upper-risk,obj:lem:parametric-floor-infrastructure,obj:lem:prior-reduction,obj:lem:honest-interval-construction,obj:thm:point-frontier,obj:thm:interval-frontier,obj:prop:complete-arrival-reduction,obj:thm:normalized-converse,obj:synth_1,obj:synth_2,obj:synth_3,obj:synth_4,obj:synth_6}, define the cellwise arrival probability and arrived-outcome conditional mean by
\[
\rho_P(x,a,s)=
\begin{cases}
P(R=1\mid X=x,A=a,S=s),
& P(X=x,A=a,S=s)>0,\\
0,
& P(X=x,A=a,S=s)=0,
\end{cases}
\]
and
\[
\mu_P(x,a,s)=
\begin{cases}
\mathbb E_P[Y\mid X=x,A=a,S=s,R=1],
& P(X=x,A=a,S=s,R=1)>0,\\
0,
& P(X=x,A=a,S=s,R=1)=0.
\end{cases}
\]
Here \(P\) is a full-data law, \(X\in\{0,\ldots,d-1\}\) is the baseline cell, \(A\in\{0,1\}\) is the treatment assignment, \(S\in\{0,1\}\) is the observed surrogate, \(R\in\{0,1\}\) is the outcome-arrival indicator, and \(Y\in\{0,1\}\) is the realized primary outcome. The definitions apply to every \((x,a,s)\in\{0,\ldots,d-1\}\times\{0,1\}^2\).
\end{definitionv}

The cellwise arrival probability \leanref{sym:rho}{\ensuremath{\rho_P(x,a,s)}} and arrived-outcome mean \leanref{sym:mu}{\ensuremath{\mu_P(x,a,s)}} are therefore defined for every cell. Their zero conventions assign values to conditional quantities at null events. The positivity restriction concerns cells with positive probability.

\begin{assumptionv}{ass:arrival-floor}[Cellwise outcome-arrival floor]
For every cell \((x,a,s)\) with positive probability under the full-data law \(P\), the cellwise arrival probability \(\rho_P(x,a,s)\) is at least the arrival floor \(q\):
\[
P(X=x,A=a,S=s)>0
\quad\Longrightarrow\quad
\rho_P(x,a,s)\ge q.
\]
\end{assumptionv}

The known uniform MAR positivity floor is the finite-cell specialization of the outcome-observation overlap condition \citep{KallusMao2025}. It ensures positive arrived mass in each occupied cell and quantifies how much primary-outcome information is retained there. The floor is a bound on the cellwise probabilities; the probabilities themselves may vary across cells. Combining it with the design and missingness conditions gives the law class over which statistical performance is evaluated.

\begin{definitionv}{def:law-class}[Law class \(\mathcal M(d,q)\)]
\[
\mathcal M(d,q)
=
\left\{
P:
|\operatorname{supp}_P(X)|\le d,\;
P\text{ satisfies }
\cref{obj:ass:randomized-independence,obj:ass:balanced-randomization,obj:ass:consistency,obj:ass:mar,obj:ass:arrival-floor}
\right\}.
\]
Here \(d\) is a positive integer, \(q\in[1/2,1]\) is the known cellwise arrival floor, and \(\operatorname{supp}_P(X)\) is the set of baseline labels with positive probability under \(P\).
\end{definitionv}

The law class \leanref{sym:M_dq}{\ensuremath{\mathcal M(d,q)}} allows arbitrary baseline masses and arbitrary binary potential-measurement distributions subject to the stated restrictions. In particular, occupied baseline labels may have small probability. At \(q=1\), every occupied cell has complete outcome arrival. The coordinate definitions also specify an empty-alphabet convention at \(d=0\); the minimax analysis uses positive \(d\).

Within this class, the population treatment effect can be expressed entirely through the observed law. The identifying functional weights arrived-outcome means by each arm's distribution of baseline and surrogate measurements.

\begin{definitionv}{def:identified-functional}[Observed-data identifying functional]
The observed-data functional is
\[
\Psi(P_O)
=
\sum_{x=0}^{d-1}\sum_{s=0}^{1}
\left\{
\frac{P_O(X=x,A=1,S=s)}{P_O(A=1)}\mu_P(x,1,s)
-
\frac{P_O(X=x,A=0,S=s)}{P_O(A=0)}\mu_P(x,0,s)
\right\},
\]
for any nonnegative integer \(d\) and any observed-record probability law \(P_O\) on
\[
O=\{0,\ldots,d-1\}\times\{0,1\}^{4},
\]
with coordinates \((X,A,S,R,RY)\). The baseline-label set and the outer sum are empty when \(d=0\). The arrived-outcome mean is defined by
\[
\mu_P(x,a,s)
=
\frac{P_O(X=x,A=a,S=s,R=1,RY=1)}
     {P_O(X=x,A=a,S=s,R=1)},
\qquad a,s\in\{0,1\}.
\]
Every ratio with zero denominator is assigned the value zero. Thus each weighted contribution is zero whenever \(P_O(X=x,A=a,S=s)=0\).
\end{definitionv}

The observed-data functional \leanref{sym:Psi}{\ensuremath{\Psi(P_O)}} identifies \(\tau(P)\) for \(P\in\mathcal M(d,q)\), as established in \cref{obj:lem:identification-infrastructure}. Balanced randomization gives positive arm probabilities, and the arrival floor gives positive arrived mass in every occupied cell. The zero-denominator convention handles the remaining null cells. The resulting finite-cell expression corresponds to the nested-regression identification formula of \href{https://academic.oup.com/jrsssb/article/87/2/480/7829031}{\citet[Lemma~1.1, equation~(10), under Assumptions~1--3]{KallusMao2025}}, specialized here to balanced randomized assignment, binary measurements, and finite baseline support.

Having fixed the target and observed experiment, we evaluate point estimation by worst-law squared error. Procedures are defined on the entire Cartesian sample space and take values in the treatment-effect range.

\begin{definitionv}{synth_6}[Measurable estimator class $\mathcal T_n$]
For integers \(n,d\ge1\), define the measurable estimator class by
\[
\mathcal T_n
=
\left\{
T:\bigl(\{0,\ldots,d-1\}\times\{0,1\}^4\bigr)^n
\longrightarrow[-1,1]
:\ T\text{ is measurable}
\right\}.
\]
Here the observed record \(O=(X,A,S,R,RY)\) takes values in the Cartesian observed-record space \(\{0,\ldots,d-1\}\times\{0,1\}^4\), with the baseline-label convention of \cref{obj:synth_2}. Thus each \(T\in\mathcal T_n\) is defined on every sample in this Cartesian space and takes values in the treatment-effect range \([-1,1]\), in the setting of \cref{obj:ass:iid,obj:ass:randomized-independence,obj:ass:balanced-randomization,obj:ass:consistency,obj:ass:mar,obj:ass:arrival-floor,obj:def:law-class,obj:def:identified-functional,obj:def:rate,obj:def:point-risk,obj:def:interval-class,obj:def:length-risk,obj:def:mm-estimator,obj:def:mm-interval,obj:def:paired-prior-handle,obj:lem:identification-infrastructure,obj:lem:upper-risk,obj:lem:parametric-floor-infrastructure,obj:lem:prior-reduction,obj:lem:honest-interval-construction,obj:thm:point-frontier,obj:thm:interval-frontier,obj:prop:complete-arrival-reduction,obj:thm:normalized-converse,obj:synth_1,obj:synth_3,obj:synth_4,obj:synth_5}.
\end{definitionv}

The estimator class \leanref{sym:T_n}{\ensuremath{\mathcal T_n}} permits any measurable use of the retained information. Since the sample space is finite with its discrete \(\sigma\)-algebra, every map into \([-1,1]\) is measurable. Taking the worst risk over the law class and then optimizing over procedures defines the point-estimation criterion.

\begin{definitionv}{def:point-risk}[Minimax squared-error risk \(\mathfrak R^*_{n,d,q}\)]
\[
\mathfrak R^*_{n,d,q}
=
\inf_{T\in\mathcal T_n}
\sup_{P\in\mathcal M(d,q)}
\mathbb E_P
\left[
\left\{
T(O_1,\ldots,O_n)-\tau(P)
\right\}^2
\right].
\]
\end{definitionv}

The minimax squared-error risk \leanref{sym:Rstar}{\ensuremath{\mathfrak R^*_{n,d,q}}} measures the smallest worst-law mean squared error attainable from \(n\) observed records. The order of optimization requires one procedure to perform across the entire law class at the specified sample size, support bound, and arrival floor.

For inference, fix a noncoverage level \leanref{sym:alpha}{\ensuremath{\alpha\in(0,1/2)}}. We require connected intervals with coverage at least \(1-\alpha\) uniformly over the same class.

\begin{definitionv}{def:interval-class}[Honest connected confidence intervals]
The class \(\mathcal C_\alpha(n,d,q)\), for nonnegative integers \(n,d\) and real parameters \(q,\alpha\), is defined using the observed-record space \(O\) with the baseline-label convention of \cref{obj:synth_2}. It consists of procedures \(I(o)=[L(o),U(o)]\) with measurable endpoints satisfying
\[
-1\le L(o)\le U(o)\le1
\qquad\text{for every }o\in O^n,
\]
and uniform coverage:
\[
\mathcal C_\alpha(n,d,q)
=
\left\{
I:
P_O^{\otimes n}\!\left\{
o\in O^n:\tau(P)\in[L(o),U(o)]
\right\}\ge1-\alpha
\quad\text{for every }P\in\mathcal M(d,q)
\right\}.
\]
Here \(\mathcal M(d,q)\) is the law class in \cref{obj:def:law-class}, and \(\tau(P)=\mathbb E_P[Y(1)-Y(0)]\) is the superpopulation average treatment effect. For \(q\in[1/2,1]\) and \(P\in\mathcal M(d,q)\), \cref{obj:lem:identification-infrastructure} gives \(\tau(P)=\Psi(P_O)\), with \(\Psi\) defined in \cref{obj:def:identified-functional}.
\end{definitionv}

The honest interval class \leanref{sym:C_alpha}{\ensuremath{\mathcal C_\alpha(n,d,q)}} imposes a finite-sample coverage requirement for every law in \(\mathcal M(d,q)\). Its endpoints define a closed connected subset of \([-1,1]\) on every possible sample. The fixed interval \([-1,1]\) belongs to the class, so the coverage requirement is feasible. Among honest procedures, expected length measures precision.

\begin{definitionv}{def:length-risk}[Minimax expected length \(\mathfrak L^*_{n,d,q,\alpha}\)]
\[
\mathfrak L^*_{n,d,q,\alpha}
=
\inf_{I\in\mathcal C_\alpha(n,d,q)}
\sup_{P\in\mathcal M(d,q)}
\mathbb E_P
\left[
\operatorname{length}\{I(O^n)\}
\right].
\]
For \(I(O^n)=[L,U]\), its length is \(U-L\).
\end{definitionv}

The minimax expected length \leanref{sym:Lstar}{\ensuremath{\mathfrak L^*_{n,d,q,\alpha}}} is the smallest worst-law expected interval length compatible with uniform coverage. Both criteria use the full observed record: missing-outcome observations contribute their baseline category, treatment, surrogate, and arrival status to any admissible procedure.

A common benchmark will express the dependence of these criteria on sample size, baseline dimension, and outcome arrival. Define the logarithmic sample-size factor \leanref{sym:ell}{\ensuremath{\ell_n}} and the arrival-deficit-weighted dimension scale \leanref{sym:g_ndq}{\ensuremath{g_{n,d,q}}} by
\[
\ell_n=\log(e+n),
\qquad
g_{n,d,q}=(1-q)\min\left\{1,\frac{d}{n\ell_n}\right\}.
\]
These quantities give the following squared-error benchmark.

\begin{definitionv}{def:rate}[Squared-error benchmark \(r_{n,d,q}\)]
\[
r_{n,d,q}
=
\frac{1}{n}
+
(1-q)^2
\min\left\{
1,
\left[\frac{d}{n\log(e+n)}\right]^2
\right\}
=
\frac{1}{n}+g_{n,d,q}^2.
\]
\end{definitionv}

The benchmark \leanref{sym:r_ndq}{\ensuremath{r_{n,d,q}}} combines a sampling term with a dimension term weighted by the outcome-arrival deficit. Its dimension component is bounded by \((1-q)^2\) and vanishes at complete arrival. The comparisons in \cref{obj:thm:point-frontier,obj:thm:interval-frontier} relate minimax squared error to \(r_{n,d,q}\) and minimax honest expected length to its square root, respectively.


\section{Minimax estimation and honest inference}

The benchmark in \cref{obj:def:rate} is attained by splitting each cell's treatment-effect contribution into information from arrived outcomes and a correction for its missing-outcome mass. Arrival reveals the outcome needed for the first piece; a missing record still reveals \((X,A,S)\), which supplies the weight for the second. Under \cref{obj:ass:mar}, the arrived-outcome mean identifies the mean used in both pieces. The moment estimator \leanref{sym:tauhat_MM}{\ensuremath{\widehat\tau_{\mathrm{MM}}}} uses separate streams to estimate them and to choose between polynomial and ratio corrections. Its final conditional average produces a procedure determined by the observed sample.

We first fix the operator and polynomial conventions used in the construction. Projection onto the treatment-effect range is
\(\Pi_{[-1,1]}(t)=\max\{-1,\min\{1,t\}\}\).
The degree-\(k\) Chebyshev polynomial \leanref{sym:Cheb}{\ensuremath{T_k}} has convention
\(T_k(\cos\theta)=\cos(k\theta)\).
For a cell count \(C_j\) and a positive integer \(v\), the falling factorial
\((C_j-1)_{v-1}\) denotes \(\prod_{t=1}^{v-1}(C_j-t)\), with an empty product equal to one. These conventions specify the projection and polynomial correction in the following construction.

\begin{algorithmv}{def:mm-estimator}[Moment estimator $\widehat\tau_{\mathrm{MM}}$]
Given observed records \(O_1,\ldots,O_n\), sample size \(n\), support bound \(d\), and arrival floor \(q\), define \(\widehat\tau_{\mathrm{MM}}\) using the tuning parameters \(m,k,B\) as follows.
\begin{enumerate}
\item Set
\[
\ell_n=\log(e+n),\qquad
m=\frac n8,\qquad
k=\max\{1,\lfloor\ell_n/64\rfloor\},\qquad
B=256\ell_n.
\]
Write the supplied fifth coordinate of each observed record as
\[
Y_i^{\mathrm{obs}}=(RY)_i.
\]
On the support of the observation map, this coordinate equals \(R_iY_i\) and is zero when the outcome is missing.

\item Draw
\[
N\sim\operatorname{Poisson}(n/2).
\]
If \(N>n\), set \(\widetilde\tau=0\). If \(N\le n\), independently assign each of the first \(N\) records a uniform stream indicator and define the stream index sets by
\[
Z_i\in\{1,2,3,4\},\qquad
\Pr(Z_i=t)=\frac14,\qquad
\mathcal I_t=\{i\in\{1,\ldots,N\}:Z_i=t\}.
\]
For this branch, perform the following three steps.

\item For each cell \(j=(x,a,s)\), write \(a_j=a\) and define
\[
M_j^{(2)}
=\sum_{i\in\mathcal I_2}
\mathbf 1\{(X_i,A_i,S_i)=j,\ R_i=0\},
\qquad
V_j=\frac{M_j^{(2)}}m,
\]
\[
C'_j
=\sum_{i\in\mathcal I_3}
\mathbf 1\{(X_i,A_i,S_i)=j,\ R_i=1\},
\qquad
C_j
=\sum_{i\in\mathcal I_4}
\mathbf 1\{(X_i,A_i,S_i)=j,\ R_i=1\},
\]
\[
U_j
=\sum_{i\in\mathcal I_4}
\mathbf 1\{(X_i,A_i,S_i)=j,\ Y_i^{\mathrm{obs}}=1\}.
\]
Thus \(U_j\) is the arrived-success count on the support of the observation map.

\item Let \(T_k\) be the degree-\(k\) Chebyshev polynomial, with convention
\[
T_k(\cos\theta)=\cos(k\theta).
\]
Define
\[
Q_k(z)=
\begin{cases}
\dfrac{1-T_k(1-2z/B)}{2k^2z/B},&z\ne0,\\[6pt]
1,&z=0,
\end{cases}
\qquad
1-Q_k(z)=\sum_{v=1}^{k-1}a_vz^v.
\]
Using the falling-factorial convention
\[
(C_j-1)_{v-1}=\prod_{t=1}^{v-1}(C_j-t),
\]
where an empty product equals one, set
\[
H_j=\sum_{v=1}^{k-1}a_v(U_j-C_j/2)(C_j-1)_{v-1},
\qquad
D_j=
\begin{cases}
U_j/C_j-1/2,&C_j>0,\\
0,&C_j=0,
\end{cases}
\]
\[
G_j
=H_j\mathbf 1\{C'_j\le B/4\}
+D_j\mathbf 1\{C'_j>B/4\}.
\]

\item Define the interval projection and randomized estimate by
\[
\Pi_{[-1,1]}(t)=\max\{-1,\min\{1,t\}\},
\]
\[
\widetilde\tau
=\Pi_{[-1,1]}\!\left[
\frac2m\sum_{i\in\mathcal I_1}(2A_i-1)R_i
\left(Y_i^{\mathrm{obs}}-\frac12\right)
+2\sum_j(2a_j-1)V_jG_j
\right].
\]

\item Output
\[
\widehat\tau_{\mathrm{MM}}
=
\begin{cases}
\displaystyle
\Pi_{[-1,1]}\!\left[
\frac2n\sum_{i=1}^n(2A_i-1)Y_i^{\mathrm{obs}}
\right],
&q=1,\\[8pt]
\displaystyle
\Pi_{[-1,1]}\!\left[
\frac2n\sum_{i=1}^n(2A_i-1)R_i
\left(Y_i^{\mathrm{obs}}-\frac12\right)
\right],
&q\ne1\ \text{and}\ (\ell_n<128\ \text{or}\ d\ge n\ell_n),\\[8pt]
\displaystyle
\mathbb E[\widetilde\tau\mid O_1,\ldots,O_n],
&q\ne1,\ \ell_n\ge128,\ d<n\ell_n.
\end{cases}
\]
The conditional expectation averages over the auxiliary Poisson draw and stream assignments.
\end{enumerate}
\end{algorithmv}

The polynomial-and-ratio branch is used when \(\ell_n=\log(e+n)\ge128\), equivalently \(n\ge e^{128}-e\), together with \(d<n\ell_n\) and \(q\ne1\); the displayed fallback branch covers the complementary regime. These are the exhibited finite-sample tuning choices supporting the order guarantee.

Centering arrived outcomes at one half aligns the first-stream contribution with the centered cell means used in the correction. The arrival indicator ensures that the centered contribution is drawn from observed outcomes. Missing records supply their cell membership through the second-stream statistic \(V_j\), which weights the correction by missing-outcome mass. Under the missing-at-random condition in \cref{obj:ass:mar}, the arrived-outcome mean identifies the outcome mean needed for this missing component.

The third-stream count selects how the fourth-stream data enter the correction. Lightly populated cells use the polynomial statistic \(H_j\), whose coefficients and falling-factorial terms implement the moment correction. Heavily populated cells use the centered empirical ratio \(D_j\). Separating the selection count from the correction data supports the risk analysis in \cref{obj:lem:upper-risk}. Conditional averaging then integrates out the auxiliary Poisson draw and stream assignments; convexity of squared loss preserves the risk bound under this averaging. The complete-arrival and fallback branches specify the estimator across all sample sizes and support bounds.

The uniform squared-risk bound also supplies a finite-sample confidence radius. Let \leanref{sym:C_U}{\ensuremath{C_U}} be a universal upper-risk constant, so that the worst-law squared risk of the estimator is at most \(C_Ur_{n,d,q}\), as established in \cref{obj:lem:upper-risk}. The confidence interval \leanref{sym:Ihat_MM}{\ensuremath{\widehat I_{\mathrm{MM},\alpha}}} uses the clipped radius \leanref{sym:h_alpha}{\ensuremath{h_{n,d,q,\alpha}}} defined below. The construction attains the minimax expected-length order for fixed \(\alpha\); with the exhibited fallback contribution \(C_{\mathrm{fall}}=4e^{128}\) and resulting choice of \(C_U\), its interval equals \([-1,1]\) whenever \(\ell_n<128\).

\begin{definitionv}{def:mm-interval}[Confidence interval $\widehat I_{\mathrm{MM},\alpha}$]
Define
\[
\widehat I_{\mathrm{MM},\alpha}(O^n)
=
\left[
\max\{-1,\widehat\tau_{\mathrm{MM}}-h_{n,d,q,\alpha}\},
\min\{1,\widehat\tau_{\mathrm{MM}}+h_{n,d,q,\alpha}\}
\right],
\]
where the clipped confidence radius is
\[
h_{n,d,q,\alpha}
=\min\left\{2,\sqrt{\frac{C_Ur_{n,d,q}}{\alpha}}\right\},
\]
and \(C_U\) is any universal constant for which the upper squared-risk bound for \(\widehat\tau_{\mathrm{MM}}\) holds.
\end{definitionv}

For a radius below two, Markov's inequality bounds the probability that the estimation error exceeds the radius by \(\alpha\), uniformly over \(\mathcal M(d,q)\). A radius of two yields the entire treatment-effect range. Clipping the endpoints preserves coverage because the target belongs to \([-1,1]\). Consequently, the interval belongs to the honest class in \cref{obj:def:interval-class}, and its length is at most \(2h_{n,d,q,\alpha}\). The coverage and length conclusions are recorded in \cref{obj:lem:honest-interval-construction}.

The attainable risk and interval length match lower bounds over the same experiment and law class. The following two results characterize the point-estimation and honest-inference criteria together.

\begin{theoremv}{thm:point-frontier}[Finite-sample minimax risk frontier]
There exist universal numerical constants \(0<c<C<\infty\) such that, for every pair of positive integers \(n,d\) and every known arrival floor \(q\in[1/2,1]\),
\[
c\,r_{n,d,q}\le \mathfrak R^*_{n,d,q}\le C\,r_{n,d,q},
\]
where \(r_{n,d,q}\) is the squared-error benchmark defined in \cref{obj:def:rate} and \(\mathfrak R^*_{n,d,q}\) is the minimax squared-error risk defined in \cref{obj:def:point-risk}.
\end{theoremv}

The comparison separates ordinary sampling error from the cost of recovering outcome means across many cells. The sampling contribution is \(n^{-1}\); the additional squared-error contribution is
\((1-q)^2\min\{1,[d/(n\log(e+n))]^2\}\).
Its dependence on missingness reflects the estimator's use of cell corrections weighted by missing-outcome mass. The matching lower bound makes this dependence a property of the minimax criterion over \(\mathcal M(d,q)\). In particular, the comparison identifies the precision scale uniformly as the arrival floor approaches one.

Honest inference is governed by the square root of the same benchmark, with comparison constants allowed to depend on the fixed coverage requirement.

\begin{theoremv}{thm:interval-frontier}[Fixed coverage length frontier]
For each fixed noncoverage level \(\alpha\in(0,1/2)\), there exist constants \(0<c_\alpha<C_\alpha<\infty\), depending only on \(\alpha\), such that, for every pair of positive integers \(n,d\) and every \(q\in[1/2,1]\), the minimax expected interval length \(\mathfrak L^*_{n,d,q,\alpha}\) from \cref{obj:def:length-risk} and the squared-error benchmark \(r_{n,d,q}\) from \cref{obj:def:rate} satisfy
\[
c_\alpha\sqrt{r_{n,d,q}}
\le \mathfrak L^*_{n,d,q,\alpha}
\le C_\alpha\sqrt{r_{n,d,q}}.
\]
\end{theoremv}

The risk-based interval construction explains the attainable square-root scale. The lower bound establishes its necessity for worst-law expected length among connected intervals satisfying uniform finite-sample coverage. Thus, for each fixed \(\alpha\), the same sampling and missingness-weighted dimension terms determine both squared-error precision and honest interval precision under the respective criteria in \cref{obj:def:point-risk,obj:def:length-risk}.

The upper-bound arguments are supplied by \cref{obj:lem:upper-risk,obj:lem:honest-interval-construction}. For the converse, \cref{obj:lem:parametric-floor-infrastructure} supplies the sampling lower bounds, while \cref{obj:thm:normalized-converse,obj:lem:prior-reduction} supplies the dimension lower bounds through normalized priors and testing reductions. Together, these arguments give the matching comparisons in \cref{obj:thm:point-frontier,obj:thm:interval-frontier}.

At complete arrival, the construction has a direct randomized-trial interpretation. It becomes the projected difference formed with the known treatment-allocation weights, and both minimax criteria have dimension-free parametric rates. The following result also quantifies the effect of projection on each sample and on worst-law squared risk.

\begin{propositionv}{prop:complete-arrival-reduction}[Complete arrival and parametric frontiers]
There exist universal constants \(0<c<C\) such that, for every noncoverage level \(\alpha\in(0,1/2)\), there exist constants \(0<c_\alpha<C_\alpha\), depending only on \(\alpha\), for which the following statements hold for every positive integer sample size \(n\) and positive integer baseline-support bound \(d\) at arrival floor \(q=1\).

The error scale and squared-error benchmark satisfy
\[
g_{n,d,1}
=(1-1)\min\left\{1,\frac{d}{n\log(e+n)}\right\}=0,
\qquad
r_{n,d,1}=\frac1n,
\]
where the benchmark is defined in \cref{obj:def:rate}.
For every sample \(O^n=(O_1,\ldots,O_n)\) in the Cartesian observed-record space, the estimator of \cref{obj:def:mm-estimator} satisfies
\[
\widehat\tau_{\mathrm{MM}}(O^n)
=
\Pi_{[-1,1]}
\left[
\frac{2}{n}\sum_{i=1}^{n}(2A_i-1)Y_i^{\mathrm{obs}}
\right],
\]
where \(A_i\) is the binary treatment coordinate, \(Y_i^{\mathrm{obs}}=(RY)_i\) is the supplied fifth coordinate of record \(i\), and
\(\Pi_{[-1,1]}(z)=\max\{-1,\min\{1,z\}\}\) is interval projection.

For every full-data law \(P\in\mathcal M(d,1)\) and every such sample, with population average treatment effect
\(\tau(P)=\mathbb E_P[Y(1)-Y(0)]\), where \(Y(1)\) and \(Y(0)\) are the binary potential outcomes,
\[
\left|\widehat\tau_{\mathrm{MM}}(O^n)-\tau(P)\right|
\le
\left|
\frac{2}{n}\sum_{i=1}^{n}(2A_i-1)Y_i^{\mathrm{obs}}-\tau(P)
\right|.
\]
Under the iid observed-sample law induced by \(P\),
\[
\sup_{P\in\mathcal M(d,1)}
\mathbb E_P\!\left[
\bigl(\widehat\tau_{\mathrm{MM}}(O^n)-\tau(P)\bigr)^2
\right]
\le
\sup_{P\in\mathcal M(d,1)}
\mathbb E_P\!\left[
\left(
\frac{2}{n}\sum_{i=1}^{n}(2A_i-1)Y_i^{\mathrm{obs}}-\tau(P)
\right)^2
\right]
\le \frac4n.
\]
The minimax squared-error risk and minimax worst-law expected honest interval length of
\cref{obj:def:point-risk,obj:def:length-risk} satisfy, uniformly over \(n,d\ge1\),
\[
\frac{c}{n}
\le \mathfrak R^*_{n,d,1}
\le \frac{C}{n},
\qquad
\frac{c_\alpha}{\sqrt n}
\le \mathfrak L^*_{n,d,1,\alpha}
\le \frac{C_\alpha}{\sqrt n}.
\]
\end{propositionv}

Under balanced randomization, the unprojected expression is the standard allocation-weighted difference between treated and control outcomes, using the known allocation probability of one half. Projection weakly reduces its absolute error for every target in the treatment-effect range, and the displayed risk bound is uniform in the baseline-support bound. Complete arrival therefore yields squared-error precision of order \(n^{-1}\) and honest expected length of order \(n^{-1/2}\) across all \(d\ge1\). The endpoint follows from the complete-arrival branch of \cref{obj:def:mm-estimator}, projection contraction, and the sampling lower bounds in \cref{obj:lem:parametric-floor-infrastructure}.


\section{Statistical implications}

The comparisons in \cref{obj:thm:point-frontier,obj:thm:interval-frontier} distinguish sampling uncertainty from the additional cost of recovering outcome means across baseline cells. In the benchmark of \cref{obj:def:rate}, sampling contributes \(n^{-1}\) to squared error, while the dimension contribution is weighted by the square of the arrival deficit \(1-q\). Accordingly, the minimax root mean squared error has order
\[
n^{-1/2}+(1-q)\min\left\{1,\frac{d}{n\ell_n}\right\},
\qquad \ell_n=\log(e+n).
\]
The comparison constants are uniform over positive integers \(n,d\) and known arrival floors \(q\in[1/2,1]\). This uniformity permits sample size, baseline dimension, and outcome arrival to vary jointly.

Before saturation, when \(d\le n\ell_n\), the dimension contribution to root mean squared error is \((1-q)d/(n\ell_n)\). It matches the sampling scale when
\[
(1-q)d=\sqrt n\,\ell_n.
\]
Below this balance, sampling determines the order of precision; above it, the arrival-weighted dimension term determines that order. For a fixed arrival floor below one, the transition therefore occurs at baseline dimension of order \(\sqrt n\,\ell_n/(1-q)\), provided this lies within the unsaturated range. As the arrival floor increases, a larger baseline dimension becomes compatible with sampling-scale precision. In particular, \(1-q\le n^{-1/2}\) guarantees precision of order \(n^{-1/2}\) uniformly over all positive baseline-support bounds. Complete arrival gives the exact endpoint described in \cref{obj:prop:complete-arrival-reduction}.

Saturation occurs at \(d\ge n\ell_n\), where the dimension contribution to squared error equals \((1-q)^2\). Increasing the support bound within this regime leaves the benchmark unchanged. With a fixed arrival floor below one, the minimax risk then remains bounded away from zero as sample size grows along such a sequence. An arrival floor approaching one instead reduces the saturated contribution directly. Thus the rate separates the effect of increasing sample size relative to dimension from the effect of improving outcome arrival.

More generally, let \(d_n\) be a sequence of positive baseline-support bounds and \(q_n\in[1/2,1]\) a sequence of known arrival floors, with \(n\to\infty\). The point-estimation comparison implies that minimax squared risk tends to zero precisely when
\[
(1-q_n)\min\left\{1,\frac{d_n}{n\ell_n}\right\}\longrightarrow0.
\]
For an arrival floor fixed below one, this condition becomes \(d_n=o(n\ell_n)\). For \(q_n\to1\), it is satisfied for every dimension sequence. Sampling-order squared risk, of order \(n^{-1}\), is attained precisely when the displayed dimension scale is \(O(n^{-1/2})\). These statements describe the joint sequences through the same finite-sample comparison, with constants independent of the sequence.

Honest inference translates these regimes into uniform coverage with corresponding expected precision. For each fixed \(\alpha\in(0,1/2)\), \cref{obj:thm:interval-frontier} gives minimax worst-law expected length of order
\[
n^{-1/2}+(1-q)\min\left\{1,\frac{d}{n\ell_n}\right\},
\]
with constants depending only on \(\alpha\). The interval criterion in \cref{obj:def:interval-class,obj:def:length-risk} requires coverage at least \(1-\alpha\) for every law in the specified class at each sample size. Consequently, the joint-sequence condition above also characterizes when this optimal worst-law expected length tends to zero. Sampling-dominated regimes yield honest expected length of order \(n^{-1/2}\), while saturation yields length of order \(n^{-1/2}+1-q\). The lower comparison concerns the worst-law expected length among honest connected intervals; the construction in \cref{obj:def:mm-interval} attains the corresponding order throughout the class.

The converse also has a useful implication for model enlargement. The sampling and dimension lower bounds in \cref{obj:lem:parametric-floor-infrastructure,obj:thm:normalized-converse,obj:lem:prior-reduction} are supported by randomized submodels with a constant surrogate. Any broader law class on the same observed experiment, with the same treatment-effect target, that contains these submodels inherits their lower bounds: its worst-law criterion includes the laws used in the converse. This establishes a necessary precision scale for such containing classes. The matching upper bounds characterize attainable precision over the finite-cell randomized-MAR model in \cref{obj:def:law-class}, with its binary measurements, known arrival floor, and baseline-support bound. Together, the two directions identify how sampling, dimension, and outcome arrival govern optimal precision in that model.


\section{Limitations and future work}

The results in \cref{obj:thm:point-frontier,obj:thm:interval-frontier} characterize minimax squared-error risk and honest expected interval length for the finite-cell randomized-MAR class in \cref{obj:def:law-class}, with a known cellwise arrival floor \(q\in[1/2,1]\) and a specified baseline-support bound \(d\). These conditions provide the reference point for several extensions concerning weaker arrival guarantees, adaptation, and practical implementation.

Arrival floors below one half are a natural direction for further analysis. Extending the comparison to \(0<q<1/2\) requires establishing upper and lower bounds that account for the dependence on the floor, including regimes in which it approaches zero as the sample size increases. The bounds established for \(q\in[1/2,1]\) do not determine that dependence. A separate question is adaptation when the floor is unknown. The construction in \cref{obj:def:mm-estimator,obj:def:mm-interval} uses the supplied floor, so a procedure that selects its degree of correction or interval radius from the data would require a new analysis of risk and uniform coverage. Such an analysis would need to address arrival probabilities in sparsely represented cells, where estimating a valid lower bound can itself be difficult.

The comparison with broader models has a precise lower-bound interpretation. Any law class on the same observed experiment, with the same treatment-effect target, that contains the submodels used in \cref{obj:lem:parametric-floor-infrastructure,obj:thm:normalized-converse,obj:lem:prior-reduction} inherits the corresponding lower bounds. For point estimation, the enlarged worst-law risk includes the laws supporting the converse. For honest intervals, coverage over the enlarged class entails coverage over those submodels, and the enlarged worst-law expected length includes their expected lengths. Matching procedures for broader classes require further analysis of their identifying restrictions and statistical complexity; the matching upper bounds here concern the class in \cref{obj:def:law-class}.

Computational implementation is another direction for development. The estimator in \cref{obj:def:mm-estimator} averages over auxiliary randomization conditional on the observed sample. This conditional average specifies a statistical procedure, while an efficient method for evaluating it remains a practical question. Possible approaches include exact evaluation where feasible and numerical approximation elsewhere. An approximation would require control of its additional error, together with an assessment of how that error affects the risk bound and the coverage of intervals centered at the computed estimate. The statistical construction supplies no computational complexity guarantee.

Practical constants also warrant investigation. The polynomial branch of \cref{obj:def:mm-estimator} requires \(n\ge e^{128}-e\), and the proof's fallback contribution is \(C_{\mathrm{fall}}=4e^{128}\). With the exhibited upper-risk constant, the interval in \cref{obj:def:mm-interval} is therefore the full range \([-1,1]\) whenever \(\log(e+n)<128\). The theorems establish universal-constant minimax orders, while these particular numerical choices do not yield an informative interval at ordinary sample sizes. Sharper constants, alternative tuning, and computational study are needed to turn the theoretical construction into a practical procedure; each change would require new risk and coverage justification.

Empirical calibration can complement this work through simulations that vary support size, cell-frequency imbalance, and arrival probabilities within the specified model. Mean squared error, coverage, and expected interval length would provide direct measures of performance, including near complete arrival and in settings with many sparsely occupied cells. Such exercises could guide implementation and identify useful tuning choices, while uniform guarantees for those choices would remain a separate theoretical task.


\appendix

\section*{Appendices}

\section{Identification and estimation bounds}

The constructive bounds first identify the superpopulation treatment effect from the observed law and then control the estimator's squared error uniformly.

\subsection{Identification under randomized MAR}

The cell means in \cref{obj:synth_5,obj:def:identified-functional} use explicit conventions at null conditioning events. The following result establishes identification with those conventions and the randomized-MAR restrictions.

\begin{lemmav}{lem:identification-infrastructure}[Randomized MAR identification]
Let \(d\) be a nonnegative integer, let \(q\in[1/2,1]\) be the cellwise arrival floor, and let \(P\) be a full-data probability law with baseline variable \(X\) taking values in a set of \(d\) labels, binary treatment \(A\), binary potential surrogates \(S(0),S(1)\), binary potential outcomes \(Y(0),Y(1)\), binary realized surrogate and outcome \(S,Y\), and binary outcome-arrival indicator \(R\). Suppose that \(P\) satisfies:
\begin{itemize}
\item \textbf{(Randomized independence.)} \cref{obj:ass:randomized-independence}.
\item \textbf{(Balanced randomization.)} \cref{obj:ass:balanced-randomization}.
\item \textbf{(Consistency.)} \cref{obj:ass:consistency}.
\item \textbf{(Missing at random.)} \cref{obj:ass:mar}.
\item \textbf{(Arrival floor.)} \cref{obj:ass:arrival-floor} with floor \(q\).
\end{itemize}
Write \(P_O\) for the image of \(P\) under the observation map \(O=(X,A,S,R,RY)\), where \(RY\) is the observed outcome mark. Then the superpopulation average treatment effect satisfies
\[
\tau(P)=\mathbb E_P[Y(1)-Y(0)]=\Psi(P_O),
\]
where \(\Psi\) is the observed-law functional in \cref{obj:def:identified-functional}, with every contribution from a null \((X,A,S)\) cell equal to zero by definition.
\end{lemmav}

\begin{proof}[Proof of \cref{obj:lem:identification-infrastructure}]
\begin{enumerate}
\item For each \(a\in\{0,1\}\), consistency in
\cref{obj:ass:consistency} gives
\[
P(A=a,Y=1)=P(A=a,Y(a)=1).
\]
To calculate the latter probability, define the potential-coordinate vector and its relevant finite subset by
\[
\mathbf V_{\mathrm{pot}}
:=(X,S(0),S(1),Y(0),Y(1)),
\qquad
\mathcal V_{\mathrm{pot},a}
:=
\left\{
(x,s_0,s_1,y_0,y_1)
\in\{0,\ldots,d-1\}\times\{0,1\}^{4}
:\ y_a=1
\right\},
\quad a\in\{0,1\}.
\]
Summing the atomwise factorization supplied by
\cref{obj:ass:randomized-independence} over this subset yields
\[
\begin{aligned}
P(A=a,Y(a)=1)
&=\sum_{\boldsymbol v\in\mathcal V_{\mathrm{pot},a}}
  P(A=a,\mathbf V_{\mathrm{pot}}=\boldsymbol v)\\
&=P(A=a)
  \sum_{\boldsymbol v\in\mathcal V_{\mathrm{pot},a}}
  P(\mathbf V_{\mathrm{pot}}=\boldsymbol v)\\
&=P(A=a)P(Y(a)=1).
\end{aligned}
\]
Here the sums partition the indicated events by their potential-coordinate tuples.
By \cref{obj:ass:balanced-randomization} and the binary treatment partition,
\[
P(A=1)=\frac12,
\qquad
P(A=0)=1-P(A=1)=\frac12.
\]
Consequently, the two arm identities are
\[
2P(A=1,Y=1)=P(Y(1)=1),
\qquad
2P(A=0,Y=1)=P(Y(0)=1).
\]
% lean: CausalSmith.Stat.MarNearcompleteFrontier.randomized_arm_outcome

\item Since both potential outcomes are binary, their numerical values equal their success indicators. Linearity of expectation therefore gives
\[
\begin{aligned}
\tau(P)
&=\mathbb E_P\!\left[
  \mathbf 1_{\{Y(1)=1\}}-\mathbf 1_{\{Y(0)=1\}}
  \right]\\
&=P(Y(1)=1)-P(Y(0)=1).
\end{aligned}
\]
% lean: CausalSmith.Stat.MarNearcompleteFrontier.tau_eq_potential_mass

\item Define the cell events, for their full range of indices, by
\[
\mathcal E_{x,a,s}:=\{X=x,A=a,S=s\},
\qquad
x\in\{0,\ldots,d-1\},\quad a,s\in\{0,1\}.
\]
The observation map preserves \(X,A,S,R\), and on \(\{R=1\}\) its outcome mark equals \(Y\). Thus its image law satisfies
\[
\begin{gathered}
P_O(\mathcal E_{x,a,s})=P(\mathcal E_{x,a,s}),
\qquad
P_O(A=a)=P(A=a),\\
P_O(\mathcal E_{x,a,s},R=1)
  =P(\mathcal E_{x,a,s},R=1),\\
P_O(\mathcal E_{x,a,s},R=1,RY=1)
  =P(\mathcal E_{x,a,s},R=1,Y=1).
\end{gathered}
\]
In particular, \cref{obj:def:identified-functional} supplies
\[
\frac{P_O(\mathcal E_{x,a,s})}{P_O(A=a)}
=\frac{P(\mathcal E_{x,a,s})}{P(A=a)},
\qquad
\mu_P(x,a,s)
=\frac{P(\mathcal E_{x,a,s},R=1,Y=1)}
       {P(\mathcal E_{x,a,s},R=1)},
\]
with zero assigned to every ratio whose denominator is zero.

We identify the weighted contribution by splitting according to the cell mass. If \(P(\mathcal E_{x,a,s})=0\), nonnegativity and event inclusion imply
\[
0\le P(\mathcal E_{x,a,s},Y=1)
\le P(\mathcal E_{x,a,s})=0.
\]
Hence both the weighted observed contribution and
\(P(\mathcal E_{x,a,s},Y=1)/P(A=a)\) are zero.

If \(P(\mathcal E_{x,a,s})>0\), event inclusion gives
\[
0<P(\mathcal E_{x,a,s})\le P(A=a).
\]
The arrival floor in \cref{obj:ass:arrival-floor} and the stipulated range of \(q\) give
\[
\frac{P(\mathcal E_{x,a,s},R=1)}
     {P(\mathcal E_{x,a,s})}
\ge q\ge\frac12>0,
\qquad
P(\mathcal E_{x,a,s},R=1)>0.
\]
Conditional independence in \cref{obj:ass:mar}, applied to \(R=1\) and \(Y=1\), gives the cross-product identity
\[
P(\mathcal E_{x,a,s},R=1,Y=1)\,
P(\mathcal E_{x,a,s})
=
P(\mathcal E_{x,a,s},R=1)\,
P(\mathcal E_{x,a,s},Y=1).
\]
Dividing by the two positive cell masses yields
\[
\mu_P(x,a,s)
=
\frac{P(\mathcal E_{x,a,s},R=1,Y=1)}
     {P(\mathcal E_{x,a,s},R=1)}
=
\frac{P(\mathcal E_{x,a,s},Y=1)}
     {P(\mathcal E_{x,a,s})}.
\]
Multiplication by the cell weight now proves, in both cases,
\[
\frac{P_O(\mathcal E_{x,a,s})}{P_O(A=a)}
\,\mu_P(x,a,s)
=
\frac{P(\mathcal E_{x,a,s},Y=1)}{P(A=a)}.
\]
% lean: CausalSmith.Stat.MarNearcompleteFrontier.cellContribution_identified

\item For either treatment arm, every realization with \(A=a\) and \(Y=1\) belongs to exactly one cell indexed by its values of \(X\) and \(S\). The disjoint finite-cell partition therefore gives
\[
\sum_{x=0}^{d-1}\sum_{s=0}^{1}
P(\mathcal E_{x,a,s},Y=1)
=
P(A=a,Y=1),
\qquad a\in\{0,1\}.
\]
% lean: CausalSmith.Stat.MarNearcompleteFrontier.cellOutcome_sum

\item Substituting the cell-contribution identity into
\cref{obj:def:identified-functional}, distributing the finite sums, and applying the arm partitions gives
\[
\begin{aligned}
\Psi(P_O)
&=\sum_{x=0}^{d-1}\sum_{s=0}^{1}
\left\{
\frac{P(\mathcal E_{x,1,s},Y=1)}{P(A=1)}
-
\frac{P(\mathcal E_{x,0,s},Y=1)}{P(A=0)}
\right\}\\
&=\frac{P(A=1,Y=1)}{P(A=1)}
  -\frac{P(A=0,Y=1)}{P(A=0)}\\
&=2P(A=1,Y=1)-2P(A=0,Y=1)\\
&=P(Y(1)=1)-P(Y(0)=1)\\
&=\tau(P).
\end{aligned}
\]
The null-cell case above establishes the stipulated zero contributions. When \(d=0\), the baseline-label set is empty, so a probability law with \(X\) taking values in that set cannot exist; the assertion for that case is vacuous.
% lean: CausalSmith.Stat.MarNearcompleteFrontier.tau_eq_psi
\end{enumerate}
\end{proof}

Randomized independence and consistency give the arm-specific outcome means their potential-outcome interpretation. Within an occupied baseline--treatment--surrogate cell, MAR makes the arrived-outcome conditional mean equal to the cell's realized-outcome mean. The arrival floor ensures that this arrived conditional mean has a positive conditioning probability. The weights in \cref{obj:def:identified-functional} are the conditional cell probabilities within each randomized arm; balanced assignment fixes the arm probabilities at one half. Together, these components give the finite-cell form of the nested-regression identification formula in \citet[Lemma~1.1, equation~(10), under Assumptions~1--3]{KallusMao2025}.

The null-cell conventions complete this interpretation over the entire baseline alphabet. A cell with zero mass receives zero weight, and its arrived-outcome mean is assigned zero. On occupied cells, the positive arrival floor makes the conditional mean well defined. Thus the totalized functional supplies the original superpopulation target for every risk and coverage calculation over \(\mathcal M(d,q)\).

\subsection{Uniform estimation control}

The estimator in \cref{obj:def:mm-estimator} combines an arrived-outcome contribution with a correction for missing outcomes. Its polynomial, ratio, complete-arrival, and fallback branches are covered by one uniform squared-error statement.

\begin{lemmav}{lem:upper-risk}[Uniform upper risk bound]
There exists a universal numerical constant \(0<C_U<\infty\) such that, for every choice satisfying
\begin{itemize}
\item \textbf{(Sample size and dimension.)} \(n,d\) are integers with \(n\ge1\) and \(d\ge1\);
\item \textbf{(Arrival floor.)} \(q\in[1/2,1]\);
\item \textbf{(Law and sampling.)} \(P\in\mathcal M(d,q)\), and \(O_1,\ldots,O_n\) are independent observed records with common law \(P_O\), the observed-law image of \(P\);
\end{itemize}
the estimator \(\widehat\tau_{\mathrm{MM}}\) satisfies
\[
\mathbb E_P\!\left[
  \bigl(\widehat\tau_{\mathrm{MM}}-\tau(P)\bigr)^2
\right]
\le C_U r_{n,d,q},
\]
where the expectation is over the stated independent sample.
\end{lemmav}

\begin{proof}[Proof of \cref{obj:lem:upper-risk}]
\begin{enumerate}
\item Define the numerical constants
\[
\begin{aligned}
C_{\mathrm{bias}}&=268435456,\\
C_{\mathrm{switch}}&=201719808,\\
C_{\mathrm{log}}&=720(4/3)^6(e+1),\\
C_{\mathrm{non}}&=
8\bigl(C_{\mathrm{switch}}^2C_{\mathrm{log}}+119\bigr)
+2C_{\mathrm{bias}}^2+104,\\
C_{\mathrm{fall}}&=4e^{128},\\
C_U&=C_{\mathrm{non}}+C_{\mathrm{fall}}.
\end{aligned}
\]
These constants are finite and positive. Fix parameters and a law satisfying the hypotheses, and put
\[
\delta=1-q\in[0,1/2],
\qquad
\ell_n=\log(e+n).
\]
By the independent sampling hypothesis,
\[
(O_1,\ldots,O_n)\sim P_O^{\otimes n}.
\]
From \cref{obj:def:rate},
\[
r_{n,d,q}\ge \frac1n>0.
\]
We establish the risk bounds \(C_{\mathrm{fall}}r_{n,d,q}\) and
\(C_{\mathrm{non}}r_{n,d,q}\) in the corresponding branches of
\cref{obj:def:mm-estimator}. Each is at most \(C_Ur_{n,d,q}\).
% lean: CausalSmith.Stat.MarNearcompleteFrontier.upper_risk

\item Suppose first that \(q=1\). The arrival floor makes the conditional arrival probability equal to one in every positive-probability cell. Consequently,
\[
P(R=0)=0.
\]
Define the complete-arrival score on the observed-record space by
\[
F_{\mathrm{complete}}:O\longrightarrow[-2,2],
\qquad
F_{\mathrm{complete}}(O)=2(2A-1)RY.
\]
The restrictions defining \cref{obj:def:law-class} give
\[
\begin{aligned}
\mathbb E_P[F_{\mathrm{complete}}(O)]
&=2P(A=1,Y=1)-2P(A=0,Y=1)\\
&=P(Y(1)=1)-P(Y(0)=1)
=\tau(P).
\end{aligned}
\]
Indeed, consistency replaces the realized outcome within each arm by its potential outcome, and randomized independence and balanced assignment give
\(2P(A=a,Y=1)=P(Y(a)=1)\) for \(a\in\{0,1\}\).
Also,
\[
\operatorname{Var}_P(F_{\mathrm{complete}}(O))
\le \mathbb E_P[F_{\mathrm{complete}}(O)^2]\le4.
\]
Independence of the observations therefore yields
\[
\mathbb E_P\!\left[
\left\{\frac1n\sum_{i=1}^nF_{\mathrm{complete}}(O_i)-\tau(P)\right\}^2
\right]
=\frac{\operatorname{Var}_P(F_{\mathrm{complete}}(O))}{n}
\le\frac4n.
\]
The binary potential outcomes imply \(\tau(P)\in[-1,1]\), and interval projection satisfies
\[
\left|\Pi_{[-1,1]}(t)-\tau(P)\right|
\le |t-\tau(P)|,
\qquad t\in\mathbb R.
\]
Thus the complete-arrival branch satisfies
\[
\mathbb E_P[(\widehat\tau_{\mathrm{MM}}-\tau(P))^2]
\le \frac4n
=4r_{n,d,1}
\le C_{\mathrm{fall}}r_{n,d,1}
\le C_Ur_{n,d,1}.
\]
% lean: CausalSmith.Stat.MarNearcompleteFrontier.upper_complete_arrival_risk_rate

\item Suppose next that
\[
q\ne1,
\qquad
\ell_n<128\ \text{or}\ d\ge n\ell_n.
\]
When \(\ell_n<128\), exponentiating its definition gives
\[
n<e+n=e^{\ell_n}<e^{128}.
\]
Both the estimator and the target lie in \([-1,1]\), so
\[
\mathbb E_P[(\widehat\tau_{\mathrm{MM}}-\tau(P))^2]
\le4
\le\frac{4e^{128}}n
\le C_{\mathrm{fall}}r_{n,d,q}.
\]

For the large-dimension branch, introduce the following quantities, defined for every cell, including null cells:
\[
\begin{gathered}
\mathcal J=\{0,\ldots,d-1\}\times\{0,1\}^2,
\qquad j=(x,a,s)\in\mathcal J,\\
\pi_j=P(X=x,A=a,S=s),\\
t_j=P(X=x,A=a,S=s,R=1),
\qquad
w_j=\pi_j-t_j=P(X=x,A=a,S=s,R=0),\\
\mu_P(j)=\mu_P(x,a,s),
\qquad
\eta_j=\mu_P(j)-\frac12,
\qquad
\epsilon_j=2a-1.
\end{gathered}
\]
The mean \(\mu_P(j)\) uses the zero-denominator convention in
\cref{obj:def:identified-functional}. In particular,
\[
|\eta_j|\le\frac12,\qquad
\sum_{j\in\mathcal J}\pi_j=1.
\]
The arrival floor gives \(q\pi_j\le t_j\) on positive cells. On a null cell,
\(\pi_j=t_j=w_j=0\), so the same inequality holds. Since \(q\ge1/2\),
\[
0\le w_j\le\delta\pi_j\le2\delta t_j,
\qquad
\sum_{j\in\mathcal J}w_j\le\delta.
\]
Balanced assignment gives
\[
\sum_{j\in\mathcal J}\epsilon_j\pi_j=0.
\]
Applying \cref{obj:lem:identification-infrastructure} and centering the cell means therefore gives
\[
\tau(P)
=2\sum_{j\in\mathcal J}\epsilon_j\pi_j\eta_j
=2\sum_{j\in\mathcal J}\epsilon_j(t_j\eta_j+w_j\eta_j).
\]

Define the fallback score by
\[
F_{\mathrm{fallback}}:O\longrightarrow[-1,1],
\qquad
F_{\mathrm{fallback}}(O)=2(2A-1)R(RY-1/2).
\]
On the support of \(P_O\), its mean is
\[
\mathbb E_P[F_{\mathrm{fallback}}(O)]
=2\sum_{j\in\mathcal J}\epsilon_jt_j\eta_j.
\]
Hence
\[
\left|\mathbb E_P[F_{\mathrm{fallback}}(O)]-\tau(P)\right|
=\left|2\sum_{j\in\mathcal J}\epsilon_jw_j\eta_j\right|
\le\sum_{j\in\mathcal J}w_j
\le\delta.
\]
The score bound and independent sampling imply
\[
\mathbb E_P\!\left[
\left\{\frac1n\sum_{i=1}^nF_{\mathrm{fallback}}(O_i)
-\mathbb E_P[F_{\mathrm{fallback}}(O)]\right\}^2
\right]\le\frac1n.
\]
Using projection contraction and
\((u+v)^2\le2u^2+2v^2\), we obtain
\[
\mathbb E_P[(\widehat\tau_{\mathrm{MM}}-\tau(P))^2]
\le\frac2n+2\delta^2.
\]
For \(d\ge n\ell_n\), the minimum in \cref{obj:def:rate} equals one. Thus
\[
\frac2n+2\delta^2
=2r_{n,d,q}
\le C_{\mathrm{fall}}r_{n,d,q}.
\]
Both fallback subcases consequently satisfy the bound with \(C_U\).
The cell quantities just defined will also be used below.
% lean: CausalSmith.Stat.MarNearcompleteFrontier.upper_fallback_branches_risk_rate

\item Consider now the remaining regime
\[
q\ne1,\qquad \ell_n\ge128,\qquad d<n\ell_n.
\]
The tuning parameters and benchmark become
\[
m=\frac n8,\qquad
k=\lfloor\ell_n/64\rfloor,\qquad
B=256\ell_n,
\]
\[
\frac{\ell_n}{128}\le k\le\frac{\ell_n}{64}
=\frac{B}{16384},
\qquad
g_{n,d,q}=\frac{\delta d}{n\ell_n},
\qquad
r_{n,d,q}=\frac1n+g_{n,d,q}^2.
\]
The lower bound on \(k\) follows from
\(\lfloor\ell_n/64\rfloor\ge\ell_n/64-1\ge\ell_n/128\).

Extend the observed sample to an infinite independent sequence with law \(P_O\). Independently draw
\[
N\sim\operatorname{Poisson}(n/2),
\qquad
Z_i\stackrel{\mathrm{iid}}{\sim}\operatorname{Unif}\{1,2,3,4\},
\qquad i\ge1,
\]
and define the uncapped stream sets
\[
\mathcal I_t^\star=\{i\le N:Z_i=t\},
\qquad t\in\{1,2,3,4\}.
\]
Write \(\mathbb E_\star\) and \(\Pr_\star\) for expectation and probability under this joint experiment. For an observed record \(o\), define its stream histogram count by
\[
K_{t,o}^{\mathrm{hist}}
=\sum_{i\in\mathcal I_t^\star}\mathbf1\{O_i=o\}.
\]
For \(u_{t,o}\in[0,1]\), conditioning on \(N\) gives the joint probability generating function
\[
\mathbb E_\star\!\left[
\prod_{t=1}^4\prod_{o\in O}
u_{t,o}^{K_{t,o}^{\mathrm{hist}}}\right]
=
\exp\!\left\{
m\sum_{t=1}^4\sum_{o\in O}P_O(\{o\})(u_{t,o}-1)
\right\}.
\]
Its factorization shows that these histogram counts are independent Poisson variables with means \(mP_O(\{o\})\).

Evaluate the correction formulas of \cref{obj:def:mm-estimator} on the uncapped streams:
\[
\begin{aligned}
V_j^\star&=\frac1m\sum_{i\in\mathcal I_2^\star}
\mathbf1\{(X_i,A_i,S_i)=j,\ R_i=0\},\\
C_j^{\prime\star}&=\sum_{i\in\mathcal I_3^\star}
\mathbf1\{(X_i,A_i,S_i)=j,\ R_i=1\},\\
C_j^\star&=\sum_{i\in\mathcal I_4^\star}
\mathbf1\{(X_i,A_i,S_i)=j,\ R_i=1\},\\
U_j^\star&=\sum_{i\in\mathcal I_4^\star}
\mathbf1\{(X_i,A_i,S_i)=j,\ Y_i^{\mathrm{obs}}=1\},\\
H_j^\star&=\sum_{v=1}^{k-1}
a_v(U_j^\star-C_j^\star/2)(C_j^\star-1)_{v-1},\\
D_j^\star&=
\begin{cases}
U_j^\star/C_j^\star-1/2,&C_j^\star>0,\\
0,&C_j^\star=0,
\end{cases}\\
G_j^\star&=
H_j^\star\mathbf1\{C_j^{\prime\star}\le B/4\}
+D_j^\star\mathbf1\{C_j^{\prime\star}>B/4\}.
\end{aligned}
\]
Define the linear term, aggregate correction, raw estimate, and projected ideal estimate by
\[
\begin{aligned}
\tau_{\mathrm{lin}}
&=\frac2m\sum_{i\in\mathcal I_1^\star}
(2A_i-1)R_i(Y_i^{\mathrm{obs}}-1/2),\\
W_{\mathrm{corr}}
&=\sum_{j\in\mathcal J}\epsilon_jV_j^\star G_j^\star,\\
\tau_{\mathrm{raw}}
&=\tau_{\mathrm{lin}}+2W_{\mathrm{corr}},\\
\widetilde\tau_\star&=\Pi_{[-1,1]}(\tau_{\mathrm{raw}}).
\end{aligned}
\]

The capped randomized estimate \(\widetilde\tau\) in
\cref{obj:def:mm-estimator} agrees with \(\widetilde\tau_\star\) when \(N\le n\).
On the complementary event their squared errors are bounded by four.
Conditional Jensen and this coupling therefore give
\[
\begin{aligned}
\mathbb E_P[(\widehat\tau_{\mathrm{MM}}-\tau(P))^2]
&\le\mathbb E_\star[(\widetilde\tau-\tau(P))^2]\\
&\le\mathbb E_\star[(\widetilde\tau_\star-\tau(P))^2]
+4\Pr_\star(N>n).
\end{aligned}
\]
Exponential Markov, applied to \(2^N\), yields
\[
\Pr_\star(N>n)
\le2^{-n}\mathbb E_\star[2^N]
=\exp\{-(\log2-1/2)n\}
\le e^{-n/8}
\le\frac8n.
\]
Here \(\log2-1/2\ge1/8\), and \(n/8\le e^{n/8}\) gives the last inequality. Consequently,
\[
\mathbb E_P[(\widehat\tau_{\mathrm{MM}}-\tau(P))^2]
\le\mathbb E_\star[(\widetilde\tau_\star-\tau(P))^2]
+\frac{32}{n}.
\]
% lean: CausalSmith.Stat.MarNearcompleteFrontier.upper_nonfallback_risk_le_ideal_add_rate

\item All the correction statistics have finite second moments: their polynomial parts are finite polynomials in Poisson counts, and their ratio parts are bounded. Projection contraction and the variance--bias identity give
\[
\mathbb E_\star[(\widetilde\tau_\star-\tau(P))^2]
\le
\operatorname{Var}_\star(\tau_{\mathrm{raw}})
+\{\mathbb E_\star[\tau_{\mathrm{raw}}]-\tau(P)\}^2.
\]
To bound the first term, define
\[
\begin{aligned}
K_{\mathrm{lin},+}
&=\sum_{i\in\mathcal I_1^\star}
\mathbf1\{R_i=1,\ A_i=Y_i^{\mathrm{obs}}\},\\
K_{\mathrm{lin},-}
&=\sum_{i\in\mathcal I_1^\star}
\mathbf1\{R_i=1,\ A_i\ne Y_i^{\mathrm{obs}}\}.
\end{aligned}
\]
Then
\[
\tau_{\mathrm{lin}}
=\frac{K_{\mathrm{lin},+}-K_{\mathrm{lin},-}}m.
\]
Each count is Poisson with mean at most \(m\). Using
\(\operatorname{Var}(U-V)\le2\operatorname{Var}(U)+2\operatorname{Var}(V)\),
\[
\operatorname{Var}_\star(\tau_{\mathrm{lin}})
\le\frac{4m}{m^2}
=\frac{32}{n}.
\]
Applying the corresponding sum inequality to
\(\tau_{\mathrm{raw}}=\tau_{\mathrm{lin}}+2W_{\mathrm{corr}}\) gives
\[
\operatorname{Var}_\star(\tau_{\mathrm{raw}})
\le\frac{64}{n}
+8\operatorname{Var}_\star(W_{\mathrm{corr}}).
\]
% lean: CausalSmith.Stat.MarNearcompleteFrontier.upper_fourStreamRawEstimate_variance_le_aggregate

\item For every \(j\in\mathcal J\), define its arrived-count intensity and light-branch probability by
\[
\zeta_j=mt_j\ge0,
\qquad
\wp_j=\Pr_\star(C_j^{\prime\star}\le B/4)\in[0,1].
\]
Poisson splitting gives
\[
\mathbb E_\star[V_j^\star]=w_j,
\qquad
\mathbb E_\star[(V_j^\star)^2]=w_j^2+\frac{w_j}{m},
\]
\[
C_j^{\prime\star},C_j^\star\sim\operatorname{Poisson}(\zeta_j),
\qquad
U_j^\star\mid C_j^\star
\sim\operatorname{Binomial}(C_j^\star,\mu_P(j)).
\]
The conditional binomial statement includes \(t_j=0\): then both counts vanish almost surely and \(\mu_P(j)=0\).

For integers \(v\ge1\), use the falling factorial
\[
(C_j^\star)_v=\prod_{t=0}^{v-1}(C_j^\star-t).
\]
Conditional expectation of \(U_j^\star\), followed by the Poisson factorial-moment identity, gives
\[
\begin{aligned}
\mathbb E_\star[
(U_j^\star-C_j^\star/2)(C_j^\star-1)_{v-1}]
&=\eta_j\mathbb E_\star[(C_j^\star)_v]\\
&=\eta_j\zeta_j^v.
\end{aligned}
\]
Thus the coefficient identity in \cref{obj:def:mm-estimator} yields
\[
\mathbb E_\star[H_j^\star]
=\eta_j\{1-Q_k(\zeta_j)\}.
\]
Similarly, conditioning on \(C_j^\star>0\) gives
\[
\mathbb E_\star[D_j^\star]
=\eta_j\Pr_\star(C_j^\star>0)
=\eta_j(1-e^{-\zeta_j}).
\]
Independence of the pilot stream from the fourth stream now implies
\[
\mathbb E_\star[G_j^\star]-\eta_j
=-\eta_j\{\wp_jQ_k(\zeta_j)+(1-\wp_j)e^{-\zeta_j}\}.
\]
The first-stream mean and independence of the second stream from the pilot and fourth streams also give
\[
\mathbb E_\star[\tau_{\mathrm{lin}}]
=2\sum_{j\in\mathcal J}\epsilon_jt_j\eta_j,
\qquad
\mathbb E_\star[V_j^\star G_j^\star]
=w_j\mathbb E_\star[G_j^\star].
\]
Combining these identities with the cell decomposition of the target,
\[
\mathbb E_\star[\tau_{\mathrm{raw}}]-\tau(P)
=
2\sum_{j\in\mathcal J}
\epsilon_jw_j\{\mathbb E_\star[G_j^\star]-\eta_j\}.
\]
% lean: CausalSmith.Stat.MarNearcompleteFrontier.upper_fourStreamRawEstimate_bias_eq

\item We next derive the polynomial and pilot bounds used to control this bias. For \(0<z\le B\), the argument \(1-2z/B\) belongs to \([-1,1]\), so the Chebyshev bound \(|T_k(1-2z/B)|\le1\) gives
\[
Q_k(z)\ge0,
\qquad
zQ_k(z)
=\frac{B\{1-T_k(1-2z/B)\}}{2k^2}
\le\frac{B}{k^2}.
\]
At \(z=0\), the convention is \(Q_k(0)=1\).
The degree bound already established implies
\[
\frac{B}{k^2}
\le\frac{4194304}{\ell_n}.
\]

For the second-moment estimates, define the coefficient norm of a polynomial
\(\mathcal P(u)\) by
\[
\|\mathcal P\|_{\mathrm{coeff},1}
=\sum_{v\ge0}|[u^v]\mathcal P|,
\]
where \([u^v]\mathcal P\) denotes its coefficient of \(u^v\); the sum is finite. Define the shifted Chebyshev polynomials by
\[
\mathcal P_v(u)=T_v(1-2u),
\qquad v\in\mathbb N_0.
\]
Their recurrence and initial values give
\[
\mathcal P_{v+1}(u)
=2(1-2u)\mathcal P_v(u)-\mathcal P_{v-1}(u)
\quad(v\ge1),
\]
\[
\|\mathcal P_0\|_{\mathrm{coeff},1}=1,
\qquad
\|\mathcal P_1\|_{\mathrm{coeff},1}=3,
\]
\[
\|\mathcal P_{v+1}\|_{\mathrm{coeff},1}
\le6\|\mathcal P_v\|_{\mathrm{coeff},1}
+\|\mathcal P_{v-1}\|_{\mathrm{coeff},1}.
\]
Induction therefore yields
\[
\|\mathcal P_v\|_{\mathrm{coeff},1}\le7^v.
\]
Indeed, the induction step uses
\(6\cdot7^v+7^{v-1}\le7^{v+1}\).
Since
\[
Q_k(Bu)=\frac{1-\mathcal P_k(u)}{2k^2u},
\]
with the removable value at zero, deleting the zero constant term and dividing by \(u\) cannot increase the coefficient norm. Hence
\[
\sum_{v=1}^{k-1}|a_v|B^v
\le\frac{1+7^k}{2k^2}
\le7^k
\le8^k.
\]

On observed-law support, \(0\le U_j^\star\le C_j^\star\), and
\[
|(U_j^\star-C_j^\star/2)(C_j^\star-1)_{v-1}|
\le\frac{(C_j^\star)_v}{2}
\le\frac{(C_j^\star)^v}{2}.
\]
For every \(z\ge0\) and \(1\le v\le k\),
\[
(z/B)^v\le\max\{1,(z/B)^k\}.
\]
The coefficient bound therefore gives the pointwise envelope
\[
(H_j^\star)^2
\le\frac{64^k}{4}
\max\left\{1,\left(\frac{C_j^\star}{B}\right)^{2k}\right\}
\le\frac{64^k}{4}
\left\{1+\left(\frac{C_j^\star}{B}\right)^{2k}\right\}.
\]
For an integer \(v\ge0\), expansion of an ordinary Poisson moment into factorial moments gives
\[
\mathbb E_\star[(C_j^\star)^v]
=\sum_{t=0}^v {v\brace t}\zeta_j^t
\le\sum_{t=0}^v\binom vt v^{\,v-t}\zeta_j^t
=(\zeta_j+v)^v.
\]
Here \({v\brace t}\) counts partitions into \(t\) nonempty blocks: choosing one representative per block and assigning the remaining elements gives
\({v\brace t}\le\binom vt v^{v-t}\). The case \(v=0\) uses the constant zeroth moment. Integrating the pointwise envelope yields
\[
\mathbb E_\star[(H_j^\star)^2]
\le\frac{64^k}{4}
\left\{1+\left(\frac{\zeta_j+2k}{B}\right)^{2k}\right\}
\le\frac{64^k}{2}
\max\left\{1,\left(\frac{\zeta_j+2k}{B}\right)^{2k}\right\}.
\]
% lean: CausalSmith.Stat.MarNearcompleteFrontier.upper_streamH_second_moment_le

\item When \(\zeta_j\le B\), the degree calibration gives
\[
0\le\frac{\zeta_j+2k}{B}\le\frac32\le e^{1/2}.
\]
Consequently,
\[
\mathbb E_\star[(H_j^\star)^2]
\le\frac{64^k}{2}e^k
\le e^{7k}
\le e^{\ell_n/8},
\]
where \(64^k\le e^{6k}\) and \(7k\le\ell_n/8\).

For \(\zeta_j\ge B\), define
\[
\upsilon_j=\frac{\zeta_j+2k}{B}.
\]
From \(k\le B/16384\) and \(B\le\zeta_j\),
\[
2k\le\zeta_j,\qquad
k\upsilon_j\le\frac{\zeta_j}{8192},
\qquad
6k+2k\upsilon_j\le\frac{\zeta_j}{8}.
\]
Using \(\upsilon_j\le e^{\upsilon_j}\) in the second-moment envelope gives
\[
\mathbb E_\star[(H_j^\star)^2]
\le
64^k\max\{1,\upsilon_j^{2k}\}
\le e^{6k+2k\upsilon_j}
\le e^{\zeta_j/8}.
\]

Exponential Markov for the Poisson pilot count, using \(4^{-C_j^{\prime\star}}\) and \(4^{C_j^{\prime\star}}\), respectively, gives
\[
\wp_j\le
\exp\{B\log4/4-3\zeta_j/4\},
\qquad
1-\wp_j\le
\exp\{3\zeta_j-B\log4/4\}.
\]
Since \(1\le\log4\le2\), for \(\zeta_j\ge B\) these imply
\[
\wp_j\le e^{-\zeta_j/4},
\qquad
\wp_j\mathbb E_\star[(H_j^\star)^2]
\le e^{-\zeta_j/8}.
\]
Also, for every \(\zeta_j\ge0\),
\[
(1-\wp_j)e^{-\zeta_j}\le e^{-B/16}.
\]
To see this, if \(\zeta_j\ge B/16\), use \(1-\wp_j\le1\). Otherwise, the upper-tail bound gives
\[
(1-\wp_j)e^{-\zeta_j}
\le e^{2\zeta_j-B\log4/4}
\le e^{-B/16}.
\]

For \(\zeta_j\ge B\), Cauchy--Schwarz and \(\wp_j^2\le\wp_j\) yield
\[
\{\wp_j|\mathbb E_\star[H_j^\star]|\}^2
\le\wp_j\mathbb E_\star[(H_j^\star)^2]
\le e^{-\zeta_j/8},
\]
and hence
\[
\wp_j|\mathbb E_\star[H_j^\star]|
\le e^{-\zeta_j/16}.
\]
Using
\(\eta_jQ_k(\zeta_j)=\eta_j-\mathbb E_\star[H_j^\star]\), we obtain
\[
\begin{aligned}
\wp_j|\eta_jQ_k(\zeta_j)|
&\le\frac{\wp_j}{2}
+\wp_j|\mathbb E_\star[H_j^\star]|\\
&\le\frac32e^{-B/16}.
\end{aligned}
\]
The pilot bias identity therefore gives, on these large-intensity cells,
\[
|\mathbb E_\star[G_j^\star]-\eta_j|
\le2e^{-B/16}.
\]
% lean: CausalSmith.Stat.MarNearcompleteFrontier.upper_stream_pilot_large_bias_abs_le

\item We can now sum the bias. If \(\zeta_j=0\), then \(t_j=0\), and
\(0\le w_j\le2\delta t_j\) forces \(w_j=0\). Thus null arrived cells contribute zero to the weighted bias.

For \(0<\zeta_j\le B\), polynomial nonnegativity and
\(w_j\le2\delta t_j=2\delta\zeta_j/m\) give
\[
w_j\wp_j|Q_k(\zeta_j)|
\le w_jQ_k(\zeta_j)
\le\frac{2\delta}{m}\frac{B}{k^2}
\le\frac{2\delta}{m}\frac{4194304}{\ell_n}.
\]
There are \(4d\) cells in \(\mathcal J\), so
\[
\sum_{\substack{j\in\mathcal J\\\zeta_j\le B}}
w_j\wp_j|Q_k(\zeta_j)|
\le
4d\,\frac{2\delta}{m}\frac{4194304}{\ell_n}
=C_{\mathrm{bias}}g_{n,d,q}.
\]
For a small-intensity cell, the pilot bias identity and
\(|\eta_j|\le1/2\) give
\[
2w_j|\mathbb E_\star[G_j^\star]-\eta_j|
\le w_j\wp_j|Q_k(\zeta_j)|+w_je^{-B/16}.
\]
For a large-intensity cell, the preceding bound gives
\[
2w_j|\mathbb E_\star[G_j^\star]-\eta_j|
\le4w_je^{-B/16}.
\]
Together, these yield the uniform cellwise bound
\[
2w_j|\mathbb E_\star[G_j^\star]-\eta_j|
\le
\mathbf1\{\zeta_j\le B\}w_j\wp_j|Q_k(\zeta_j)|
+4w_je^{-B/16}.
\]
The triangle inequality and \(\sum_jw_j\le\delta\) therefore imply
\[
|\mathbb E_\star[\tau_{\mathrm{raw}}]-\tau(P)|
\le C_{\mathrm{bias}}g_{n,d,q}
+4\delta e^{-16\ell_n}.
\]
The remainder satisfies
\[
(4\delta e^{-16\ell_n})^2
=16\delta^2e^{-32\ell_n}
\le4e^{-32\ell_n}
\le4e^{-\ell_n}
=\frac4{e+n}
\le\frac4n.
\]
Squaring the bias bound with \((u+v)^2\le2u^2+2v^2\) gives
\[
\{\mathbb E_\star[\tau_{\mathrm{raw}}]-\tau(P)\}^2
\le
2C_{\mathrm{bias}}^2g_{n,d,q}^2+\frac8n
\le(2C_{\mathrm{bias}}^2+8)r_{n,d,q}.
\]
Combining projection, the raw variance bound, and the capping penalty established above,
\[
\begin{aligned}
\mathbb E_P[(\widehat\tau_{\mathrm{MM}}-\tau(P))^2]
&\le
8\operatorname{Var}_\star(W_{\mathrm{corr}})
+\frac{96}{n}
+(2C_{\mathrm{bias}}^2+8)r_{n,d,q}\\
&\le
8\operatorname{Var}_\star(W_{\mathrm{corr}})
+(2C_{\mathrm{bias}}^2+104)r_{n,d,q}.
\end{aligned}
\]
% lean: CausalSmith.Stat.MarNearcompleteFrontier.upper_nonfallback_risk_le_aggregate

\item The histogram coordinates used by two distinct cells are disjoint. Their independence, together with \(\epsilon_j^2=1\), gives
\[
\operatorname{Var}_\star(W_{\mathrm{corr}})
=\sum_{j\in\mathcal J}
\operatorname{Var}_\star(V_j^\star G_j^\star).
\]
We first bound the ratio component.

Since \(|D_j^\star|\le1/2\), its variance is at most \(1/4\). Thus, when \(\zeta_j\le15\),
\[
\operatorname{Var}_\star(D_j^\star)
\le\frac14\le\frac4{1+\zeta_j}.
\]
For \(\zeta_j>15\), define the arrived-success and arrived-failure fractions by
\[
F_{j,+}=
\begin{cases}
U_j^\star/C_j^\star,&C_j^\star>0,\\
0,&C_j^\star=0,
\end{cases}
\qquad
F_{j,-}=
\begin{cases}
(C_j^\star-U_j^\star)/C_j^\star,&C_j^\star>0,\\
0,&C_j^\star=0.
\end{cases}
\]
Then
\[
D_j^\star=\frac{F_{j,+}-F_{j,-}}2.
\]
Conditioning on the arrived count gives
\[
\begin{aligned}
\mathbb E_\star[(F_{j,+}-\mu_P(j))^2]
&\le
\Pr_\star(C_j^\star=0)
+\mathbb E_\star\!\left[
\frac{\mathbf1\{C_j^\star>0\}}{C_j^\star}\right],\\
\mathbb E_\star[(F_{j,-}-(1-\mu_P(j)))^2]
&\le
\Pr_\star(C_j^\star=0)
+\mathbb E_\star\!\left[
\frac{\mathbf1\{C_j^\star>0\}}{C_j^\star}\right].
\end{aligned}
\]
On a positive count, the conditional binomial variance divided by the squared count is at most its reciprocal; on a zero count, the squared error is at most one.

For a Poisson count of intensity \(\zeta_j\), the zero-safe reciprocal on the right is at most one. If \(\zeta_j\le1\), this is at most \(4/(1+\zeta_j)\). If \(\zeta_j>1\), the pointwise bounds
\[
\mathbf1\{C_j^\star=0\}
+\frac{\mathbf1\{C_j^\star>0\}}{C_j^\star}
\le\frac2{C_j^\star+1}
\]
and the Poisson identity
\[
\mathbb E_\star\!\left[\frac1{C_j^\star+1}\right]
=\frac{1-e^{-\zeta_j}}{\zeta_j}
\]
give
\[
\Pr_\star(C_j^\star=0)
+\mathbb E_\star\!\left[
\frac{\mathbf1\{C_j^\star>0\}}{C_j^\star}\right]
\le\frac2{\zeta_j}
\le\frac4{1+\zeta_j}.
\]
The reciprocal identity has value one at \(\zeta_j=0\), by its continuous extension. Using the center
\(\eta_j=\{\mu_P(j)-(1-\mu_P(j))\}/2\), we conclude
\[
\begin{aligned}
\operatorname{Var}_\star(D_j^\star)
&\le\mathbb E_\star[(D_j^\star-\eta_j)^2]\\
&\le\frac12\mathbb E_\star[(F_{j,+}-\mu_P(j))^2]
+\frac12\mathbb E_\star[(F_{j,-}-(1-\mu_P(j)))^2]\\
&\le\frac4{1+\zeta_j}.
\end{aligned}
\]
Together with the small-intensity case, this holds for every cell.

The second and fourth streams are independent. Their first and second moments consequently give
\[
\begin{aligned}
\operatorname{Var}_\star(V_j^\star D_j^\star)
&=\left(w_j^2+\frac{w_j}{m}\right)
\operatorname{Var}_\star(D_j^\star)
+\frac{w_j}{m}\{\mathbb E_\star[D_j^\star]\}^2\\
&\le\frac{4w_j^2}{1+\zeta_j}
+\frac{17w_j}{4m},
\end{aligned}
\]
where \(\{\mathbb E_\star[D_j^\star]\}^2\le1/4\).
Also,
\[
\frac{w_j^2}{1+mt_j}
\le\frac{2\delta}{m}w_j,
\qquad
\sum_{j\in\mathcal J}\frac{w_j^2}{1+mt_j}
\le\frac{2\delta^2}{m}.
\]
Since \(\delta+2\delta^2\le1\),
\[
\sum_{j\in\mathcal J}
\left(\frac{w_j}{m}+\frac{w_j^2}{1+mt_j}\right)
\le\frac1m.
\]
Replacing the coefficient \(4\) of the nonnegative quadratic term by \(17/4\) therefore yields
\[
\sum_{j\in\mathcal J}
\operatorname{Var}_\star(V_j^\star D_j^\star)
\le\frac{17}{4m}
=\frac{34}{n}.
\]
% lean: CausalSmith.Stat.MarNearcompleteFrontier.upper_streamVD_variance_sum_le

\item Define the total switching second moment by
\[
\mathcal V_{\mathrm{switch}}
=
\sum_{j\in\mathcal J}
\mathbb E_\star[
\{V_j^\star(G_j^\star-D_j^\star)\}^2].
\]
The identity
\[
G_j^\star-D_j^\star
=\mathbf1\{C_j^{\prime\star}\le B/4\}(H_j^\star-D_j^\star)
\]
and independence of the second, third, and fourth streams imply
\[
\mathbb E_\star[
\{V_j^\star(G_j^\star-D_j^\star)\}^2]
=
\left(w_j^2+\frac{w_j}{m}\right)\wp_j
\mathbb E_\star[(H_j^\star-D_j^\star)^2].
\]
Since \((u-v)^2\le2u^2+2v^2\) and
\(\mathbb E_\star[(D_j^\star)^2]\le1/4\),
\[
\mathbb E_\star[
\{V_j^\star(G_j^\star-D_j^\star)\}^2]
\le
\left(w_j^2+\frac{w_j}{m}\right)
\left\{2\wp_j\mathbb E_\star[(H_j^\star)^2]+\frac{\wp_j}{2}\right\}.
\]

For \(\zeta_j\le B\), use \(\wp_j\le1\), the light second-moment bound, and
\(w_j\le2\delta B/m\). Since \(e^{\ell_n/8}\ge1\), this gives
\[
\mathbb E_\star[
\{V_j^\star(G_j^\star-D_j^\star)\}^2]
\le
3e^{\ell_n/8}
\frac{4\delta^2B^2+2\delta B}{m^2}.
\]
Summing over at most \(4d\) cells,
\[
\sum_{\substack{j\in\mathcal J\\\zeta_j\le B}}
\mathbb E_\star[
\{V_j^\star(G_j^\star-D_j^\star)\}^2]
\le
12d\,e^{\ell_n/8}
\frac{4\delta^2B^2+2\delta B}{m^2}.
\]

For \(\zeta_j>B\), the pilot-weighted polynomial bound and
\(\wp_j\le e^{-\zeta_j/4}\le e^{-\zeta_j/8}\) give
\[
\mathbb E_\star[
\{V_j^\star(G_j^\star-D_j^\star)\}^2]
\le
\frac52\left(w_j^2+\frac{w_j}{m}\right)e^{-B/8}.
\]
Nonnegativity and \(\sum_jw_j\le\delta\) imply
\(\sum_jw_j^2\le\delta^2\). Hence
\[
\sum_{\substack{j\in\mathcal J\\\zeta_j>B}}
\mathbb E_\star[
\{V_j^\star(G_j^\star-D_j^\star)\}^2]
\le\frac52\left(\delta^2+\frac{\delta}{m}\right)e^{-B/8}
\le\frac{25}{n}.
\]
For the last inequality, use \(n\ge1\), \(\delta\le1\),
\(\delta/m\le8\), and
\[
e^{-B/8}=e^{-32\ell_n}\le e^{-\ell_n}
=\frac1{e+n}\le\frac1n.
\]
Thus
\[
\mathcal V_{\mathrm{switch}}
\le
12d\,e^{\ell_n/8}
\frac{4\delta^2B^2+2\delta B}{m^2}
+\frac{25}{n}.
\]

To convert the first term into the benchmark, define
\[
\Lambda_n=\frac{e^{\ell_n/8}\ell_n^3}{n}.
\]
Since \(\delta^2\le\delta\), \(\ell_n\le\ell_n^2\), and \(B=256\ell_n\),
\[
4\delta^2B^2+2\delta B
\le262656\,\delta\ell_n^2.
\]
Substitution of \(m=n/8\) gives the calibrated bound
\[
12d\,e^{\ell_n/8}
\frac{4\delta^2B^2+2\delta B}{m^2}
\le
C_{\mathrm{switch}}g_{n,d,q}\Lambda_n,
\qquad
C_{\mathrm{switch}}=12\cdot64\cdot262656.
\]
The sixth term of the exponential series yields
\[
\frac{(3\ell_n/4)^6}{720}\le e^{3\ell_n/4},
\qquad
\ell_n^6\le720(4/3)^6e^{3\ell_n/4}.
\]
Since \(e^{\ell_n}=e+n\le(e+1)n\),
\[
\Lambda_n^2
=\frac{e^{\ell_n/4}\ell_n^6}{n^2}
\le\frac{720(4/3)^6e^{\ell_n}}{n^2}
\le\frac{C_{\mathrm{log}}}{n}.
\]
Finally, the nonnegative square
\((g_{n,d,q}-C_{\mathrm{switch}}\Lambda_n)^2\) gives
\[
C_{\mathrm{switch}}g_{n,d,q}\Lambda_n
\le\frac{g_{n,d,q}^2}{2}
+\frac{C_{\mathrm{switch}}^2\Lambda_n^2}{2}
\le\frac{g_{n,d,q}^2}{2}
+\frac{C_{\mathrm{switch}}^2C_{\mathrm{log}}}{2n}.
\]
Consequently,
\[
\mathcal V_{\mathrm{switch}}
\le
\frac{g_{n,d,q}^2}{2}
+\frac{C_{\mathrm{switch}}^2C_{\mathrm{log}}/2+25}{n}.
\]
% lean: CausalSmith.Stat.MarNearcompleteFrontier.upper_stream_correction_second_moment_sum_rate

\item The decomposition
\[
V_j^\star G_j^\star
=V_j^\star D_j^\star
+V_j^\star(G_j^\star-D_j^\star)
\]
and the variance sum inequality give
\[
\operatorname{Var}_\star(V_j^\star G_j^\star)
\le
2\operatorname{Var}_\star(V_j^\star D_j^\star)
+2\mathbb E_\star[
\{V_j^\star(G_j^\star-D_j^\star)\}^2].
\]
Summing and using the two preceding bounds,
\[
\begin{aligned}
\operatorname{Var}_\star(W_{\mathrm{corr}})
&\le\frac{68}{n}+2\mathcal V_{\mathrm{switch}}\\
&\le
g_{n,d,q}^2+
\frac{C_{\mathrm{switch}}^2C_{\mathrm{log}}+118}{n}\\
&\le
\bigl(C_{\mathrm{switch}}^2C_{\mathrm{log}}+119\bigr)
r_{n,d,q}.
\end{aligned}
\]
Substitution into the aggregate risk inequality yields
\[
\begin{aligned}
\mathbb E_P[(\widehat\tau_{\mathrm{MM}}-\tau(P))^2]
&\le
\left\{
8\bigl(C_{\mathrm{switch}}^2C_{\mathrm{log}}+119\bigr)
+2C_{\mathrm{bias}}^2+104
\right\}r_{n,d,q}\\
&=C_{\mathrm{non}}r_{n,d,q}
\le C_Ur_{n,d,q}.
\end{aligned}
\]
Together with the complete-arrival and fallback branches, this establishes the asserted uniform bound for every admissible \(n,d,q,P\) and the stated independent sample.
% lean: CausalSmith.Stat.MarNearcompleteFrontier.upper_risk
\end{enumerate}
\end{proof}

The bound accommodates arbitrary cell-frequency imbalance within \(\mathcal M(d,q)\). Its benchmark combines the sampling contribution \(n^{-1}\) with the squared missingness-weighted dimension scale and supplies the upper bounds for point estimation and honest intervals.

To describe the components of the estimation argument, fix a cell \(j=(x,a,s)\) and write
\[
p_j=P(X=x,A=a,S=s),\qquad
b_j=P(X=x,A=a,S=s,R=0),
\]
\[
\theta_j=\mu_P(x,a,s)-\frac12,\qquad
\lambda_j=mP(X=x,A=a,S=s,R=1),
\]
where \(m=n/8\) is the mean stream-size parameter in \cref{obj:def:mm-estimator}. Here \(p_j\) is total cell mass, \(b_j\) is missing-outcome mass, \(\theta_j\) is the centered arrived-outcome mean, and \(\lambda_j\) is the expected arrived count in one stream of the auxiliary Poisson experiment. These quantities are defined on null cells as well, using the conventions of \cref{obj:synth_5}.

The scientific decomposition separates the treatment contrast into arrived outcomes and the outcomes represented by missing records. Centering at one half aligns the first-stream statistic with the centered cell means \(\theta_j\); balanced randomization gives the centered contrast the same population target. MAR identifies the mean needed for the missing component, while the second-stream statistic \(V_j\) estimates its cell weight \(b_j\). The arrival floor controls the relation between missing and arrived cell mass, explaining why the additional error scale is weighted by \(1-q\).

For sparsely observed cells, the polynomial statistic \(H_j\) in \cref{obj:def:mm-estimator} supplies the correction through arrived-count moments. The polynomial \(Q_k\), its specified value at zero, and the falling-factorial convention are all fixed in that construction. The falling-factorial terms implement unbiased estimation of the polynomial terms in the auxiliary Poisson experiment. This use of polynomial approximation and unbiased moment estimation follows the methodological framework of \citet[Sections~III--IV]{JiaoVenkatHanWeissman2015}; related polynomial constructions for discrete treatment-effect estimation appear in \citet{CausalSmith2026,CausalSmith2026Annotation}. The light-cell component addresses aggregate missing-outcome correction when individual cell means have few arrived observations.

Cells selected as heavily populated use the centered ratio \(D_j\), with its specified zero-count value. The selection count \(C'_j\) comes from a separate stream from the counts \(C_j,U_j\) used to form either correction. In the auxiliary Poisson experiment, this separation supports control of the selection rule together with the correction's sampling variation. The thresholded statistic \(G_j\) therefore joins polynomial approximation in light cells with ratio estimation in heavy cells.

Projection and conditional averaging turn the auxiliary construction into an observed-sample estimator in \([-1,1]\) without increasing squared error. The branch definitions and their full parameter ranges are stated once in \cref{obj:def:mm-estimator}; \cref{obj:lem:upper-risk} controls all of them with the same universal constant.

\subsection{Honest clipped intervals}

The uniform risk bound supplies a finite-sample confidence radius through \cref{obj:def:mm-interval}. Clipping the resulting interval to the treatment-effect range gives simultaneous coverage and length guarantees.

\begin{lemmav}{lem:honest-interval-construction}[Honest clipped confidence intervals]
Suppose that:
\begin{itemize}
\item \textbf{(Sample size and dimension.)} \(n,d\) are integers with \(n,d\ge 1\).
\item \textbf{(Arrival floor.)} \(q\in[1/2,1]\).
\item \textbf{(Noncoverage level.)} \(\alpha\in(0,1/2)\).
\item \textbf{(Uniform risk certificate.)} \(C_U>0\) certifies the uniform squared-error bound
\[
\mathbb E_{P_O^{\otimes n}}
\left[
  \bigl(\widehat\tau_{\mathrm{MM}}(O^n)-\tau(P)\bigr)^2
\right]
\le C_U r_{n,d,q}
\]
for every \(n,d\ge1\), every \(q\in[1/2,1]\), and every \(P\in\mathcal M(d,q)\), under any iid observed-sample law. Here \(r_{n,d,q}\) is defined in \cref{obj:def:rate}, and
\(\tau(P)=\mathbb E_P[Y(1)-Y(0)]\) is the population average treatment effect.
\end{itemize}
Let \(\widehat I_{\mathrm{MM},\alpha}\) and its radius \(h_{n,d,q,\alpha}\) be as in \cref{obj:def:mm-interval}, using this \(C_U\). The interval has uniform coverage:
\[
P_O^{\otimes n}
\left\{
  \tau(P)\in\widehat I_{\mathrm{MM},\alpha}(O^n)
\right\}
\ge 1-\alpha
\qquad
\text{for every }P\in\mathcal M(d,q).
\]
For every observed sample \(O^n\), its length satisfies
\[
\operatorname{length}
\left\{\widehat I_{\mathrm{MM},\alpha}(O^n)\right\}
\le
\min\{2,2h_{n,d,q,\alpha}\}
=
\min\left\{2,2\sqrt{\frac{C_Ur_{n,d,q}}{\alpha}}\right\}.
\]
Consequently, for every \(P\in\mathcal M(d,q)\),
\[
\mathbb E_{P_O^{\otimes n}}
\left[
  \operatorname{length}
  \left\{\widehat I_{\mathrm{MM},\alpha}(O^n)\right\}
\right]
\le
\min\left\{2,2\sqrt{\frac{C_Ur_{n,d,q}}{\alpha}}\right\}.
\]
If \(h_{n,d,q,\alpha}<2\), then for every \(P\in\mathcal M(d,q)\),
\[
P_O^{\otimes n}
\left\{
  \tau(P)\notin\widehat I_{\mathrm{MM},\alpha}(O^n)
\right\}
\le\alpha.
\]
If \(h_{n,d,q,\alpha}=2\), then for every \(P\in\mathcal M(d,q)\) and every observed sample \(O^n\),
\[
\tau(P)\in\widehat I_{\mathrm{MM},\alpha}(O^n).
\]
\end{lemmav}

\begin{proof}[Proof of \cref{obj:lem:honest-interval-construction}]
Fix the parameters in the statement. For every observed sample \(o\in O^n\), write the endpoints from \cref{obj:def:mm-interval} as
\[
L(o)=\max\{-1,\widehat\tau_{\mathrm{MM}}(o)-h_{n,d,q,\alpha}\},
\qquad
U(o)=\min\{1,\widehat\tau_{\mathrm{MM}}(o)+h_{n,d,q,\alpha}\}.
\]
The radius satisfies
\[
0\le h_{n,d,q,\alpha}
=\min\left\{2,\sqrt{\frac{C_Ur_{n,d,q}}{\alpha}}\right\}
\le2.
\]

\begin{enumerate}
\item The endpoint formulas give
\[
L(o)\ge-1,\qquad U(o)\le1,
\]
and
\[
L(o)\ge\widehat\tau_{\mathrm{MM}}(o)-h_{n,d,q,\alpha},
\qquad
U(o)\le\widehat\tau_{\mathrm{MM}}(o)+h_{n,d,q,\alpha}.
\]
Subtracting yields both bounds
\[
U(o)-L(o)\le2,
\qquad
U(o)-L(o)\le2h_{n,d,q,\alpha}.
\]
Consequently,
\[
U(o)-L(o)\le\min\{2,2h_{n,d,q,\alpha}\}.
\]
If \(2\le\sqrt{C_Ur_{n,d,q}/\alpha}\), then the radius equals \(2\), and both sides of
\[
\min\{2,2h_{n,d,q,\alpha}\}
=
\min\left\{2,2\sqrt{\frac{C_Ur_{n,d,q}}{\alpha}}\right\}
\]
equal \(2\). If \(\sqrt{C_Ur_{n,d,q}/\alpha}\le2\), the identity follows by substituting that square root for the radius. This proves the pointwise length bound and its displayed equality.
% lean: CausalSmith.Stat.MarNearcompleteFrontier.honest_interval_construction

\item We establish the target range used in the coverage argument. For a fixed \(P\in\mathcal M(d,q)\), define
\[
\mathcal W_P=\operatorname{supp}(P),
\qquad
p_w=P(\{w\}),
\qquad
\Delta_w=Y(1)(w)-Y(0)(w)
\quad(w\in\mathcal W_P).
\]
The support is finite, and the binary potential outcomes give
\[
p_w\ge0,\qquad
\sum_{w\in\mathcal W_P}p_w=1,
\qquad
-1\le\Delta_w\le1.
\]
Thus
\[
-p_w\le p_w\Delta_w\le p_w.
\]
Summing these inequalities proves
\[
-1
=-\sum_{w\in\mathcal W_P}p_w
\le
\tau(P)=\sum_{w\in\mathcal W_P}p_w\Delta_w
\le
\sum_{w\in\mathcal W_P}p_w
=1.
\]
% lean: CausalSmith.Stat.MarNearcompleteFrontier.tau_range

\item The estimator also satisfies
\[
-1\le\widehat\tau_{\mathrm{MM}}(o)\le1
\qquad(o\in O^n).
\]
To derive this from \cref{obj:def:mm-estimator}, first observe that, for every real \(t\),
\[
-1\le\Pi_{[-1,1]}(t)=\max\{-1,\min\{1,t\}\}\le1:
\]
the lower inequality follows from the maximum, and the upper inequality follows because both arguments of that maximum are at most \(1\). This establishes the range when \(q=1\). When \(q\ne1\) and
\(\ell_n<128\) or \(d\ge n\ell_n\), the fallback branch is likewise a projection and has the same range.

For the remaining branch, define the Poisson weights and conditional stream averages by
\[
\omega_N=e^{-n/2}\frac{(n/2)^N}{N!}
\quad(N\in\mathbb N_0),
\]
\[
\overline\tau_N(o)
=
\frac{1}{4^N}
\sum_{\boldsymbol\zeta\in\{1,2,3,4\}^N}
\widetilde\tau(o;N,\boldsymbol\zeta)
\quad(0\le N\le n,\ o\in O^n),
\]
where \(\widetilde\tau(o;N,\boldsymbol\zeta)\) is the projected randomized estimate in \cref{obj:def:mm-estimator} with count \(N\) and stream assignment \(\boldsymbol\zeta\). There are \(4^N\) assignments, including one empty assignment when \(N=0\), and each estimate lies in \([-1,1]\). Therefore
\[
-4^N
\le
\sum_{\boldsymbol\zeta\in\{1,2,3,4\}^N}
\widetilde\tau(o;N,\boldsymbol\zeta)
\le4^N,
\qquad
-1\le\overline\tau_N(o)\le1.
\]
Moreover,
\[
\omega_N\ge0,\qquad
\sum_{N=0}^{\infty}\omega_N
=e^{-n/2}e^{n/2}=1,
\qquad
\sum_{N=0}^{n}\omega_N\le1.
\]
Since the randomized estimate equals zero when \(N>n\), this branch satisfies
\[
\widehat\tau_{\mathrm{MM}}(o)
=\sum_{N=0}^{n}\omega_N\overline\tau_N(o).
\]
Multiplying the stream-average bounds by the nonnegative weights and summing gives
\[
-1
\le-\sum_{N=0}^{n}\omega_N
\le\widehat\tau_{\mathrm{MM}}(o)
\le\sum_{N=0}^{n}\omega_N
\le1.
\]
% lean: CausalSmith.Stat.MarNearcompleteFrontier.tauhatMM_range

In particular, if \(h_{n,d,q,\alpha}=2\), the target and estimator ranges imply
\[
\widehat\tau_{\mathrm{MM}}(o)-2\le-1\le\tau(P),
\qquad
\tau(P)\le1\le\widehat\tau_{\mathrm{MM}}(o)+2.
\]
Taking the maximum in the lower endpoint and the minimum in the upper endpoint yields
\[
L(o)\le\tau(P)\le U(o)
\qquad(P\in\mathcal M(d,q),\ o\in O^n).
\]
% lean: CausalSmith.Stat.MarNearcompleteFrontier.honest_interval_construction

\item Suppose now that \(h_{n,d,q,\alpha}<2\), and fix \(P\in\mathcal M(d,q)\). Define the sample probability measure and squared-error function by
\[
\mu_{\mathrm{obs}}=P_O^{\otimes n},
\qquad
\mathcal E_P(o)
=\bigl(\widehat\tau_{\mathrm{MM}}(o)-\tau(P)\bigr)^2
\quad(o\in O^n).
\]
If \(\sqrt{C_Ur_{n,d,q}/\alpha}\le2\), the radius definition gives
\(h_{n,d,q,\alpha}=\sqrt{C_Ur_{n,d,q}/\alpha}\).
If that square root exceeds \(2\), the radius equals \(2\), contradicting the present case. Hence
\[
h_{n,d,q,\alpha}
=\sqrt{\frac{C_Ur_{n,d,q}}{\alpha}},
\qquad
h_{n,d,q,\alpha}^2=\frac{C_Ur_{n,d,q}}{\alpha}.
\]
By \cref{obj:def:rate} and \(n\ge1\),
\[
r_{n,d,q}\ge\frac1n>0.
\]
Together with \(C_U>0\) and \(\alpha>0\), this proves
\[
h_{n,d,q,\alpha}>0,
\qquad
h_{n,d,q,\alpha}^2>0,
\qquad
C_Ur_{n,d,q}=\alpha h_{n,d,q,\alpha}^2.
\]

Define the noncoverage and squared-error tail events by
\[
\mathcal B_P=\{o\in O^n:\tau(P)\notin[L(o),U(o)]\},
\qquad
\mathcal T_P=\{o\in O^n:h_{n,d,q,\alpha}^2\le\mathcal E_P(o)\}.
\]
We claim that
\[
\mathcal B_P\subseteq\mathcal T_P.
\]
Indeed, noncoverage gives either \(\tau(P)<L(o)\) or \(U(o)<\tau(P)\). In the first case, if
\(-1\le\widehat\tau_{\mathrm{MM}}(o)-h_{n,d,q,\alpha}\), then
\[
L(o)=\widehat\tau_{\mathrm{MM}}(o)-h_{n,d,q,\alpha},
\qquad
h_{n,d,q,\alpha}
<\widehat\tau_{\mathrm{MM}}(o)-\tau(P).
\]
Otherwise \(L(o)=-1\), contradicting \(\tau(P)\ge-1\). In the second case, if
\(\widehat\tau_{\mathrm{MM}}(o)+h_{n,d,q,\alpha}\le1\), then
\[
U(o)=\widehat\tau_{\mathrm{MM}}(o)+h_{n,d,q,\alpha},
\qquad
h_{n,d,q,\alpha}
<\tau(P)-\widehat\tau_{\mathrm{MM}}(o).
\]
Otherwise \(U(o)=1\), contradicting \(\tau(P)\le1\). Squaring the positive inequalities in the two possible cases gives
\[
h_{n,d,q,\alpha}^2
\le\bigl(\widehat\tau_{\mathrm{MM}}(o)-\tau(P)\bigr)^2,
\]
as claimed.

The sample space is finite, so the squared-error function is integrable. Its nonnegativity and Markov's inequality give
\[
h_{n,d,q,\alpha}^2\,\mu_{\mathrm{obs}}(\mathcal T_P)
\le
\int\mathcal E_P(o)\,d\mu_{\mathrm{obs}}(o).
\]
Applying the assumed uniform risk certificate to this iid sample law, and using event inclusion, yields
\[
h_{n,d,q,\alpha}^2\,\mu_{\mathrm{obs}}(\mathcal B_P)
\le
h_{n,d,q,\alpha}^2\,\mu_{\mathrm{obs}}(\mathcal T_P)
\le
\int\mathcal E_P(o)\,d\mu_{\mathrm{obs}}(o)
\le
C_Ur_{n,d,q}
=\alpha h_{n,d,q,\alpha}^2.
\]
Dividing by the strictly positive squared radius proves
\[
P_O^{\otimes n}\{\tau(P)\notin\widehat I_{\mathrm{MM},\alpha}(O^n)\}
\le\alpha.
\]
% lean: CausalSmith.Stat.MarNearcompleteFrontier.honest_interval_construction

\item To obtain honesty, fix any \(P\in\mathcal M(d,q)\) and split according to whether \(h_{n,d,q,\alpha}=2\). In that case, the pointwise inclusion proved above gives
\[
P_O^{\otimes n}\{\tau(P)\in\widehat I_{\mathrm{MM},\alpha}(O^n)\}
=1\ge1-\alpha.
\]
In the other case, \(h_{n,d,q,\alpha}\le2\) implies \(h_{n,d,q,\alpha}<2\). The coverage event is measurable, and the preceding noncoverage bound therefore gives
\[
P_O^{\otimes n}\{\tau(P)\in\widehat I_{\mathrm{MM},\alpha}(O^n)\}
=
1-P_O^{\otimes n}\{\tau(P)\notin\widehat I_{\mathrm{MM},\alpha}(O^n)\}
\ge1-\alpha.
\]
Since \(P\) was arbitrary, the interval has the asserted uniform coverage.
% lean: CausalSmith.Stat.MarNearcompleteFrontier.honest_interval_construction

\item Finally, the estimator range and nonnegative radius ensure that the endpoints form a nonempty interval, so
\[
0\le\operatorname{length}\{\widehat I_{\mathrm{MM},\alpha}(o)\}
=U(o)-L(o).
\]
These functions are integrable on the finite sample space. Integrating the pointwise bound from the first step under the probability measure \(P_O^{\otimes n}\) gives
\[
\begin{aligned}
\mathbb E_{P_O^{\otimes n}}
\!\left[\operatorname{length}\{\widehat I_{\mathrm{MM},\alpha}(O^n)\}\right]
&\le
\int\min\{2,2h_{n,d,q,\alpha}\}\,dP_O^{\otimes n}\\
&=\min\{2,2h_{n,d,q,\alpha}\}\\
&=\min\left\{2,2\sqrt{\frac{C_Ur_{n,d,q}}{\alpha}}\right\}.
\end{aligned}
\]
This holds for every \(P\in\mathcal M(d,q)\).
% lean: CausalSmith.Stat.MarNearcompleteFrontier.honest_interval_construction
\end{enumerate}
\end{proof}

For a radius below two, the squared-error bound supports the uniform noncoverage bound through Markov's inequality. At radius two, the interval equals the entire treatment-effect range because its center belongs to \([-1,1]\). Endpoint clipping preserves inclusion of a target in that range. These two branches yield the finite-sample coverage requirement in \cref{obj:def:interval-class}.

The interval's radius bounds its length by twice that radius, while the treatment-effect range bounds it by two. This deterministic cap and \cref{obj:lem:upper-risk} give the constructive expected-length bound in \cref{obj:thm:interval-frontier}.


\section{Normalized priors and the dimension lower bound}

The full-range converse uses two branches. When the squared arrival-weighted dimension scale is at least \(n^{-1}\), paired-label priors over admissible causal laws have close observed-sample distributions and separated population treatment effects. Signed interpolation weights coordinate opposite orientations within each pair, and a fixed filler label normalizes every latent draw. When that squared scale is below \(n^{-1}\), the parametric floors supply the required lower bounds. Together, the normalized prior comparison and parametric floors establish lower bounds across the full parameter range on the complete observed-record experiment used by \cref{obj:def:point-risk,obj:def:length-risk}.

Polynomial approximation and moment matching provide the general tools for constructing such mixture comparisons \citep{WuYang2016,Wu2020}. Here their implementation incorporates the randomized-MAR restrictions directly into the observed atom probabilities. The following construction collects the interpolation scheme, latent distribution, normalized causal laws, and auxiliary count experiment.

\begin{definitionv}{def:paired-prior-handle}[Normalized paired prior construction]
The paired-prior handle consists of the normalized priors \(\Pi_-,\Pi_+\), their observed-sample mixtures \(Q_-^{(n)},Q_+^{(n)}\), their associated unnormalized count laws, a fixed fallback record \(o_{\circ}\), and the count-ordering kernel \(\kappa_n\), constructed below for parameters satisfying:
\begin{itemize}
\item \textbf{(Sample size.)} \(n\) is an integer with \(n\ge1\).
\item \textbf{(Baseline dimension.)} \(d\) is an integer with \(d\ge1\).
\item \textbf{(Arrival floor.)} \(q\in[1/2,1]\).
\end{itemize}
The observed-record space \(O\) consists of tuples \((X,A,S,R,RY)\), with \(X\) taking values in a fixed set of \(d\) baseline labels and the remaining coordinates binary. Fix \(o_{\circ}\in O\).

Set
\[
\delta=1-q,\qquad
L=\log(e+n),\qquad
K=\lceil8L\rceil,\qquad
h=K^{-2},\qquad
m=K-1,\qquad
B_0=\frac{K}{1024n},
\]
\[
M=\min\left\{
\left\lfloor\frac{d-1}{2}\right\rfloor,
\lfloor nL\rfloor
\right\},
\qquad
c_0=\frac1{16},\qquad
b=1-\frac{\delta}{2},\qquad
\mu_0=\frac{b^2}{2q}.
\]
For \(0\le i\le m\), define the interpolation nodes and their Lagrange cardinal polynomials by
\[
x_i=\frac{1+h}{2}+\frac{1-h}{2}\cos\frac{i\pi}{m},
\qquad
l_i(t)=\prod_{\substack{0\le t'\le m\\t'\ne i}}
\frac{t-x_{t'}}{x_i-x_{t'}}.
\]
Set
\[
D=\sum_{i=0}^{m}|l_i(0)|,
\qquad
\nu_i=\frac{l_i(0)}{D}.
\]
The latent pair \((p,z)\in[0,B_0]\times\{-1,1\}\) has distribution
\[
\Pr\{(p,z)=(B_0x_i,\operatorname{sign}(\nu_i))\}
=\frac{h|\nu_i|}{x_i},
\qquad 0\le i\le m,
\]
\[
\Pr\{(p,z)=(0,1)\}
=1-\sum_{i=0}^{m}\frac{h|\nu_i|}{x_i}.
\]
Draw independent copies \((p_j,z_j)\), \(1\le j\le M\), and independent fair orientation signs, independent also of these copies, with
\[
\Pr(r_j=1)=\Pr(r_j=-1)=\frac12.
\]

For each \(\sigma\in\{-1,1\}\), pair \(j\) consists of two distinct baseline labels, each with unnormalized mass \(p_j\), with orientations
\[
u\in\{r_j,-r_j\}.
\]
At each label, \(X\) equals that label, and the potential surrogates and control outcome satisfy
\[
S(0)=S(1)=0,\qquad Y(0)=0.
\]
Treatment \(A\) is an independent fair coin, and control arrival is complete. The treated arrival probability and treated outcome mean are
\[
\rho_{\sigma,j,u}
=1-\frac{\delta}{2}(1+\sigma z_ju),
\qquad
\mu_{\sigma,j,u}
=\frac{b/2+c_0u}{\rho_{\sigma,j,u}}.
\]
Draw \(Y(1)\) as a Bernoulli variable with mean \(\mu_{\sigma,j,u}\), and draw treated arrival independently of \(Y(1)\) with probability \(\rho_{\sigma,j,u}\). In the coordinate order
\[
(A,R,RY)=(0,1,0),(1,0,0),(1,1,1),(1,1,0),
\]
the four observed atom coefficients at orientation \(u\) are
\[
v_{\sigma}(z_j,u)
=
\left(
\frac12,
\frac{\delta(1+\sigma z_ju)}4,
\frac b4+\frac{c_0u}{2},
\frac b4-\frac{c_0u}{2}-\frac{\sigma\delta z_ju}{4}
\right).
\]
The two orientations specify the complete eight-coordinate observed atom vector for pair \(j\).

Add one distinct fixed filler label with unnormalized mass
\[
f=1-2MhB_0.
\]
At this label, set
\[
S(0)=S(1)=0,\qquad Y(0)=0,
\]
let treatment be an independent fair coin, make arrival complete in both treatment arms, and draw \(Y(1)\) as a Bernoulli variable with mean \(\mu_0\). Its observed atom coefficients, in the same coordinate order, are
\[
v_{\mathrm{fill}}
=\left(\frac12,0,\frac{\mu_0}{2},\frac{1-\mu_0}{2}\right).
\]
At every label, the realized variables and observed outcome mark satisfy
\[
S=S(A),\qquad Y=Y(A),\qquad RY=R\,Y.
\]

Normalize the label masses by
\[
J=f+2\sum_{j=1}^{M}p_j,
\qquad
f\longmapsto\frac fJ,
\qquad
p_j\longmapsto\frac{p_j}{J}
\]
for each of the two labels in pair \(j\). Let \(P_{\sigma}\) be the resulting normalized causal law and \(P_{\sigma,O}\) its observed-law image. Define the prior and complete observed-sample mixture by
\[
\Pi_{\sigma}=\operatorname{Law}(P_{\sigma}),
\qquad
Q_{\sigma}^{(n)}
=\int P_{\sigma,O}^{\otimes n}\,d\Pi_{\sigma}(P_{\sigma}).
\]

For the associated unnormalized count experiment, set
\[
\lambda=4n.
\]
Conditional on the latent variables, assign independent Poisson counts to the eight atoms of each pair and the filler atoms, with respective mean vectors
\[
\lambda p_jv_{\sigma}(z_j,r_j),
\qquad
\lambda p_jv_{\sigma}(z_j,-r_j),
\qquad
\lambda f v_{\mathrm{fill}}.
\]
All other observed atoms have count zero. The associated count law is the mixture of this conditional count experiment over the latent draws.

Define the Markov kernel \(\kappa_n\) by uniformly ordering the counted multiset and returning its first \(n\) records when the total count is at least \(n\); otherwise it returns
\[
(o_{\circ},\ldots,o_{\circ})\in O^n.
\]
The kernel \(\kappa_n\) is independent of \(\sigma\) and of every latent variable.
\end{definitionv}

The restrictions defining \cref{obj:def:law-class} are built into \cref{obj:def:paired-prior-handle}. Each pair uses two labels, and the filler uses one, respecting the baseline-support bound. Independent fair treatment supplies randomized independence and balanced allocation. The potential surrogates are constant, and arrival is drawn independently of the treated potential outcome conditional on the label. Together with complete control arrival and the stipulated realization of the potential measurements, these choices supply consistency and MAR. Treated arrival takes values \(q\) and \(1\); the outcome means and observed atom coefficients are valid throughout the stated range of the arrival floor.

The opposite orientations serve complementary roles. They preserve a common total unnormalized mass for each pair while placing the sign-dependent arrival perturbation in both labels' observed distributions. Keeping the full eight-coordinate vector accounts jointly for control records, missing treated outcomes, arrived treated successes, and arrived treated failures at the two labels. Thus the mixture comparison concerns the complete observed record. The arrival deficit \leanref{sym:delta}{\ensuremath{\delta}}, defined in \cref{obj:def:paired-prior-handle}, also governs the attenuation of the treatment-effect perturbation.

Exact normalization is central to this construction. The filler mass is fixed across latent draws and signs, whereas the total mass \(J\) varies with the latent label masses. Dividing all label masses by that same total produces a probability law for every draw while preserving the conditional randomization and arrival restrictions. The normalized priors \leanref{sym:PiPair}{\ensuremath{\Pi_-,\Pi_+}} and complete sample mixtures \leanref{sym:QPair}{\ensuremath{Q_-^{(n)},Q_+^{(n)}}} in \cref{obj:def:paired-prior-handle} therefore belong to the causal model and its original sampling experiment.

The converse argument has three components. First, the interpolation weights support a comparison of the auxiliary Poisson mixtures, using the approximation and moment-matching tools developed in \citet{WuYang2016} and the Poisson-mixture framework described by \citet{Wu2020}. Second, the count-ordering transformation in \cref{obj:def:paired-prior-handle} connects this comparison to samples from the normalized laws. Its independence from the sign and latent variables makes it a common statistical transformation for both alternatives; the fixed fallback specifies its output for every count realization. Third, concentration of the treatment effects under the priors supplies separation around a deterministic center. The same-experiment normalization and the attenuation by the arrival deficit are the construction-specific ingredients linking these components to the randomized-MAR model.

The resulting statement fixes the comparison tolerance \leanref{sym:beta}{\ensuremath{\beta\in(0,1/4)}} before ranging over the sample size, dimension, and arrival floor. Its regime requires the squared missingness-and-dimension scale to be at least the sampling scale.

\begin{theoremv}{thm:normalized-converse}[Normalized priors with separated effects]
For every fixed \(\beta\in(0,1/4)\), there exists a constant \(c_\beta>0\) such that the following holds for every sample size \(n\), baseline-support bound \(d\), and arrival floor \(q\) satisfying:
\begin{itemize}
\item \textbf{(Sample size and dimension.)} \(n,d\) are integers with \(n,d\ge1\).
\item \textbf{(Arrival floor.)} \(q\in[1/2,1]\).
\item \textbf{(Separation scale.)} The scale
\[
g_{n,d,q}=(1-q)\min\left\{1,\frac{d}{n\log(e+n)}\right\}
\]
satisfies \(g_{n,d,q}^{\,2}\ge n^{-1}\).
\end{itemize}
There exist discrete probability priors \(\Pi_-\) and \(\Pi_+\) supported on \(\mathcal M(d,q)\) from \cref{obj:def:law-class}, whose laws satisfy \cref{obj:ass:randomized-independence,obj:ass:balanced-randomization,obj:ass:consistency,obj:ass:mar,obj:ass:arrival-floor}. Their complete observed-sample mixtures, defined by
\[
Q_-^{(n)}=\int P_O^{\otimes n}\,\Pi_-(dP),
\qquad
Q_+^{(n)}=\int P_O^{\otimes n}\,\Pi_+(dP),
\]
where \(P_O\) is the observed-law image of the full-data law \(P\), satisfy
\[
\operatorname{TV}\!\left(Q_-^{(n)},Q_+^{(n)}\right)
=\sup_{B\subseteq O^n}
\left|Q_-^{(n)}(B)-Q_+^{(n)}(B)\right|
\le\beta.
\]
Here \(O^n\) is the complete observed-record sample space, and the supremum is over measurable events. There also exists a deterministic center \(m_0\in\mathbb R\) such that, for the population average treatment effect
\[
\tau(P)=\mathbb E_P\!\left[Y(1)-Y(0)\right],
\]
where \(Y(1)\) and \(Y(0)\) are the binary potential outcomes,
\[
\Pi_-\!\left\{P:\tau(P)\le m_0-c_\beta g_{n,d,q}\right\}\ge1-\beta,
\qquad
\Pi_+\!\left\{P:m_0+c_\beta g_{n,d,q}\le\tau(P)\right\}\ge1-\beta.
\]
\end{theoremv}

The statistical content of \cref{obj:thm:normalized-converse} is the coexistence of observational proximity and target separation. Every event in the complete observed sample has probabilities differing by at most the fixed tolerance, while each prior places at least \(1-\beta\) probability on its respective side of the deterministic center. Many small baseline masses carry the signed perturbations; their aggregate effect yields separation at the missingness-and-dimension scale. The factor \(1-q\) records how that separation contracts as arrival becomes complete.

The positive separation constant \leanref{sym:c_beta}{\ensuremath{c_\beta}} in \cref{obj:thm:normalized-converse} may depend on the fixed tolerance and is uniform over the parameter triples satisfying the theorem's conditions. In the dimension-dominated regime \(g_{n,d,q}^{\,2}\ge n^{-1}\), the theorem supplies the prior comparison underlying the dimension terms in the estimation and honest-length bounds of \cref{obj:thm:point-frontier,obj:thm:interval-frontier}. Its two probability guarantees keep the observed-mixture comparison and the concentration of the population effects distinct, providing the inputs for the testing reductions.


\section{Testing reductions and proofs of the main results}

The normalized comparison in \cref{obj:thm:normalized-converse} supplies separated treatment effects under nearby observed-sample mixtures. The auxiliary results below translate that comparison into lower bounds for squared-error risk and honest expected interval length. A parametric subexperiment supplies the sampling floors throughout the parameter range, while the normalized priors supply the dimension contribution under their stated regime condition.

\subsection{Parametric sampling floors}

A one-label, complete-arrival Bernoulli subexperiment belongs to \(\mathcal M(d,q)\), the law class in \cref{obj:def:law-class}, for every \(d\ge1\) and \(q\in[1/2,1]\). Its sampling uncertainty supplies the parametric component of both statistical criteria. The underlying two-hypothesis testing comparison is standard \citep[Theorem~2.2]{Tsybakov2009}. The following statement records the resulting floors uniformly over the support bound and arrival floor.

\begin{lemmav}{lem:parametric-floor-infrastructure}[Parametric risk and length floors]
There exists a universal constant \(c_0>0\) such that, for every pair of positive integers \(n,d\) and every \(q\in[1/2,1]\), the minimax squared-error risk \(\mathfrak R^*_{n,d,q}\) defined in \cref{obj:def:point-risk} satisfies
\[
\mathfrak R^*_{n,d,q}\ge \frac{c_0}{n}.
\]
For each noncoverage level \(\alpha\in(0,1/2)\), there exists a constant \(c_{0,\alpha}>0\), depending only on \(\alpha\), such that, for every pair of positive integers \(n,d\) and every \(q\in[1/2,1]\), the minimax worst-law expected honest interval length \(\mathfrak L^*_{n,d,q,\alpha}\) defined in \cref{obj:def:length-risk} satisfies
\[
\mathfrak L^*_{n,d,q,\alpha}\ge \frac{c_{0,\alpha}}{\sqrt n}.
\]
\end{lemmav}

\begin{proof}[Proof of \cref{obj:lem:parametric-floor-infrastructure}]
Fix positive integers \(n,d\) and \(q\in[1/2,1]\).

\begin{enumerate}
\item \emph{Risk calibration.}
Define
\[
a_{\mathrm{risk}}
:=
\min\left\{\frac14,\frac{\sqrt{\log 2}}4\right\},
\qquad
u_{\mathrm{risk}}
:=
\frac{a_{\mathrm{risk}}}{\sqrt n},
\qquad
c_0:=\frac{a_{\mathrm{risk}}^2}{16}.
\]
Since \(\log 2>0\),
\[
0<a_{\mathrm{risk}}\le\frac14,
\qquad
a_{\mathrm{risk}}^2\le\frac{\log 2}{16},
\qquad
c_0>0.
\]
Moreover, \(n\ge1\) implies \(\sqrt n\ge1\) and \((\sqrt n)^2=n\). Consequently,
\[
0<u_{\mathrm{risk}}\le\frac14,
\qquad
n u_{\mathrm{risk}}^2=a_{\mathrm{risk}}^2,
\qquad
16n u_{\mathrm{risk}}^2\le\log 2.
\]
Thus the same positive constant \(a_{\mathrm{risk}}\) meets this budget for every positive sample size.
% lean: CausalSmith.Stat.MarNearcompleteFrontier.parametric_radius_budget

\item \emph{A one-cell pair.}
For each parameter \(\theta\in[0,1]\), define a full-data law \(P_\theta\) by
\[
\begin{gathered}
X=0,\qquad S(0)=S(1)=0,\qquad Y(0)=0,\qquad R=1
\quad\text{almost surely},\\
A\sim\operatorname{Bernoulli}(1/2),\qquad
Y(1)\sim\operatorname{Bernoulli}(\theta),\qquad
A\perp Y(1),\\
(S,Y)=(S(A),Y(A)),
\qquad
Q_\theta:=(P_\theta)_O .
\end{gathered}
\]
The label \(0\) is available because \(d\ge1\). The construction gives randomized treatment independence, balanced assignment, and consistency. Since \(R=1\), outcome arrival is conditionally independent of \(Y\), and every occupied cell has arrival probability one. Hence the requirements in \cref{obj:def:law-class} are satisfied:
\[
P_\theta\in\mathcal M(d,q),
\qquad
\tau(P_\theta)
=\mathbb E_{P_\theta}\{Y(1)-Y(0)\}
=\theta.
\]
The observed law has exactly the following possible nonzero masses:
\[
\begin{aligned}
Q_\theta\{(0,0,0,1,0)\}&=\frac12,\\
Q_\theta\{(0,1,0,1,0)\}&=\frac{1-\theta}{2},\\
Q_\theta\{(0,1,0,1,1)\}&=\frac{\theta}{2},
\end{aligned}
\]
and assigns mass zero to every other record.

Define the risk pair by
\[
P_{\mathrm{risk},0}:=P_{1/2},
\qquad
P_{\mathrm{risk},1}:=P_{1/2+u_{\mathrm{risk}}},
\qquad
Q_{\mathrm{risk},\iota}:=(P_{\mathrm{risk},\iota})_O
\quad(\iota\in\{0,1\}).
\]
Both parameters belong to \([0,1]\), and
\[
\tau(P_{\mathrm{risk},0})=\frac12,
\qquad
\tau(P_{\mathrm{risk},1})=\frac12+u_{\mathrm{risk}},
\qquad
Q_{\mathrm{risk},1}\ll Q_{\mathrm{risk},0}.
\]
For this finite, absolutely continuous pair, the chi-square divergence is
\[
\chi^2(Q_{\mathrm{risk},1}\Vert Q_{\mathrm{risk},0})
:=
\sum_{\substack{o\in O\\Q_{\mathrm{risk},0}\{o\}>0}}
\frac{\bigl(Q_{\mathrm{risk},1}\{o\}
             -Q_{\mathrm{risk},0}\{o\}\bigr)^2}
     {Q_{\mathrm{risk},0}\{o\}}.
\]
In sums over the entire space, a summand with zero denominator is assigned zero; absolute continuity makes its numerator zero as well. Expanding the square and using that both laws have total mass one gives
\[
\begin{aligned}
1+\chi^2(Q_{\mathrm{risk},1}\Vert Q_{\mathrm{risk},0})
&=
\sum_{o\in O}
\frac{Q_{\mathrm{risk},1}\{o\}^2}
     {Q_{\mathrm{risk},0}\{o\}}\\
&=
\frac12+
\left(\frac12-u_{\mathrm{risk}}\right)^2+
\left(\frac12+u_{\mathrm{risk}}\right)^2\\
&=1+2u_{\mathrm{risk}}^2.
\end{aligned}
\]
In particular,
\[
\chi^2(Q_{\mathrm{risk},1}\Vert Q_{\mathrm{risk},0})
=2u_{\mathrm{risk}}^2
\le16u_{\mathrm{risk}}^2.
\]
% lean: CausalSmith.Stat.MarNearcompleteFrontier.parametric_oneCell_pair

\item \emph{Product total variation.}
We use total variation with the convention
\[
\operatorname{TV}
 (Q_{\mathrm{risk},1}^{\otimes n},Q_{\mathrm{risk},0}^{\otimes n})
=
\sup_{\mathcal E\subseteq O^n}
\left|
Q_{\mathrm{risk},1}^{\otimes n}(\mathcal E)
-
Q_{\mathrm{risk},0}^{\otimes n}(\mathcal E)
\right|.
\]
Absolute continuity is preserved by these finite products. Factoring the finite sum of squared product masses divided by reference masses yields
\[
1+\chi^2
 (Q_{\mathrm{risk},1}^{\otimes n}\Vert Q_{\mathrm{risk},0}^{\otimes n})
=
\left\{
1+\chi^2(Q_{\mathrm{risk},1}\Vert Q_{\mathrm{risk},0})
\right\}^{n}.
\]
The inequality \(1+t\le e^t\) for \(t\ge0\), followed by the one-record bound, therefore gives
\[
\begin{aligned}
1+\chi^2
 (Q_{\mathrm{risk},1}^{\otimes n}\Vert Q_{\mathrm{risk},0}^{\otimes n})
&\le
\exp\!\left\{
n\chi^2(Q_{\mathrm{risk},1}\Vert Q_{\mathrm{risk},0})
\right\}\\
&\le \exp(16n u_{\mathrm{risk}}^2),
\end{aligned}
\]
so
\[
\chi^2
 (Q_{\mathrm{risk},1}^{\otimes n}\Vert Q_{\mathrm{risk},0}^{\otimes n})
\le \exp(16n u_{\mathrm{risk}}^2)-1.
\]
On a finite space, total variation equals half the sum of the absolute mass differences. Applying Cauchy--Schwarz to that sum, with the reference masses as weights, gives
\[
\begin{aligned}
\operatorname{TV}
 (Q_{\mathrm{risk},1}^{\otimes n},Q_{\mathrm{risk},0}^{\otimes n})
&\le
\frac12\sqrt{
\chi^2
 (Q_{\mathrm{risk},1}^{\otimes n}\Vert Q_{\mathrm{risk},0}^{\otimes n})
}\\
&\le
\frac12\sqrt{\exp(16n u_{\mathrm{risk}}^2)-1}.
\end{aligned}
\]
Finally, the calibrated budget implies
\[
\exp(16n u_{\mathrm{risk}}^2)\le e^{\log 2}=2,
\qquad
\operatorname{TV}
 (Q_{\mathrm{risk},1}^{\otimes n},Q_{\mathrm{risk},0}^{\otimes n})
\le\frac12.
\]
% lean: CausalSmith.Stat.MarNearcompleteFrontier.parametric_product_tv_at_radius

\item \emph{The squared-error lower bound.}
The estimator class in \cref{obj:synth_6} is nonempty, since it contains
\[
T_{\mathrm{zero}}(o):=0
\qquad(o\in O^n).
\]
Fix \(T\in\mathcal T_n\), and define its two error events by
\[
\mathcal E_{\mathrm{risk},\iota}
:=
\left\{
o\in O^n:
\frac{u_{\mathrm{risk}}}{2}
\le
\left|T(o)-\tau(P_{\mathrm{risk},\iota})\right|
\right\},
\qquad \iota\in\{0,1\}.
\]
The targets are separated by \(u_{\mathrm{risk}}\). The triangle inequality therefore gives
\[
\mathcal E_{\mathrm{risk},0}^{\,c}
\subseteq \mathcal E_{\mathrm{risk},1}.
\]
Using this inclusion and the total-variation bound,
\[
\begin{aligned}
&Q_{\mathrm{risk},0}^{\otimes n}(\mathcal E_{\mathrm{risk},0})
+
Q_{\mathrm{risk},1}^{\otimes n}(\mathcal E_{\mathrm{risk},1})\\
&\qquad\ge
Q_{\mathrm{risk},0}^{\otimes n}(\mathcal E_{\mathrm{risk},0})
+
Q_{\mathrm{risk},0}^{\otimes n}(\mathcal E_{\mathrm{risk},1})
-
\operatorname{TV}
 (Q_{\mathrm{risk},0}^{\otimes n},Q_{\mathrm{risk},1}^{\otimes n})\\
&\qquad\ge\frac12.
\end{aligned}
\]
Consequently,
\[
\max_{\iota\in\{0,1\}}
Q_{\mathrm{risk},\iota}^{\otimes n}(\mathcal E_{\mathrm{risk},\iota})
\ge\frac14.
\]
For each \(\iota\in\{0,1\}\), nonnegativity of squared error gives the threshold bound
\[
\left(\frac{u_{\mathrm{risk}}}{2}\right)^2
Q_{\mathrm{risk},\iota}^{\otimes n}(\mathcal E_{\mathrm{risk},\iota})
\le
\mathbb E_{Q_{\mathrm{risk},\iota}^{\otimes n}}
\left[
\{T-\tau(P_{\mathrm{risk},\iota})\}^2
\right].
\]
If the error probability under the first law is at most that under the second, then
\[
\mathbb E_{Q_{\mathrm{risk},1}^{\otimes n}}
\left[
\{T-\tau(P_{\mathrm{risk},1})\}^2
\right]
\ge
\left(\frac{u_{\mathrm{risk}}}{2}\right)^2\frac14
=
\frac{u_{\mathrm{risk}}^2}{16}.
\]
If that ordering is reversed, the corresponding bound is
\[
\mathbb E_{Q_{\mathrm{risk},0}^{\otimes n}}
\left[
\{T-\tau(P_{\mathrm{risk},0})\}^2
\right]
\ge
\left(\frac{u_{\mathrm{risk}}}{2}\right)^2\frac14
=
\frac{u_{\mathrm{risk}}^2}{16}.
\]
Thus, in either case,
\[
\max_{\iota\in\{0,1\}}
\mathbb E_{Q_{\mathrm{risk},\iota}^{\otimes n}}
\left[
\{T-\tau(P_{\mathrm{risk},\iota})\}^2
\right]
\ge\frac{u_{\mathrm{risk}}^2}{16}.
\]

For every \(P\in\mathcal M(d,q)\), binary potential outcomes imply
\[
-1\le Y(1)-Y(0)\le1,
\qquad
-1\le\tau(P)\le1.
\]
Since \(T\) also takes values in \([-1,1]\),
\[
0\le
\mathbb E_{P_O^{\otimes n}}\{T-\tau(P)\}^2
\le4.
\]
The worst-law supremum therefore exists as a finite real number and dominates the risks at both constructed laws. Taking the estimator infimum in \cref{obj:def:point-risk} gives
\[
\mathfrak R^*_{n,d,q}
\ge\frac{u_{\mathrm{risk}}^2}{16}
=
\frac{a_{\mathrm{risk}}^2}{16n}
=
\frac{c_0}{n}.
\]
% lean: CausalSmith.Stat.MarNearcompleteFrontier.parametric_point_pair_lower

\item \emph{Coverage-dependent calibration.}
Fix \(\alpha\in(0,1/2)\), and define
\[
b_{\mathrm{cov}}:=1-2\alpha,
\qquad
a_{\mathrm{len}}
:=
\min\left\{
\frac14,\frac{\sqrt{\log(1+b_{\mathrm{cov}}^2)}}4
\right\},
\qquad
c_{0,\alpha}:=\frac{a_{\mathrm{len}}b_{\mathrm{cov}}}{2}.
\]
Here
\[
b_{\mathrm{cov}}>0,
\qquad
\log(1+b_{\mathrm{cov}}^2)>0,
\qquad
0<a_{\mathrm{len}}\le\frac14,
\qquad
c_{0,\alpha}>0.
\]
The minimum also gives
\[
a_{\mathrm{len}}^2
\le\frac{\log(1+b_{\mathrm{cov}}^2)}{16}.
\]
Define
\[
u_{\mathrm{len}}:=\frac{a_{\mathrm{len}}}{\sqrt n}.
\]
As before, \(\sqrt n\ge1\) and \((\sqrt n)^2=n\), so
\[
0<u_{\mathrm{len}}\le\frac14,
\qquad
n u_{\mathrm{len}}^2=a_{\mathrm{len}}^2,
\qquad
16n u_{\mathrm{len}}^2\le\log(1+b_{\mathrm{cov}}^2).
\]
Both \(a_{\mathrm{len}}\) and \(c_{0,\alpha}\) depend only on \(\alpha\).
% lean: CausalSmith.Stat.MarNearcompleteFrontier.parametric_radius_budget_at_tolerance

\item \emph{The interval pair and its total variation.}
Using the same one-cell family, define
\[
P_{\mathrm{len},0}:=P_{1/2},
\qquad
P_{\mathrm{len},1}:=P_{1/2+u_{\mathrm{len}}},
\qquad
Q_{\mathrm{len},\iota}:=(P_{\mathrm{len},\iota})_O
\quad(\iota\in\{0,1\}).
\]
The construction and the same three-atom calculation give
\[
\begin{gathered}
P_{\mathrm{len},0},P_{\mathrm{len},1}\in\mathcal M(d,q),\\
\tau(P_{\mathrm{len},0})=\frac12,
\qquad
\tau(P_{\mathrm{len},1})=\frac12+u_{\mathrm{len}},\\
Q_{\mathrm{len},1}\ll Q_{\mathrm{len},0},
\qquad
\chi^2(Q_{\mathrm{len},1}\Vert Q_{\mathrm{len},0})
=2u_{\mathrm{len}}^2\le16u_{\mathrm{len}}^2.
\end{gathered}
\]
The product calculation above applies to this pair. The new budget gives
\[
\exp(16n u_{\mathrm{len}}^2)
\le
\exp\{\log(1+b_{\mathrm{cov}}^2)\}
=
1+b_{\mathrm{cov}}^2.
\]
Using symmetry of total variation and \(b_{\mathrm{cov}}>0\),
\[
\begin{aligned}
\operatorname{TV}
 (Q_{\mathrm{len},0}^{\otimes n},Q_{\mathrm{len},1}^{\otimes n})
&=
\operatorname{TV}
 (Q_{\mathrm{len},1}^{\otimes n},Q_{\mathrm{len},0}^{\otimes n})\\
&\le\frac12\sqrt{\exp(16n u_{\mathrm{len}}^2)-1}\\
&\le\frac12\sqrt{b_{\mathrm{cov}}^2}
=\frac{b_{\mathrm{cov}}}{2}
=\frac{1-2\alpha}{2}.
\end{aligned}
\]
% lean: CausalSmith.Stat.MarNearcompleteFrontier.parametric_product_tv_at_tolerance_radius

\item \emph{Coverage forces expected length.}
Fix \(I\in\mathcal C_\alpha(n,d,q)\), with endpoints as in \cref{obj:def:interval-class}, and define
\[
\mathcal E_{\mathrm{cov},\iota}
:=
\left\{
o\in O^n:
\tau(P_{\mathrm{len},\iota})\in[L(o),U(o)]
\right\},
\qquad \iota\in\{0,1\}.
\]
Honesty gives
\[
Q_{\mathrm{len},0}^{\otimes n}(\mathcal E_{\mathrm{cov},0})
\ge1-\alpha,
\qquad
Q_{\mathrm{len},1}^{\otimes n}(\mathcal E_{\mathrm{cov},1})
\ge1-\alpha.
\]
All these events are measurable because the sample space is finite. The definition of total variation gives
\[
\left|
Q_{\mathrm{len},0}^{\otimes n}(\mathcal E_{\mathrm{cov},1})
-
Q_{\mathrm{len},1}^{\otimes n}(\mathcal E_{\mathrm{cov},1})
\right|
\le\frac{b_{\mathrm{cov}}}{2},
\]
and hence
\[
Q_{\mathrm{len},0}^{\otimes n}(\mathcal E_{\mathrm{cov},1})
\ge1-\alpha-\frac{b_{\mathrm{cov}}}{2}.
\]
Taking complements,
\[
Q_{\mathrm{len},0}^{\otimes n}(\mathcal E_{\mathrm{cov},0}^{\,c})
\le\alpha,
\qquad
Q_{\mathrm{len},0}^{\otimes n}(\mathcal E_{\mathrm{cov},1}^{\,c})
\le\alpha+\frac{b_{\mathrm{cov}}}{2}.
\]
The union bound therefore yields
\[
\begin{aligned}
Q_{\mathrm{len},0}^{\otimes n}
 (\mathcal E_{\mathrm{cov},0}\cap\mathcal E_{\mathrm{cov},1})
&=
1-
Q_{\mathrm{len},0}^{\otimes n}
 (\mathcal E_{\mathrm{cov},0}^{\,c}
  \cup\mathcal E_{\mathrm{cov},1}^{\,c})\\
&\ge1-2\alpha-\frac{b_{\mathrm{cov}}}{2}
=\frac{b_{\mathrm{cov}}}{2}.
\end{aligned}
\]
On the intersection, the interval contains both \(1/2\) and \(1/2+u_{\mathrm{len}}\), so
\[
L(o)\le\frac12,
\qquad
U(o)\ge\frac12+u_{\mathrm{len}},
\qquad
U(o)-L(o)\ge u_{\mathrm{len}}.
\]
Consequently,
\[
Q_{\mathrm{len},0}^{\otimes n}
 (\mathcal E_{\mathrm{cov},0}\cap\mathcal E_{\mathrm{cov},1})
\le
Q_{\mathrm{len},0}^{\otimes n}
 \{o:u_{\mathrm{len}}\le U(o)-L(o)\}.
\]
Since interval length is nonnegative, its expectation satisfies
\[
u_{\mathrm{len}}\,
Q_{\mathrm{len},0}^{\otimes n}
 \{o:u_{\mathrm{len}}\le U(o)-L(o)\}
\le
\mathbb E_{Q_{\mathrm{len},0}^{\otimes n}}[U-L].
\]
Combining these bounds,
\[
\begin{aligned}
\frac{u_{\mathrm{len}}b_{\mathrm{cov}}}{2}
&\le
u_{\mathrm{len}}\,
Q_{\mathrm{len},0}^{\otimes n}
 (\mathcal E_{\mathrm{cov},0}\cap\mathcal E_{\mathrm{cov},1})\\
&\le
u_{\mathrm{len}}\,
Q_{\mathrm{len},0}^{\otimes n}
 \{o:u_{\mathrm{len}}\le U(o)-L(o)\}\\
&\le
\mathbb E_{Q_{\mathrm{len},0}^{\otimes n}}[U-L].
\end{aligned}
\]

The honest interval class is nonempty: the procedure
\[
I_{\mathrm{full}}(o):=[-1,1]
\qquad(o\in O^n)
\]
covers every target because \(\tau(P)\in[-1,1]\). For every admissible interval and law,
\[
0\le U(o)-L(o)\le2,
\qquad
0\le\mathbb E_{P_O^{\otimes n}}[U-L]\le2.
\]
Thus the worst-law expected length is finite and dominates its value at \(P_{\mathrm{len},0}\). Taking the interval infimum in \cref{obj:def:length-risk} gives
\[
\mathfrak L^*_{n,d,q,\alpha}
\ge
\frac{u_{\mathrm{len}}b_{\mathrm{cov}}}{2}
=
\frac{a_{\mathrm{len}}b_{\mathrm{cov}}}{2\sqrt n}
=
\frac{c_{0,\alpha}}{\sqrt n}.
\]
Together with the squared-error bound, this proves both assertions.
% lean: CausalSmith.Stat.MarNearcompleteFrontier.parametric_interval_pair_lower
% lean: CausalSmith.Stat.MarNearcompleteFrontier.parametric_floor
\end{enumerate}
\end{proof}

These bounds retain ordinary sampling uncertainty even at complete arrival. For estimation, the floor is universal; for honest intervals, its constant accommodates the fixed coverage requirement. They supply the lower bounds when the missingness-and-dimension contribution lies below the sampling scale and remain available in the regime covered by the normalized converse.

\subsection{From normalized priors to statistical lower bounds}

For each fixed mixture tolerance \(\beta\in(0,1/4)\), choose a positive separation constant \(c_\beta\) satisfying \cref{obj:thm:normalized-converse}. Fix this choice before varying \(n,d,q\); its uniformity is over the parameter triples satisfying that statement, including the condition \(g_{n,d,q}^{\,2}\ge n^{-1}\).

The testing-error comparison applies to the two complete observed-sample mixtures as probability laws \citep[Theorem~2.2]{Tsybakov2009}. The model-specific input is their construction and target separation in \cref{obj:def:paired-prior-handle,obj:thm:normalized-converse}. For squared-error risk, separation is measured in squared units. For honest connected intervals, coverage of separated target values links the same comparison to expected length. The next statement fixes the tolerances for these two uses and records their numerical losses.

\begin{lemmav}{lem:prior-reduction}[Prior reduction and lower bounds]
Let \(c_\beta\) denote the fixed separation constant selected from the normalized converse at mixture tolerance \(\beta\in(0,1/4)\), and set \(\beta_0=1/16\). Write
\[
\ell_n=\log(e+n),
\qquad
g_{n,d,q}=(1-q)\min\left\{1,\frac{d}{n\ell_n}\right\}.
\]
Use the minimax risks and rate in \cref{obj:def:point-risk,obj:def:length-risk,obj:def:rate}.

There exist universal constants \(c>0\) and \(c_0>0\) such that, for every \(\alpha\in(0,1/2)\), there exist constants \(c_\alpha>0\) and \(c_{0,\alpha}>0\) with the following guarantees. Set
\[
\beta_\alpha=\frac{1-2\alpha}{8}.
\]
For every \(n,d,q\) satisfying
\begin{itemize}
\item \textbf{(Sample size.)} \(n\) is an integer with \(n\ge1\).
\item \textbf{(Support bound.)} \(d\) is an integer with \(d\ge1\).
\item \textbf{(Arrival floor.)} \(q\in[1/2,1]\).
\end{itemize}
whenever \(g_{n,d,q}^2\ge n^{-1}\), both bounds
\[
\mathfrak R^*_{n,d,q}
\ge \frac{13}{32}c_{\beta_0}^{\,2}g_{n,d,q}^{\,2},
\qquad
\mathfrak L^*_{n,d,q,\alpha}
\ge \frac54 c_{\beta_\alpha}(1-2\alpha)g_{n,d,q}
\]
hold. Throughout the stated parameter range, the parametric bounds are
\[
\mathfrak R^*_{n,d,q}\ge\frac{c_0}{n},
\qquad
\mathfrak L^*_{n,d,q,\alpha}\ge\frac{c_{0,\alpha}}{\sqrt n}.
\]
The combined guarantees are
\(\mathfrak R^*_{n,d,q}\ge c\,r_{n,d,q}\) and
\(\mathfrak L^*_{n,d,q,\alpha}\ge c_\alpha\sqrt{r_{n,d,q}}\).
\end{lemmav}

\begin{proof}[Proof of \cref{obj:lem:prior-reduction}]
\begin{enumerate}
\item Choose the universal constant \(c_{\mathrm{floor}}>0\) supplied by
\cref{obj:lem:parametric-floor-infrastructure}, so that
\[
\mathfrak R^*_{n,d,q}\ge \frac{c_{\mathrm{floor}}}{n}
\]
throughout the stated parameter range. Define
\[
c_{\mathrm{sep,pt}}
=\frac{13}{32}c_{\beta_0}^{\,2},
\qquad
c=\frac12\min\{c_{\mathrm{floor}},c_{\mathrm{sep,pt}}\}.
\]
Since \(\beta_0=1/16\in(0,1/4)\), the selected constant in
\cref{obj:thm:normalized-converse} satisfies \(c_{\beta_0}>0\).
Consequently,
\[
c_{\mathrm{sep,pt}}>0,
\qquad c>0.
\]

Fix \(\alpha\in(0,1/2)\). Choose \(c_{0,\alpha}>0\) from
\cref{obj:lem:parametric-floor-infrastructure}, giving
\[
\mathfrak L^*_{n,d,q,\alpha}
\ge \frac{c_{0,\alpha}}{\sqrt n}
\]
throughout the same parameter range. Define
\[
\beta_\alpha=\frac{1-2\alpha}{8},
\qquad
c_{\mathrm{sep,len},\alpha}
=\frac54c_{\beta_\alpha}(1-2\alpha),
\qquad
c_\alpha
=\frac12\min\{c_{0,\alpha},c_{\mathrm{sep,len},\alpha}\}.
\]
The inequalities
\[
0<1-2\alpha<1,
\qquad
0<\beta_\alpha<\frac18<\frac14
\]
ensure \(c_{\beta_\alpha}>0\), hence
\[
c_{\mathrm{sep,len},\alpha}>0,
\qquad c_\alpha>0.
\]
The constants \(c,c_{\mathrm{floor}}\) are universal, and \(c_\alpha,c_{0,\alpha}\)
depend only on \(\alpha\).
% lean: CausalSmith.Stat.MarNearcompleteFrontier.prior_reduction

\item Fix admissible \(n,d,q\). The scale is nonnegative because
\[
1-q\ge0,\qquad
\ell_n=\log(e+n)>0,\qquad
\frac{d}{n\ell_n}\ge0,
\qquad
g_{n,d,q}=(1-q)\min\left\{1,\frac{d}{n\ell_n}\right\}\ge0.
\]
Suppose \(g_{n,d,q}^{\,2}\ge n^{-1}\), and define
\[
a_{\mathrm{pt}}=c_{\beta_0}g_{n,d,q}\ge0.
\]
By \cref{obj:thm:normalized-converse}, there are probability priors
\(\Pi_{-,\mathrm{pt}},\Pi_{+,\mathrm{pt}}\) and a center
\(m_{0,\mathrm{pt}}\) satisfying
\[
\Pi_{\sigma,\mathrm{pt}}\bigl(\mathcal M(d,q)\bigr)=1,
\qquad
m_{0,\mathrm{pt}}\in\mathbb R,
\qquad \sigma\in\{-,+\}.
\]
Define their sample-mixture laws by
\[
Q_{\sigma,\mathrm{pt}}^{(n)}
=\int_{\mathcal M(d,q)}P_O^{\otimes n}\,
  \Pi_{\sigma,\mathrm{pt}}(dP),
\qquad \sigma\in\{-,+\}.
\]
The supplied bounds are
\[
\operatorname{TV}\!\left(
Q_{-,\mathrm{pt}}^{(n)},Q_{+,\mathrm{pt}}^{(n)}
\right)\le\beta_0,
\]
\[
\Pi_{-,\mathrm{pt}}\{\tau(P)\le m_{0,\mathrm{pt}}-a_{\mathrm{pt}}\}
\ge1-\beta_0,
\qquad
\Pi_{+,\mathrm{pt}}\{m_{0,\mathrm{pt}}+a_{\mathrm{pt}}\le\tau(P)\}
\ge1-\beta_0.
\]

For an arbitrary estimator \(T\in\mathcal T_n\), define
\[
E_{\mathrm{pt}}
=\{o\in O^n:m_{0,\mathrm{pt}}\le T(o)\},
\]
\[
\mathcal B_{\sigma,\mathrm{pt}}(T)
=\int_{\mathcal M(d,q)}
  \mathbb E_P\!\left[(T(O^n)-\tau(P))^2\right]\,
  \Pi_{\sigma,\mathrm{pt}}(dP),
\qquad \sigma\in\{-,+\},
\]
and the target-failure indicators
\[
D_{-,\mathrm{pt}}(P)
=\mathbf1\{\tau(P)>m_{0,\mathrm{pt}}-a_{\mathrm{pt}}\},
\qquad
D_{+,\mathrm{pt}}(P)
=\mathbf1\{\tau(P)<m_{0,\mathrm{pt}}+a_{\mathrm{pt}}\}.
\]
On the left target region, \(o\in E_{\mathrm{pt}}\) gives
\[
T(o)-\tau(P)\ge a_{\mathrm{pt}},
\qquad
(T(o)-\tau(P))^2\ge a_{\mathrm{pt}}^2.
\]
On the right target region, \(o\in E_{\mathrm{pt}}^c\) gives
\[
\tau(P)-T(o)\ge a_{\mathrm{pt}},
\qquad
(T(o)-\tau(P))^2\ge a_{\mathrm{pt}}^2.
\]
Integrating these pointwise bounds over the corresponding sample events
gives the good-target risk bounds. On each failure region, the sample-event
probability is at most one and squared loss is nonnegative. Thus, for every
\(P\in\mathcal M(d,q)\),
\[
\begin{aligned}
a_{\mathrm{pt}}^2
\left[P_O^{\otimes n}(E_{\mathrm{pt}})-D_{-,\mathrm{pt}}(P)\right]
&\le \mathbb E_P[(T(O^n)-\tau(P))^2],\\
a_{\mathrm{pt}}^2
\left[P_O^{\otimes n}(E_{\mathrm{pt}}^c)-D_{+,\mathrm{pt}}(P)\right]
&\le \mathbb E_P[(T(O^n)-\tau(P))^2].
\end{aligned}
\]
The prior failure probabilities satisfy
\[
\int D_{\sigma,\mathrm{pt}}(P)\,
  \Pi_{\sigma,\mathrm{pt}}(dP)\le\beta_0,
\qquad \sigma\in\{-,+\}.
\]
Integration over the two priors therefore yields
\[
\begin{aligned}
\mathcal B_{-,\mathrm{pt}}(T)
&\ge a_{\mathrm{pt}}^2
 \left[Q_{-,\mathrm{pt}}^{(n)}(E_{\mathrm{pt}})-\beta_0\right],\\
\mathcal B_{+,\mathrm{pt}}(T)
&\ge a_{\mathrm{pt}}^2
 \left[Q_{+,\mathrm{pt}}^{(n)}(E_{\mathrm{pt}}^c)-\beta_0\right].
\end{aligned}
\]
Total variation gives
\[
\begin{aligned}
Q_{-,\mathrm{pt}}^{(n)}(E_{\mathrm{pt}})
+Q_{+,\mathrm{pt}}^{(n)}(E_{\mathrm{pt}}^c)
&=1+Q_{-,\mathrm{pt}}^{(n)}(E_{\mathrm{pt}})
     -Q_{+,\mathrm{pt}}^{(n)}(E_{\mathrm{pt}})\\
&\ge1-\beta_0.
\end{aligned}
\]
Adding the two risk bounds and substituting \(\beta_0=1/16\), we obtain
\[
\mathcal B_{-,\mathrm{pt}}(T)+\mathcal B_{+,\mathrm{pt}}(T)
\ge (1-3\beta_0)a_{\mathrm{pt}}^2
=\frac{13}{16}c_{\beta_0}^{\,2}g_{n,d,q}^{\,2}.
\]
Each prior average is at most the worst-law risk, since its prior is
supported on \(\mathcal M(d,q)\). Hence
\[
2\sup_{P\in\mathcal M(d,q)}
  \mathbb E_P[(T(O^n)-\tau(P))^2]
\ge
\mathcal B_{-,\mathrm{pt}}(T)+\mathcal B_{+,\mathrm{pt}}(T)
\ge 2c_{\mathrm{sep,pt}}g_{n,d,q}^{\,2}.
\]
The estimator and target both lie in \([-1,1]\), so the squared loss lies
in \([0,4]\); these bounds justify the prior integrations and the finite
worst-law risks. The estimator class contains the constant zero estimator.
Taking its infimum as in \cref{obj:def:point-risk} proves
\[
\mathfrak R^*_{n,d,q}
\ge c_{\mathrm{sep,pt}}g_{n,d,q}^{\,2}
=\frac{13}{32}c_{\beta_0}^{\,2}g_{n,d,q}^{\,2}.
\]
% lean: CausalSmith.Stat.MarNearcompleteFrontier.priorPointRisk_separation

\item Continue under \(g_{n,d,q}^{\,2}\ge n^{-1}\), and define
\[
a_{\mathrm{len}}=c_{\beta_\alpha}g_{n,d,q}\ge0.
\]
Apply \cref{obj:thm:normalized-converse} at \(\beta_\alpha\).
It supplies probability priors and a center with
\[
\Pi_{\sigma,\mathrm{len}}\bigl(\mathcal M(d,q)\bigr)=1,
\qquad
m_{0,\mathrm{len}}\in\mathbb R,
\qquad \sigma\in\{-,+\}.
\]
Define
\[
Q_{\sigma,\mathrm{len}}^{(n)}
=\int_{\mathcal M(d,q)}P_O^{\otimes n}\,
  \Pi_{\sigma,\mathrm{len}}(dP),
\qquad \sigma\in\{-,+\}.
\]
Then
\[
\operatorname{TV}\!\left(
Q_{-,\mathrm{len}}^{(n)},Q_{+,\mathrm{len}}^{(n)}
\right)\le\beta_\alpha,
\]
\[
\Pi_{-,\mathrm{len}}\{\tau(P)\le m_{0,\mathrm{len}}-a_{\mathrm{len}}\}
\ge1-\beta_\alpha,
\qquad
\Pi_{+,\mathrm{len}}\{m_{0,\mathrm{len}}+a_{\mathrm{len}}\le\tau(P)\}
\ge1-\beta_\alpha.
\]

For an arbitrary honest interval, write
\[
I\in\mathcal C_\alpha(n,d,q),
\qquad
I(o)=[L_I(o),U_I(o)],
\qquad o\in O^n,
\]
and define the endpoint events
\[
E_{\mathrm L,\alpha}
=\{o\in O^n:L_I(o)\le m_{0,\mathrm{len}}-a_{\mathrm{len}}\},
\qquad
E_{\mathrm R,\alpha}
=\{o\in O^n:m_{0,\mathrm{len}}+a_{\mathrm{len}}\le U_I(o)\}.
\]
Also define
\[
D_{-,\mathrm{len}}(P)
=\mathbf1\{\tau(P)>m_{0,\mathrm{len}}-a_{\mathrm{len}}\},
\qquad
D_{+,\mathrm{len}}(P)
=\mathbf1\{\tau(P)<m_{0,\mathrm{len}}+a_{\mathrm{len}}\}.
\]
On a left-target law, coverage implies
\(L_I(o)\le\tau(P)\le m_{0,\mathrm{len}}-a_{\mathrm{len}}\);
on a right-target law, it implies
\(m_{0,\mathrm{len}}+a_{\mathrm{len}}\le\tau(P)\le U_I(o)\).
Thus honesty from \cref{obj:def:interval-class} gives the respective
sample-event probabilities at least \(1-\alpha\).
On a failure law, \(1-\alpha-D_{\sigma,\mathrm{len}}(P)=-\alpha\le0\).
Consequently, for every \(P\in\mathcal M(d,q)\),
\[
\begin{aligned}
P_O^{\otimes n}(E_{\mathrm L,\alpha})
&\ge1-\alpha-D_{-,\mathrm{len}}(P),\\
P_O^{\otimes n}(E_{\mathrm R,\alpha})
&\ge1-\alpha-D_{+,\mathrm{len}}(P).
\end{aligned}
\]
Integrating and using the target-failure bounds gives
\[
\begin{aligned}
Q_{-,\mathrm{len}}^{(n)}(E_{\mathrm L,\alpha})
&\ge1-\alpha-\int D_{-,\mathrm{len}}\,d\Pi_{-,\mathrm{len}}
 \ge1-\alpha-\beta_\alpha,\\
Q_{+,\mathrm{len}}^{(n)}(E_{\mathrm R,\alpha})
&\ge1-\alpha-\int D_{+,\mathrm{len}}\,d\Pi_{+,\mathrm{len}}
 \ge1-\alpha-\beta_\alpha.
\end{aligned}
\]
The total-variation bound transfers the first estimate:
\[
Q_{+,\mathrm{len}}^{(n)}(E_{\mathrm L,\alpha})
\ge Q_{-,\mathrm{len}}^{(n)}(E_{\mathrm L,\alpha})-\beta_\alpha
\ge1-\alpha-2\beta_\alpha.
\]
Using the union-intersection identity and the fact that a union has
probability at most one, we obtain
\[
\begin{aligned}
Q_{+,\mathrm{len}}^{(n)}
 (E_{\mathrm L,\alpha}\cap E_{\mathrm R,\alpha})
&=
Q_{+,\mathrm{len}}^{(n)}(E_{\mathrm L,\alpha})
+Q_{+,\mathrm{len}}^{(n)}(E_{\mathrm R,\alpha})
-Q_{+,\mathrm{len}}^{(n)}
 (E_{\mathrm L,\alpha}\cup E_{\mathrm R,\alpha})\\
&\ge1-2\alpha-3\beta_\alpha.
\end{aligned}
\]
On the intersection, the endpoint inequalities imply
\[
U_I(o)-L_I(o)
\ge (m_{0,\mathrm{len}}+a_{\mathrm{len}})
   -(m_{0,\mathrm{len}}-a_{\mathrm{len}})
=2a_{\mathrm{len}}.
\]
Since interval length is nonnegative everywhere,
\[
2a_{\mathrm{len}}\,
\mathbf1_{E_{\mathrm L,\alpha}\cap E_{\mathrm R,\alpha}}(o)
\le U_I(o)-L_I(o).
\]
Integration under the plus mixture, followed by its defining mixture
identity, therefore gives
\[
\begin{aligned}
2a_{\mathrm{len}}\,
Q_{+,\mathrm{len}}^{(n)}
 (E_{\mathrm L,\alpha}\cap E_{\mathrm R,\alpha})
&\le
\int_{O^n}(U_I(o)-L_I(o))\,Q_{+,\mathrm{len}}^{(n)}(do)\\
&=
\int_{\mathcal M(d,q)}
\mathbb E_P[U_I(O^n)-L_I(O^n)]\,\Pi_{+,\mathrm{len}}(dP)\\
&\le
\sup_{P\in\mathcal M(d,q)}
\mathbb E_P[U_I(O^n)-L_I(O^n)].
\end{aligned}
\]
Multiplication of the intersection bound by \(2a_{\mathrm{len}}\ge0\)
and substitution of \(\beta_\alpha\) yield
\[
\begin{aligned}
\sup_{P\in\mathcal M(d,q)}
\mathbb E_P[U_I(O^n)-L_I(O^n)]
&\ge2a_{\mathrm{len}}(1-2\alpha-3\beta_\alpha)\\
&=2a_{\mathrm{len}}\frac58(1-2\alpha)\\
&=\frac54c_{\beta_\alpha}(1-2\alpha)g_{n,d,q}.
\end{aligned}
\]
Lengths lie in \([0,2]\), justifying these integrations, and the constant
interval \([-1,1]\) is honest because \(\tau(P)\in[-1,1]\).
Taking the infimum in \cref{obj:def:length-risk} proves
\[
\mathfrak L^*_{n,d,q,\alpha}
\ge c_{\mathrm{sep,len},\alpha}g_{n,d,q}
=\frac54c_{\beta_\alpha}(1-2\alpha)g_{n,d,q}.
\]
% lean: CausalSmith.Stat.MarNearcompleteFrontier.priorIntervalRisk_separation

\item We now combine the point-risk separation bound with the parametric
floor. By \cref{obj:def:rate},
\[
r_{n,d,q}=\frac1n+g_{n,d,q}^{\,2}.
\]
If \(g_{n,d,q}^{\,2}\ge n^{-1}\), then
\[
c\,r_{n,d,q}
\le 2c\,g_{n,d,q}^{\,2}
\le c_{\mathrm{sep,pt}}g_{n,d,q}^{\,2}
\le\mathfrak R^*_{n,d,q},
\]
where \(2c=\min\{c_{\mathrm{floor}},c_{\mathrm{sep,pt}}\}\le c_{\mathrm{sep,pt}}\).
If \(g_{n,d,q}^{\,2}<n^{-1}\), then
\[
c\,r_{n,d,q}
\le \frac{2c}{n}
\le \frac{c_{\mathrm{floor}}}{n}
\le\mathfrak R^*_{n,d,q},
\]
using \(2c\le c_{\mathrm{floor}}\) and the floor from
\cref{obj:lem:parametric-floor-infrastructure}.
Thus the combined point-risk bound holds throughout the parameter range.
% lean: CausalSmith.Stat.MarNearcompleteFrontier.pointRate_from_floor_and_separation

\item For the length bound, \(n\ge1\) ensures
\[
\sqrt n>0,
\qquad
(\sqrt n)^2=n,
\qquad
\frac1n=\left(\frac1{\sqrt n}\right)^2.
\]
Also \(r_{n,d,q}\ge0\), so
\[
\sqrt{r_{n,d,q}}\ge0,
\qquad
\bigl(\sqrt{r_{n,d,q}}\bigr)^2=r_{n,d,q}.
\]
If \(g_{n,d,q}^{\,2}\ge n^{-1}\), then
\[
\bigl(\sqrt{r_{n,d,q}}\bigr)^2
=\frac1n+g_{n,d,q}^{\,2}
\le2g_{n,d,q}^{\,2}
\le(2g_{n,d,q})^2.
\]
Both square roots being compared are nonnegative, hence
\[
\sqrt{r_{n,d,q}}\le2g_{n,d,q}.
\]
It follows that
\[
c_\alpha\sqrt{r_{n,d,q}}
\le2c_\alpha g_{n,d,q}
\le c_{\mathrm{sep,len},\alpha}g_{n,d,q}
\le\mathfrak L^*_{n,d,q,\alpha}.
\]
If \(g_{n,d,q}^{\,2}<n^{-1}\), then
\[
\bigl(\sqrt{r_{n,d,q}}\bigr)^2
=\frac1n+g_{n,d,q}^{\,2}
\le\frac2n
\le\left(\frac2{\sqrt n}\right)^2,
\]
and nonnegativity gives
\[
\sqrt{r_{n,d,q}}\le\frac2{\sqrt n}.
\]
Therefore
\[
c_\alpha\sqrt{r_{n,d,q}}
\le\frac{2c_\alpha}{\sqrt n}
=\frac{\min\{c_{0,\alpha},c_{\mathrm{sep,len},\alpha}\}}{\sqrt n}
\le\frac{c_{0,\alpha}}{\sqrt n}
\le\mathfrak L^*_{n,d,q,\alpha},
\]
using \cref{obj:lem:parametric-floor-infrastructure}.
These two cases cover every admissible parameter choice, including
\(q=1\), for which \(g_{n,d,q}=0\) and the second case applies.
Together with the separation bounds and the parametric floors established
above, this proves all the asserted guarantees.
% lean: CausalSmith.Stat.MarNearcompleteFrontier.intervalRate_from_floor_and_separation
\end{enumerate}
\end{proof}

The fixed tolerance \(\beta_0\) yields a universal squared-risk constant. By contrast, the coverage-dependent tolerance \(\beta_\alpha\) allocates the mixture-comparison error relative to the positive coverage margin \(1-2\alpha\). For each fixed \(\alpha\in(0,1/2)\), it belongs to the tolerance range required by \cref{obj:thm:normalized-converse}, and the resulting length constant depends only on \(\alpha\).

The displayed coefficients retain the losses associated with observed-mixture proximity and the prior probabilities of target separation. Their role is to make the conversion from the normalized converse to the statistical criteria explicit. Connectedness is relevant to the length criterion because an interval containing separated values spans their distance.

The dimension bounds retain the condition \(g_{n,d,q}^{\,2}\ge n^{-1}\) from \cref{obj:thm:normalized-converse}. The parametric floors in \cref{obj:lem:parametric-floor-infrastructure} supply the sampling contribution throughout the parameter range, including \(q=1\). Together, these inputs give the combined guarantees for the benchmark in \cref{obj:def:rate}, with universal constants for squared risk and constants depending only on fixed noncoverage for expected length.

\subsection{Completion of the main comparisons}

The upper bounds come from \cref{obj:lem:upper-risk,obj:lem:honest-interval-construction}; the normalized-prior reduction and parametric floors in \cref{obj:lem:prior-reduction,obj:lem:parametric-floor-infrastructure} give the matching lower bounds. The formal proofs of the two frontier theorems and the complete-arrival endpoint appear in the following proof appendix, so no second derivation is repeated here.


\section{Verification note}

The theorem statements and proofs are machine-checked in Lean 4 under the frozen statistical assumptions. The iid sampling, randomized independence, balanced assignment, consistency, missing-at-random condition, and arrival floor in \cref{obj:ass:iid,obj:ass:randomized-independence,obj:ass:balanced-randomization,obj:ass:consistency,obj:ass:mar,obj:ass:arrival-floor} enter as hypotheses. Verification covers the formal definitions of the law class, identified functional, statistical criteria, and benchmark in \cref{obj:def:law-class,obj:def:identified-functional,obj:def:point-risk,obj:def:interval-class,obj:def:length-risk,obj:def:rate}, together with the estimator, interval, and normalized-prior constructions in \cref{obj:def:mm-estimator,obj:def:mm-interval,obj:def:paired-prior-handle}.

The checked results comprise the identification and upper-risk bounds and honest interval construction in \cref{obj:lem:identification-infrastructure,obj:lem:upper-risk,obj:lem:honest-interval-construction}; the normalized converse in \cref{obj:thm:normalized-converse}; the sampling floors and testing reductions in \cref{obj:lem:parametric-floor-infrastructure,obj:lem:prior-reduction}; and the headline minimax comparisons and complete-arrival reduction in \cref{obj:thm:point-frontier,obj:thm:interval-frontier,obj:prop:complete-arrival-reduction}. These mapped results are machine-checked under the statistical hypotheses displayed in their statements. The cellwise quantities and estimator class identified below are presentation-level definitions; they organize notation and do not supply an unformalized premise to a mapped theorem. Bibliographic comparisons are supported by the cited publications and are outside the machine-checked mathematical layer.

The frozen artifact uses Lean toolchain \texttt{leanprover/lean4:v4.33.0} at repository commit \texttt{d07574d7eb0a903c2a049c6ae28702de46193ec0}. Its sources are under \texttt{CausalSmith/Stat/STAT\_MarNearcompleteFrontier\_Research/}; \texttt{presentation\_crosswalk.json} and \texttt{lean\_snippets.json} in this presentation bundle provide machine-readable object links and excerpts. The cellwise quantities \cref{obj:synth_5} and estimator class \cref{obj:synth_6} are presentation-level definitions; the remaining principal objects map as follows.

\begin{center}
\small
\begin{tabular}{p{0.29\linewidth}p{0.34\linewidth}p{0.25\linewidth}}
\toprule
Paper objects & Lean declarations & Source files \\
\midrule
Sampling and causal assumptions & \texttt{IidSample}, \texttt{RandomizedIndependence}, \texttt{BalancedRandomization}, \texttt{Consistency}, \texttt{MissingAtRandom}, \texttt{ArrivalFloor} & \texttt{Basic.lean} \\
Law, target, criteria, and rate & \texttt{LawClass}, \texttt{psi}, \texttt{pointMinimaxRisk}, \texttt{HonestIntervalClass}, \texttt{lengthMinimaxRisk}, \texttt{rate} & \texttt{Basic.lean} \\
Estimator and interval & \texttt{tauhatMM}, \texttt{intervalMM} & \texttt{Helpers/Estimator.lean} \\
Identification and prior construction & \texttt{tau\_eq\_psi}, \texttt{pairedPriorHandle} & \texttt{Helpers/Identification.lean}; \texttt{Helpers/PairedPrior.lean} \\
Upper bound and interval construction & \texttt{upper\_risk}, \texttt{honest\_interval\_construction} & \texttt{TUpperRisk.lean}; \texttt{THonestInterval.lean} \\
Lower-bound infrastructure & \texttt{parametric\_floor}, \texttt{prior\_reduction}, \texttt{normalized\_converse} & \texttt{TParametricFloor.lean}; \texttt{TPriorReduction.lean}; \texttt{TNormalizedConverse.lean} \\
Headline comparisons & \texttt{point\_frontier}, \texttt{interval\_frontier}, \texttt{complete\_arrival\_reduction} & \texttt{TPointFrontier.lean}; \texttt{TIntervalFrontier.lean}; \texttt{TCompleteArrival.lean} \\
\bottomrule
\end{tabular}
\end{center}


\section{Proofs of the main results}\label{sec:deferred-proofs}

\begin{proof}[Proof of \cref{obj:thm:point-frontier}]
\begin{enumerate}
\item Choose the universal positive squared-risk constant from \cref{obj:lem:prior-reduction}, denoting it by \(c_{\mathrm{point}}\), and the universal upper-bound constant \(C_U>0\) from \cref{obj:lem:upper-risk}. Define
\[
c=\min\{c_{\mathrm{point}},C_U/2\},
\qquad C=C_U.
\]
Both entries in the minimum are positive, and therefore
\[
0<c\le C_U/2<C_U=C<\infty.
\]
These constants are independent of \(n,d,q\).
% lean: CausalSmith.Stat.MarNearcompleteFrontier.point_frontier

\item Fix positive integers \(n,d\) and \(q\in[1/2,1]\). Apply \cref{obj:lem:prior-reduction} with \(\alpha=1/4\in(0,1/2)\). Its combined squared-risk guarantee supplies
\[
c_{\mathrm{point}}\,r_{n,d,q}
\le \mathfrak R^*_{n,d,q}.
\]
By \cref{obj:def:rate},
\[
r_{n,d,q}=\frac1n+g_{n,d,q}^{\,2}\ge0.
\]
Consequently, the choice of \(c\) gives
\[
c\,r_{n,d,q}
\le c_{\mathrm{point}}\,r_{n,d,q}
\le \mathfrak R^*_{n,d,q}.
\]
% lean: CausalSmith.Stat.MarNearcompleteFrontier.prior_reduction

\item We verify that the estimator in \cref{obj:def:mm-estimator} belongs to \(\mathcal T_n\). Write
\[
o=(O_1,\ldots,O_n)\in O^n
\]
for the observed sample tuple; the range calculations below apply to every such tuple. The projection used there satisfies, for every real argument,
\[
-1\le \Pi_{[-1,1]}(t)
=\max\{-1,\min\{1,t\}\}\le1
\qquad(t\in\mathbb R).
\]
Thus the branch \(q=1\) takes values in \([-1,1]\), as does the fallback branch when \(q\ne1\) and either \(\ell_n<128\) or \(d\ge n\ell_n\).

For the remaining branch, define the Poisson weights and the averages over stream assignments by
\[
\omega_{n,N}
=e^{-n/2}\frac{(n/2)^N}{N!}
\qquad(N\in\{0,1,2,\ldots\}),
\]
\[
\overline\tau_N(o)
=4^{-N}
\sum_{\mathbf Z\in\{1,2,3,4\}^{N}}
\widetilde\tau(o;N,\mathbf Z)
\qquad(0\le N\le n).
\]
Here \(\widetilde\tau(o;N,\mathbf Z)\) is the projected randomized estimate specified in \cref{obj:def:mm-estimator} for the fixed sample, auxiliary count, and stream assignments. Each summand lies in \([-1,1]\), and there are \(4^N\) assignments. Hence
\[
-4^N
\le
\sum_{\mathbf Z\in\{1,2,3,4\}^{N}}
\widetilde\tau(o;N,\mathbf Z)
\le4^N,
\qquad
-1\le\overline\tau_N(o)\le1.
\]
For \(N=0\), the assignment set consists of the single empty tuple, so the same calculation applies.

The Poisson weights satisfy
\[
\omega_{n,N}\ge0,
\qquad
0\le\sum_{N=0}^{n}\omega_{n,N}
\le\sum_{N=0}^{\infty}\omega_{n,N}=1.
\]
Since the randomized estimate is zero when \(N>n\), its conditional average is
\[
\widehat\tau_{\mathrm{MM}}(o)
=\sum_{N=0}^{n}\omega_{n,N}\overline\tau_N(o)
\]
in this branch. Multiplying the bounds on each average by its nonnegative weight and summing gives
\[
-1
\le-\sum_{N=0}^{n}\omega_{n,N}
\le\widehat\tau_{\mathrm{MM}}(o)
\le\sum_{N=0}^{n}\omega_{n,N}
\le1.
\]
The projected formulas and finite averages are measurable. Thus, on the entire Cartesian sample space,
\[
\widehat\tau_{\mathrm{MM}}\in\mathcal T_n,
\]
with \(\mathcal T_n\) as defined in \cref{obj:synth_6}.
% lean: CausalSmith.Stat.MarNearcompleteFrontier.tauhatMM_range

\item The zero-pair-mass specialization of the construction in \cref{obj:def:paired-prior-handle} places all normalized mass at its filler label. It has independent fair treatment and complete arrival, and supplies
\[
\mathcal M(d,q)\ne\varnothing
\qquad(d\ge1,\ q\in[1/2,1]),
\]
including \(d=1\). For the real suprema below, we use the conditional real supremum: a nonempty set bounded above has its usual least upper bound, while a set unbounded above is assigned supremum zero. Fix \(T\in\mathcal T_n\). If its set of squared-error risks is bounded above, nonnegativity of squared loss and the supremum property give
\[
0\le
\mathbb E_P\!\left[
\{T(o)-\tau(P)\}^{2}
\right]
\le
\sup_{P'\in\mathcal M(d,q)}
\mathbb E_{P'}\!\left[
\{T(o)-\tau(P')\}^{2}
\right]
\qquad(P\in\mathcal M(d,q)).
\]
The law class is nonempty, so this bounds the supremum below by zero. If the set of squared-error risks is unbounded above, the stated convention gives
\[
\sup_{P\in\mathcal M(d,q)}
\mathbb E_P\!\left[
\{T(o)-\tau(P)\}^{2}
\right]=0.
\]
In either case,
\[
0\le
\sup_{P\in\mathcal M(d,q)}
\mathbb E_P\!\left[
\{T(o)-\tau(P)\}^{2}
\right]
\qquad(T\in\mathcal T_n).
\]
The worst-case risks are therefore bounded below by zero. Evaluating the infimum in \cref{obj:def:point-risk} at the admissible estimator established above yields
\[
\mathfrak R^*_{n,d,q}
\le
\sup_{P\in\mathcal M(d,q)}
\mathbb E_P\!\left[
\{\widehat\tau_{\mathrm{MM}}(o)-\tau(P)\}^{2}
\right].
\]
For every law in this supremum, the sample consists of independent records with its observed-law distribution. Consequently, \cref{obj:lem:upper-risk} gives
\[
\sup_{P\in\mathcal M(d,q)}
\mathbb E_P\!\left[
\{\widehat\tau_{\mathrm{MM}}(o)-\tau(P)\}^{2}
\right]
\le C_U r_{n,d,q}
=C r_{n,d,q}.
\]
Together with the lower bound, this proves the claimed inequalities throughout the stated parameter range.
% lean: CausalSmith.Stat.MarNearcompleteFrontier.point_frontier
\end{enumerate}
\end{proof}

\begin{proof}[Proof of \cref{obj:thm:interval-frontier}]
\begin{enumerate}
\item Fix \(\alpha\in(0,1/2)\). Choose the universal risk constant \(C_U\) supplied by \cref{obj:lem:upper-risk}, so that
\[
0<C_U<\infty,
\qquad
\mathbb E_{P_O^{\otimes n}}
\left[
\bigl(\widehat\tau_{\mathrm{MM}}-\tau(P)\bigr)^2
\right]
\le C_U r_{n,d,q}
\]
for every \(n,d\ge1\), every \(q\in[1/2,1]\), and every \(P\in\mathcal M(d,q)\).
The combined interval bound in \cref{obj:lem:prior-reduction} supplies a constant \(c_{\mathrm{lower},\alpha}\), depending only on \(\alpha\), satisfying
\[
c_{\mathrm{lower},\alpha}>0,
\qquad
c_{\mathrm{lower},\alpha}\sqrt{r_{n,d,q}}
\le \mathfrak L^*_{n,d,q,\alpha}
\]
throughout the same parameter range.
% lean: CausalSmith.Stat.MarNearcompleteFrontier.prior_reduction; CausalSmith.Stat.MarNearcompleteFrontier.upper_risk

\item Define the comparison constants by
\[
c_\alpha:=c_{\mathrm{lower},\alpha},
\qquad
C_\alpha:=c_{\mathrm{lower},\alpha}
       +2\sqrt{C_U/\alpha}+1.
\]
Since \(\sqrt{C_U/\alpha}\ge0\), these constants satisfy
\[
0<c_\alpha<C_\alpha<\infty.
\]
Now fix arbitrary positive integers \(n,d\) and \(q\in[1/2,1]\). By \cref{obj:def:rate},
\[
r_{n,d,q}=\frac1n+g_{n,d,q}^{\,2}\ge0,
\qquad
\sqrt{r_{n,d,q}}\ge0.
\]
The preceding lower bound therefore gives the required inequality
\[
c_\alpha\sqrt{r_{n,d,q}}
\le \mathfrak L^*_{n,d,q,\alpha}.
\]
% lean: CausalSmith.Stat.MarNearcompleteFrontier.interval_frontier

\item Apply \cref{obj:lem:honest-interval-construction} with the chosen \(C_U\). For the procedure in \cref{obj:def:mm-interval}, it yields
\[
\widehat I_{\mathrm{MM},\alpha}\in\mathcal C_\alpha(n,d,q)
\]
and, for every \(P\in\mathcal M(d,q)\),
\[
\mathbb E_{P_O^{\otimes n}}
\left[
\operatorname{length}\{\widehat I_{\mathrm{MM},\alpha}(O^n)\}
\right]
\le
\min\left\{2,2\sqrt{\frac{C_Ur_{n,d,q}}{\alpha}}\right\}.
\]

The law class is nonempty. Indeed, the paired construction with all latent masses zero concentrates on its filler label. Write \(P_{\mathrm{fill}}\) for this law, specified by
\[
\begin{gathered}
X=0,\qquad S(0)=S(1)=Y(0)=0,\qquad R=1,\\
A\sim\operatorname{Bernoulli}(1/2),\qquad
Y(1)\sim\operatorname{Bernoulli}\!\left(\frac{(1+q)^2}{8q}\right),
\qquad A\perp Y(1),\\
S=S(A),\qquad Y=Y(A),\qquad RY=R\,Y.
\end{gathered}
\]
The Bernoulli parameter belongs to \([0,1]\), since
\[
0<(1+q)^2\le4\le8q.
\]
This law has one occupied baseline label, independent balanced treatment, consistent realized measurements, and complete arrival. Hence, by \cref{obj:def:law-class},
\[
P_{\mathrm{fill}}\in\mathcal M(d,q).
\]
Taking the supremum of the expected-length bound gives
\[
\sup_{P\in\mathcal M(d,q)}
\mathbb E_{P_O^{\otimes n}}
\left[
\operatorname{length}\{\widehat I_{\mathrm{MM},\alpha}(O^n)\}
\right]
\le
\min\left\{2,2\sqrt{\frac{C_Ur_{n,d,q}}{\alpha}}\right\}.
\]
% lean: CausalSmith.Stat.MarNearcompleteFrontier.honest_interval_construction; CausalSmith.Stat.MarNearcompleteFrontier.pairedClassLaw

\item Because \(C_U/\alpha\ge0\) and \(r_{n,d,q}\ge0\), square-root multiplicativity gives
\[
\frac{C_Ur_{n,d,q}}{\alpha}
=\frac{C_U}{\alpha}\,r_{n,d,q},
\qquad
\sqrt{\frac{C_Ur_{n,d,q}}{\alpha}}
=\sqrt{\frac{C_U}{\alpha}}\sqrt{r_{n,d,q}}.
\]
Consequently,
\[
\min\left\{2,2\sqrt{\frac{C_Ur_{n,d,q}}{\alpha}}\right\}
\le
2\sqrt{\frac{C_U}{\alpha}}\sqrt{r_{n,d,q}}.
\]
The definition of \(C_\alpha\) and positivity of \(c_{\mathrm{lower},\alpha}\) imply
\[
2\sqrt{\frac{C_U}{\alpha}}\le C_\alpha.
\]
Multiplying by the nonnegative quantity \(\sqrt{r_{n,d,q}}\) yields
\[
2\sqrt{\frac{C_U}{\alpha}}\sqrt{r_{n,d,q}}
\le C_\alpha\sqrt{r_{n,d,q}}.
\]
% lean: CausalSmith.Stat.MarNearcompleteFrontier.interval_frontier

\item For every \(I\in\mathcal C_\alpha(n,d,q)\), the endpoint ordering in \cref{obj:def:interval-class} gives
\[
\operatorname{length}\{I(o)\}\ge0
\quad\text{for every }o\in O^n.
\]
If the law-indexed expected lengths of this candidate interval are bounded above, then
\[
0\le
\mathbb E_{P_{\mathrm{fill},O}^{\otimes n}}
\left[\operatorname{length}\{I(O^n)\}\right]
\le
\sup_{P\in\mathcal M(d,q)}
\mathbb E_{P_O^{\otimes n}}
\left[\operatorname{length}\{I(O^n)\}\right].
\]
For an unbounded-above collection, the real-valued supremum is assigned the value zero. Thus, in the unbounded case,
\[
0=
\sup_{P\in\mathcal M(d,q)}
\mathbb E_{P_O^{\otimes n}}
\left[\operatorname{length}\{I(O^n)\}\right].
\]
In either case, the worst-law expected length is nonnegative. The collection of worst-law expected lengths is therefore bounded below by zero and contains the value attained by \(\widehat I_{\mathrm{MM},\alpha}\). Using the infimum in \cref{obj:def:length-risk}, followed by the preceding bounds, gives
\[
\begin{aligned}
\mathfrak L^*_{n,d,q,\alpha}
&\le
\sup_{P\in\mathcal M(d,q)}
\mathbb E_{P_O^{\otimes n}}
\left[
\operatorname{length}\{\widehat I_{\mathrm{MM},\alpha}(O^n)\}
\right]\\
&\le
\min\left\{2,2\sqrt{\frac{C_Ur_{n,d,q}}{\alpha}}\right\}\\
&\le
2\sqrt{\frac{C_U}{\alpha}}\sqrt{r_{n,d,q}}\\
&\le C_\alpha\sqrt{r_{n,d,q}}.
\end{aligned}
\]
Together with the lower bound, this proves the assertion uniformly over the stated parameter range.
% lean: CausalSmith.Stat.MarNearcompleteFrontier.interval_frontier
\end{enumerate}
\end{proof}

\begin{proof}[Proof of \cref{obj:prop:complete-arrival-reduction}]
\begin{enumerate}
\item Choose universal constants \(0<c<C\) from \cref{obj:thm:point-frontier}. Fix \(\alpha\in(0,1/2)\), and choose constants \(0<c_\alpha<C_\alpha\), depending only on \(\alpha\), from \cref{obj:thm:interval-frontier}. For positive integers \(n,d\), the inclusion \(1\in[1/2,1]\) gives
\[
c\,r_{n,d,1}
\le \mathfrak R^*_{n,d,1}
\le C\,r_{n,d,1},
\qquad
c_\alpha\sqrt{r_{n,d,1}}
\le \mathfrak L^*_{n,d,1,\alpha}
\le C_\alpha\sqrt{r_{n,d,1}}.
\]
We retain these inequalities while establishing the estimator and risk identities at complete arrival.
% lean: CausalSmith.Stat.MarNearcompleteFrontier.complete_arrival_reduction

\item The law class is nonempty. Indeed, the paired-law construction with zero perturbation and every latent draw at the residual zero-mass atom reduces, at \(q=1\), to the law
\[
P_{\mathrm{fill}}
:=
\text{the full-data law specified by }
\left\{
\begin{aligned}
&A,Y(1)\text{ independent, each with law }\operatorname{Bernoulli}(1/2),\\
&X=0,\qquad S(0)=S(1)=S=Y(0)=0,\\
&R=1,\qquad Y=A\,Y(1).
\end{aligned}
\right.
\]
This law has one occupied baseline label, independent balanced treatment assignment, and consistent realized measurements. Its constant arrival indicator satisfies conditional independence and has arrival probability one in every occupied cell. Thus
\[
P_{\mathrm{fill}}\in\mathcal M(d,1)
\qquad(d\ge1).
\]
This construction also applies when \(d=1\), when the collection of baseline pairs is empty.
% lean: CausalSmith.Stat.MarNearcompleteFrontier.pairedClassLaw

\item Fix \(P\in\mathcal M(d,1)\). On the entire Cartesian observed-record space, define the complete-arrival score by
\[
\zeta_{\mathrm{HT}}(O)
:=2(2A-1)(RY),
\qquad
O=(X,A,S,R,RY)\in\{0,\ldots,d-1\}\times\{0,1\}^{4}.
\]
Under the observed-data law, the fifth coordinate is the product of arrival and realized outcome. Splitting the finite expectation according to the treatment and outcome-mark coordinates therefore gives
\[
\mathbb E_P[\zeta_{\mathrm{HT}}(O)]
=
2\bigl\{
P(A=1,R=1,Y=1)-P(A=0,R=1,Y=1)
\bigr\}.
\]

To remove the arrival restriction, define, for every
\(x\in\{0,\ldots,d-1\}\) and \(a,s\in\{0,1\}\),
\[
\begin{aligned}
\mathsf p_{\mathrm{cell}}(x,a,s)
&:=P(X=x,A=a,S=s),\\
\mathsf p_{\mathrm{arr}}(x,a,s)
&:=P(X=x,A=a,S=s,R=1),\\
\mathsf p_{\mathrm{mis}}(x,a,s)
&:=P(X=x,A=a,S=s,R=0).
\end{aligned}
\]
These masses satisfy
\[
\mathsf p_{\mathrm{cell}}(x,a,s)
=
\mathsf p_{\mathrm{arr}}(x,a,s)
+\mathsf p_{\mathrm{mis}}(x,a,s),
\qquad
\mathsf p_{\mathrm{arr}}(x,a,s),
\mathsf p_{\mathrm{mis}}(x,a,s)\ge0.
\]
If the cell mass is zero, this identity forces its missing mass to be zero. If the cell mass is positive, \cref{obj:ass:arrival-floor} at \(q=1\) gives
\[
1\le
\frac{\mathsf p_{\mathrm{arr}}(x,a,s)}
     {\mathsf p_{\mathrm{cell}}(x,a,s)}
\quad\Longrightarrow\quad
\mathsf p_{\mathrm{cell}}(x,a,s)
\le\mathsf p_{\mathrm{arr}}(x,a,s)
\quad\Longrightarrow\quad
\mathsf p_{\mathrm{mis}}(x,a,s)=0.
\]
Consequently, every missing-cell mass is zero, including null cells. Every nonnegative full-data atom mass contributing to such a cell is also zero. Separating the cases \(R=1\) and \(R=0\) in the finite atom sums yields
\[
P(A=a,R=1,Y=1)=P(A=a,Y=1),
\qquad a\in\{0,1\}.
\]

By \cref{obj:ass:balanced-randomization},
\[
P(A=0)=1-P(A=1)=\frac12.
\]
Using \cref{obj:ass:consistency,obj:ass:randomized-independence}, we obtain, for each \(a\in\{0,1\}\),
\[
\begin{aligned}
2P(A=a,Y=1)
&=2P(A=a,Y(a)=1)\\
&=2P(A=a)P(Y(a)=1)
=P(Y(a)=1).
\end{aligned}
\]
Substitution into the score expectation, followed by the binary-outcome identity, gives
\[
\begin{aligned}
\mathbb E_P[\zeta_{\mathrm{HT}}(O)]
&=2\{P(A=1,Y=1)-P(A=0,Y=1)\}\\
&=P(Y(1)=1)-P(Y(0)=1)\\
&=\mathbb E_P[Y(1)-Y(0)]
=\tau(P).
\end{aligned}
\]
% lean: CausalSmith.Stat.MarNearcompleteFrontier.complete_arrival_score_mean

\item For every Cartesian sample \(O^n=(O_1,\ldots,O_n)\), define
\[
\overline\zeta_{n,\mathrm{HT}}(O^n)
:=
\frac1n\sum_{i=1}^{n}\zeta_{\mathrm{HT}}(O_i)
=
\frac2n\sum_{i=1}^{n}(2A_i-1)Y_i^{\mathrm{obs}},
\qquad
Y_i^{\mathrm{obs}}:=(RY)_i.
\]
The binary coordinates imply
\[
\zeta_{\mathrm{HT}}(O)\in\{-2,0,2\}\subseteq[-2,2].
\]
The bounded-variable variance inequality therefore gives
\[
\operatorname{Var}_P\{\zeta_{\mathrm{HT}}(O)\}\le4.
\]
All these variables are square integrable because their domain is finite. Under \cref{obj:ass:iid}, the iid average identities and the score mean established above give
\[
\mathbb E_P[\overline\zeta_{n,\mathrm{HT}}(O^n)]=\tau(P),
\qquad
\operatorname{Var}_P\{\overline\zeta_{n,\mathrm{HT}}(O^n)\}
=
\frac1n\operatorname{Var}_P\{\zeta_{\mathrm{HT}}(O)\}.
\]
Thus
\[
\begin{aligned}
\mathbb E_P\!\left[
\{\overline\zeta_{n,\mathrm{HT}}(O^n)-\tau(P)\}^{2}
\right]
&=\operatorname{Var}_P\{\overline\zeta_{n,\mathrm{HT}}(O^n)\}\\
&=\frac1n\operatorname{Var}_P\{\zeta_{\mathrm{HT}}(O)\}
\le\frac4n.
\end{aligned}
\]
Taking the supremum over the nonempty law class gives
\[
\sup_{P\in\mathcal M(d,1)}
\mathbb E_P\!\left[
\{\overline\zeta_{n,\mathrm{HT}}(O^n)-\tau(P)\}^{2}
\right]
\le\frac4n.
\]
% lean: CausalSmith.Stat.MarNearcompleteFrontier.complete_arrival_unprojected_risk_of_mean

\item The \(q=1\) branch of \cref{obj:def:mm-estimator} gives, for every Cartesian sample,
\[
\widehat\tau_{\mathrm{MM}}(O^n)
=
\Pi_{[-1,1]}\!\left[\overline\zeta_{n,\mathrm{HT}}(O^n)\right].
\]
Binary potential outcomes have difference in \([-1,1]\). Averaging these pointwise bounds against the full-data probability law gives
\[
-1=\mathbb E_P[-1]
\le \mathbb E_P[Y(1)-Y(0)]
=\tau(P)
\le\mathbb E_P[1]=1.
\]

The projection inequality follows by its three defining cases. If
\(\overline\zeta_{n,\mathrm{HT}}(O^n)\le-1\), then
\[
\left|\widehat\tau_{\mathrm{MM}}(O^n)-\tau(P)\right|
=\tau(P)+1
\le\tau(P)-\overline\zeta_{n,\mathrm{HT}}(O^n)
=
\left|\overline\zeta_{n,\mathrm{HT}}(O^n)-\tau(P)\right|.
\]
If \(-1\le\overline\zeta_{n,\mathrm{HT}}(O^n)\le1\), then
\[
\left|\widehat\tau_{\mathrm{MM}}(O^n)-\tau(P)\right|
=
\left|\overline\zeta_{n,\mathrm{HT}}(O^n)-\tau(P)\right|.
\]
If \(1\le\overline\zeta_{n,\mathrm{HT}}(O^n)\), then
\[
\left|\widehat\tau_{\mathrm{MM}}(O^n)-\tau(P)\right|
=1-\tau(P)
\le\overline\zeta_{n,\mathrm{HT}}(O^n)-\tau(P)
=
\left|\overline\zeta_{n,\mathrm{HT}}(O^n)-\tau(P)\right|.
\]
This proves the asserted pointwise comparison for every sample.
% lean: CausalSmith.Stat.MarNearcompleteFrontier.complete_arrival_projection_error

\item Squaring the pointwise comparison and integrating gives, for each \(P\in\mathcal M(d,1)\),
\[
\mathbb E_P\!\left[
\{\widehat\tau_{\mathrm{MM}}(O^n)-\tau(P)\}^{2}
\right]
\le
\mathbb E_P\!\left[
\{\overline\zeta_{n,\mathrm{HT}}(O^n)-\tau(P)\}^{2}
\right].
\]
The unprojected risks form a nonempty family bounded above by \(4/n\). Each projected risk is therefore bounded by their supremum, and taking the supremum of the projected risks yields
\[
\sup_{P\in\mathcal M(d,1)}
\mathbb E_P\!\left[
\{\widehat\tau_{\mathrm{MM}}(O^n)-\tau(P)\}^{2}
\right]
\le
\sup_{P\in\mathcal M(d,1)}
\mathbb E_P\!\left[
\{\overline\zeta_{n,\mathrm{HT}}(O^n)-\tau(P)\}^{2}
\right]
\le\frac4n.
\]
% lean: CausalSmith.Stat.MarNearcompleteFrontier.complete_arrival_projection_risk

\item Finally, the error-scale definition and \cref{obj:def:rate} give
\[
g_{n,d,1}
=(1-1)\min\left\{1,\frac{d}{n\log(e+n)}\right\}=0,
\qquad
r_{n,d,1}=\frac1n+g_{n,d,1}^{2}=\frac1n.
\]
Since \(n\ge1\),
\[
\sqrt{r_{n,d,1}}
=\sqrt{\frac1n}
=\frac1{\sqrt n}.
\]
Substituting these identities into the frontier inequalities selected at the start proves
\[
\frac cn\le\mathfrak R^*_{n,d,1}\le\frac Cn,
\qquad
\frac{c_\alpha}{\sqrt n}
\le\mathfrak L^*_{n,d,1,\alpha}
\le\frac{C_\alpha}{\sqrt n}.
\]
The constants retain their stated dependence uniformly over all positive integers \(n,d\).
% lean: CausalSmith.Stat.MarNearcompleteFrontier.complete_arrival_reduction
\end{enumerate}
\end{proof}

\begin{proof}[Proof of \cref{obj:thm:normalized-converse}]
Fix \(\beta\in(0,1/4)\). We first construct the required priors for all sufficiently large sample sizes, with a threshold depending only on \(\beta\), and then extend the construction to every positive sample size.

\begin{enumerate}
\item
Use the construction and notation of
\cref{obj:def:paired-prior-handle}. Throughout the paired construction, expectations and probabilities refer to its independent latent draws. The elementary scale bounds are
\[
L\ge1,\qquad 2\le K,\qquad 8L\le K\le9L,\qquad M\le nL.
\]
Indeed, the upper bound for \(K=\lceil8L\rceil\) follows from
\(K<8L+1\le9L\). Consequently,
\[
0\le 2MhB_0=\frac{M}{512nK}\le\frac12,
\qquad
f=1-2MhB_0\ge\frac12,
\qquad
J=f+2\sum_{j=1}^{M}p_j\ge\frac12.
\]
Thus normalization is well defined, and
\[
\frac fJ+2\sum_{j=1}^{M}\frac{p_j}{J}=1,
\qquad
2M+1\le d.
\]
These conclusions also cover \(M=0\), when the law consists of the filler label.

For every \(\sigma,z_j,u\in\{-1,1\}\), the probabilities specified in the construction satisfy
\[
q\le \rho_{\sigma,j,u}\le1,
\qquad
0\le \frac b2+c_0u\le q\le\rho_{\sigma,j,u},
\qquad
0\le\mu_{\sigma,j,u}\le1.
\]
Here the middle bound follows by substituting
\(b=(1+q)/2\), \(c_0=1/16\), and \(q\ge1/2\).
The filler mean is valid as well:
\[
0\le b^2=\frac{(1+q)^2}{4}
\le\frac{1+3q}{4}\le2q,
\qquad
0\le\mu_0=\frac{b^2}{2q}\le1.
\]
Treatment is an independent fair coin, and the realized measurements are selected from their potential measurements as specified in the construction. Conditional on an occupied baseline label and treatment arm, arrival is independent of the outcome: it is complete at the filler label and in the control arm, and has probability \(\rho_{\sigma,j,u}\) in a treated pair cell. Unoccupied cells have zero mass. Hence every constructed normalized law belongs to
\(\mathcal M(d,q)\) in \cref{obj:def:law-class}. The finitely many latent configurations therefore induce discrete probability priors supported on this class.
% lean: CausalSmith.Stat.MarNearcompleteFrontier.pairedFullLaw_legal

\item
We record the interpolation identities needed for the count comparison. Polynomial interpolation at zero gives
\[
\varphi_{\mathrm{pol}}(0)
=\sum_{i=0}^{m}\varphi_{\mathrm{pol}}(x_i)l_i(0)
\qquad
\text{for every polynomial \(\varphi_{\mathrm{pol}}\) of degree at most \(m=K-1\)}.
\]
Applying this identity to the constant polynomial and to positive powers yields
\[
\sum_{i=0}^{m}l_i(0)=1,
\qquad
\sum_{i=0}^{m}\nu_i=\frac1D,
\qquad
\sum_{i=0}^{m}\nu_i x_i^{t_{\mathrm{pol}}}=0
\quad(1\le t_{\mathrm{pol}}\le K-1).
\]
Also,
\[
D=\sum_{i=0}^{m}|l_i(0)|\ge1,
\qquad
\sum_{i=0}^{m}|\nu_i|=1.
\]
Since \(x_i\ge h>0\), the latent probabilities obey
\[
\sum_{i=0}^{m}\frac{h|\nu_i|}{x_i}\le1.
\]
Their residual mass is therefore nonnegative. Substitution into the discrete latent distribution gives
\[
\begin{aligned}
\mathbb E p&=hB_0,
&
\mathbb E(pz)&=\frac{hB_0}{D},\\
\mathbb E(p^{t_{\mathrm{pol}}}z)
&=hB_0^{t_{\mathrm{pol}}}
  \sum_{i=0}^{m}\nu_i x_i^{t_{\mathrm{pol}}-1}
 =0
&& (2\le t_{\mathrm{pol}}\le K),\\
\mathbb E p^2
&=\mathbb E[(pz)^2]
 =hB_0^2\sum_{i=0}^{m}|\nu_i|x_i
 \le hB_0^2
 =\frac1{1024^2n^2}.
\end{aligned}
\]
The atom with \(p=0\) contributes zero to each displayed positive-power moment.
% lean: CausalSmith.Stat.MarNearcompleteFrontier.latent_signed_power_moment_zero

\item
Let the eight-coordinate pair vector be
\[
V_\sigma(z,r)
=\bigl(v_\sigma(z,r),v_\sigma(z,-r)\bigr),
\qquad
\sigma,z,r\in\{-1,1\}.
\]
Its coordinates are nonnegative and satisfy
\[
\sum_{\ell_{\mathrm{atom}}=1}^{8}
V_{\sigma,\ell_{\mathrm{atom}}}(z,r)=2.
\]
For a count vector
\(\boldsymbol{k}_{\mathrm{pair}}
=(k_{\mathrm{pair},1},\ldots,k_{\mathrm{pair},8})\in\mathbb N^8\),
define
\[
t_{\mathrm{count}}
=\sum_{\ell_{\mathrm{atom}}=1}^{8}k_{\mathrm{pair},\ell_{\mathrm{atom}}},
\qquad
\mathcal H_\sigma(z;\boldsymbol{k}_{\mathrm{pair}})
=\frac12\sum_{r\in\{-1,1\}}
 \prod_{\ell_{\mathrm{atom}}=1}^{8}
 V_{\sigma,\ell_{\mathrm{atom}}}(z,r)^{
 k_{\mathrm{pair},\ell_{\mathrm{atom}}}}
\]
Products use the convention \(0^0=1\). The likelihood of this count vector in one mixed pair block is
\[
\mathcal L_{\sigma,\mathrm{pair}}
(\boldsymbol{k}_{\mathrm{pair}})
=
\mathbb E\!\left[
e^{-8np}\,
\frac{(4np)^{t_{\mathrm{count}}}}
     {\prod_{\ell_{\mathrm{atom}}=1}^{8}
       (k_{\mathrm{pair},\ell_{\mathrm{atom}}})!}
\,
\mathcal H_\sigma(z;\boldsymbol{k}_{\mathrm{pair}})
\right],
\]
where the expectation is over the latent pair \((p,z)\).

Changing \(z\) to \(-z\) exchanges the two signs, so
\[
\mathcal H_-(z;\boldsymbol{k}_{\mathrm{pair}})
-\mathcal H_+(z;\boldsymbol{k}_{\mathrm{pair}})
=
z\bigl[
\mathcal H_-(1;\boldsymbol{k}_{\mathrm{pair}})
-\mathcal H_+(1;\boldsymbol{k}_{\mathrm{pair}})
\bigr].
\]
This difference is zero when \(t_{\mathrm{count}}=0\). It is also zero when
\(t_{\mathrm{count}}=1\), because averaging over the fair orientation gives, at either label, the four coefficients
\[
\left(\frac12,\frac\delta4,\frac b4,\frac b4\right),
\]
independently of \(\sigma z\).

For \(v_{\mathrm{exp}}\in\mathbb N\), define the signed Taylor coefficient
\[
\Gamma_{v_{\mathrm{exp}}}
(\boldsymbol{k}_{\mathrm{pair}})
=
\mathbb E\!\left[
\frac{(-8np)^{v_{\mathrm{exp}}}}{v_{\mathrm{exp}}!}
\frac{(4np)^{t_{\mathrm{count}}}}
     {\prod_{\ell_{\mathrm{atom}}=1}^{8}
       (k_{\mathrm{pair},\ell_{\mathrm{atom}}})!}
\bigl(\mathcal H_-(z;\boldsymbol{k}_{\mathrm{pair}})
-\mathcal H_+(z;\boldsymbol{k}_{\mathrm{pair}})\bigr)
\right].
\]
The exponential series, together with \(0\le p\le B_0\), gives
\[
\mathcal L_{-,\mathrm{pair}}(\boldsymbol{k}_{\mathrm{pair}})
-\mathcal L_{+,\mathrm{pair}}(\boldsymbol{k}_{\mathrm{pair}})
=
\sum_{v_{\mathrm{exp}}=0}^{\infty}
\Gamma_{v_{\mathrm{exp}}}(\boldsymbol{k}_{\mathrm{pair}}).
\]
The low-count cancellations above and the signed moment identities imply
\[
\Gamma_{v_{\mathrm{exp}}}(\boldsymbol{k}_{\mathrm{pair}})=0
\qquad
\text{if }v_{\mathrm{exp}}+t_{\mathrm{count}}\le K.
\]
Indeed, when \(t_{\mathrm{count}}\ge2\), this coefficient is a fixed multiple of
\(\mathbb E[p^{v_{\mathrm{exp}}+t_{\mathrm{count}}}z]\).

For fixed count degree, the multinomial identity gives
\[
\sum_{\boldsymbol{k}_{\mathrm{pair}}:\,
      t_{\mathrm{count}}=t_{\mathrm{degree}}}
\frac{\mathcal H_\sigma(z;\boldsymbol{k}_{\mathrm{pair}})}
     {\prod_{\ell_{\mathrm{atom}}=1}^{8}
       (k_{\mathrm{pair},\ell_{\mathrm{atom}}})!}
=\frac{2^{t_{\mathrm{degree}}}}{t_{\mathrm{degree}}!},
\qquad t_{\mathrm{degree}}\in\mathbb N.
\]
Bounding the absolute mark difference by the sum of its two nonnegative terms therefore yields
\[
\sum_{\boldsymbol{k}_{\mathrm{pair}}:\,
      t_{\mathrm{count}}=t_{\mathrm{degree}}}
|\Gamma_{v_{\mathrm{exp}}}(\boldsymbol{k}_{\mathrm{pair}})|
\le
2\frac{(8nB_0)^{v_{\mathrm{exp}}}}{v_{\mathrm{exp}}!}
 \frac{(8nB_0)^{t_{\mathrm{degree}}}}{t_{\mathrm{degree}}!}.
\]
Thus, for each total degree \(s_{\mathrm{tot}}\in\mathbb N\),
\[
\sum_{t_{\mathrm{degree}}=0}^{s_{\mathrm{tot}}}
\ \sum_{\boldsymbol{k}_{\mathrm{pair}}:\,
      t_{\mathrm{count}}=t_{\mathrm{degree}}}
|\Gamma_{s_{\mathrm{tot}}-t_{\mathrm{degree}}}(\boldsymbol{k}_{\mathrm{pair}})|
\le
\begin{cases}
0,&s_{\mathrm{tot}}\le K,\\[2mm]
\displaystyle
2\frac{(16nB_0)^{s_{\mathrm{tot}}}}{s_{\mathrm{tot}}!},
&s_{\mathrm{tot}}>K.
\end{cases}
\]
The second line follows by factorial-series convolution. Its summable majorant justifies the absolute summations and gives
\[
\sum_{\boldsymbol{k}_{\mathrm{pair}}\in\mathbb N^8}
\left|
\mathcal L_{-,\mathrm{pair}}(\boldsymbol{k}_{\mathrm{pair}})
-\mathcal L_{+,\mathrm{pair}}(\boldsymbol{k}_{\mathrm{pair}})
\right|
\le
2\sum_{s_{\mathrm{tot}}>K}
\frac{(16nB_0)^{s_{\mathrm{tot}}}}{s_{\mathrm{tot}}!}.
\]

Denote the two complete mixed count laws by
\[
\mathsf C_\sigma
=\operatorname{Law}(\text{the unnormalized observed-atom counts
under sign }\sigma),
\qquad \sigma\in\{-1,1\}.
\]
The independent pair blocks have probability likelihoods, and the filler likelihood is common to both signs. Write the tuple of pair count vectors as
\[
\boldsymbol{k}_{\mathrm{blocks}}
=\bigl(\boldsymbol{k}_{\mathrm{pair}}^{(1)},\ldots,
       \boldsymbol{k}_{\mathrm{pair}}^{(M)}\bigr)
\in(\mathbb N^8)^M.
\]
Each block likelihood is nonnegative and normalized:
\[
\sum_{\boldsymbol{k}_{\mathrm{pair}}\in\mathbb N^8}
\mathcal L_{\sigma,\mathrm{pair}}(\boldsymbol{k}_{\mathrm{pair}})=1,
\qquad \sigma\in\{-1,1\}.
\]
The product difference has the telescoping identity
\[
\begin{aligned}
&\prod_{j=1}^{M}
\mathcal L_{-,\mathrm{pair}}(\boldsymbol{k}_{\mathrm{pair}}^{(j)})
-\prod_{j=1}^{M}
\mathcal L_{+,\mathrm{pair}}(\boldsymbol{k}_{\mathrm{pair}}^{(j)})\\
&=\sum_{j=1}^{M}
\Bigl[
\mathcal L_{-,\mathrm{pair}}(\boldsymbol{k}_{\mathrm{pair}}^{(j)})
-\mathcal L_{+,\mathrm{pair}}(\boldsymbol{k}_{\mathrm{pair}}^{(j)})
\Bigr]
\prod_{\ell_{\mathrm{block}}=1}^{j-1}
\mathcal L_{-,\mathrm{pair}}
 (\boldsymbol{k}_{\mathrm{pair}}^{(\ell_{\mathrm{block}})})
\prod_{\ell_{\mathrm{block}}=j+1}^{M}
\mathcal L_{+,\mathrm{pair}}
 (\boldsymbol{k}_{\mathrm{pair}}^{(\ell_{\mathrm{block}})}).
\end{aligned}
\]
Empty products equal one. Taking absolute values and summing over all block count vectors gives
\[
\begin{aligned}
&\sum_{\boldsymbol{k}_{\mathrm{blocks}}\in(\mathbb N^8)^M}
\left|
\prod_{j=1}^{M}
\mathcal L_{-,\mathrm{pair}}(\boldsymbol{k}_{\mathrm{pair}}^{(j)})
-\prod_{j=1}^{M}
\mathcal L_{+,\mathrm{pair}}(\boldsymbol{k}_{\mathrm{pair}}^{(j)})
\right|\\
&\le\sum_{j=1}^{M}
\sum_{\boldsymbol{k}_{\mathrm{pair}}\in\mathbb N^8}
\left|
\mathcal L_{-,\mathrm{pair}}(\boldsymbol{k}_{\mathrm{pair}})
-\mathcal L_{+,\mathrm{pair}}(\boldsymbol{k}_{\mathrm{pair}})
\right|\\
&\le 2M\sum_{s_{\mathrm{tot}}>K}
\frac{(16nB_0)^{s_{\mathrm{tot}}}}{s_{\mathrm{tot}}!}.
\end{aligned}
\]
Indeed, in each telescoping summand, summing every block other than the selected block gives one by normalization. The sum over the \(M\) selected blocks therefore gives the factor \(M\). The common filler likelihood also sums to one. Splitting a full count vector into these blocks consequently gives
\[
\operatorname{TV}(\mathsf C_-,\mathsf C_+)
\le
\frac12\sum_{\text{full count vectors}}
|\mathsf C_-(\{\text{vector}\})-\mathsf C_+(\{\text{vector}\})|
\le
M\sum_{s_{\mathrm{tot}}>K}
\frac{(16nB_0)^{s_{\mathrm{tot}}}}{s_{\mathrm{tot}}!}.
\]
Counts outside the specified atoms are zero under both laws. For \(M=0\), both count laws equal the common filler law and the bound is zero.
% lean: CausalSmith.Stat.MarNearcompleteFrontier.mixedCountLaw_tv_le_factorial_tail

\item
For \(\xi\ge0\), define
\[
\mathcal E_K(\xi)
=\sum_{s_{\mathrm{tot}}>K}
\frac{\xi^{s_{\mathrm{tot}}}}{s_{\mathrm{tot}}!}.
\]
Termwise exponential weighting gives
\[
\mathcal E_K(\xi)
\le e^{-K}\sum_{s_{\mathrm{tot}}=0}^{\infty}
\frac{(e\xi)^{s_{\mathrm{tot}}}}{s_{\mathrm{tot}}!}
=\exp(-K+e\xi),
\]
because \(e^{s_{\mathrm{tot}}-K}\ge1\) on the tail. Since
\(16nB_0=K/64\), \(e\le3\), and \(K\ge8L\),
\[
\mathcal E_K(16nB_0)
\le\exp(-K+eK/64)\le e^{-4L}.
\]
Using \(M\le nL\), \(n\le e^L\), and \(L\le e^L\), we obtain
\[
\operatorname{TV}(\mathsf C_-,\mathsf C_+)
\le Me^{-4L}
\le nLe^{-4L}
\le e^{-2L}
\le e^{-L}
=\frac1{e+n}
\le\frac1n.
\]
Choose an integer threshold satisfying
\[
N_{\mathrm{count}}\ge2,
\qquad
N_{\mathrm{count}}>\frac2\beta.
\]
Then, uniformly in \(d,q\),
\[
n\ge N_{\mathrm{count}}
\quad\Longrightarrow\quad
\operatorname{TV}(\mathsf C_-,\mathsf C_+)\le\frac\beta2.
\]
% lean: CausalSmith.Stat.MarNearcompleteFrontier.mixedCountLaw_tv_eventual

\item
We next transfer the count comparison to the complete fixed-size sample. Conditional on the latent variables, define the total count and its mean by
\[
N_{\mathrm{Pois}}
=\sum_{o\in O}\text{count at }o,
\qquad
\Lambda_{\mathrm{count}}=4nJ.
\]
The exact observed-atom probabilities identify the count experiment as independent Poisson counts from \(P_{\sigma,O}\) at total intensity
\(\Lambda_{\mathrm{count}}\). For example,
\[
4np_jv_\sigma(z_j,u)
=(4nJ)\frac{p_j}{J}v_\sigma(z_j,u),
\]
and the same identity holds for the filler coefficients. Hence
\[
N_{\mathrm{Pois}}\mid\text{latent variables}
\sim\operatorname{Poisson}(\Lambda_{\mathrm{count}}),
\qquad
\Lambda_{\mathrm{count}}\ge2n.
\]

For this conditional Poisson law, put
\[
x_{\mathrm{tail}}
=\sqrt{\frac{\Lambda_{\mathrm{count}}n}{4}},
\qquad
\vartheta_{\mathrm{tail}}
=\frac{x_{\mathrm{tail}}}{\Lambda_{\mathrm{count}}}.
\]
Both are positive, and
\[
x_{\mathrm{tail}}\le\frac{\Lambda_{\mathrm{count}}}{2}
\le\Lambda_{\mathrm{count}}-n.
\]
The centered Poisson exponential moment is
\[
\mathbb E\!\left[
e^{-\vartheta_{\mathrm{tail}}
(N_{\mathrm{Pois}}-\Lambda_{\mathrm{count}})}
\,\middle|\,\text{latent variables}\right]
=
\exp\!\left\{
\Lambda_{\mathrm{count}}
(e^{-\vartheta_{\mathrm{tail}}}-1+\vartheta_{\mathrm{tail}})
\right\}.
\]
For every nonnegative real \(\vartheta_{\mathrm{tail}}\),
\[
e^{-\vartheta_{\mathrm{tail}}}-1+\vartheta_{\mathrm{tail}}
\le\frac{\vartheta_{\mathrm{tail}}^2}{2}.
\]
One way to verify this bound is to combine
\(e^{\vartheta_{\mathrm{tail}}}\ge
1+\vartheta_{\mathrm{tail}}+\vartheta_{\mathrm{tail}}^2/2\)
with
\[
\left(1+\vartheta_{\mathrm{tail}}
+\frac{\vartheta_{\mathrm{tail}}^2}{2}\right)
\left(1-\vartheta_{\mathrm{tail}}
+\frac{\vartheta_{\mathrm{tail}}^2}{2}\right)
=1+\frac{\vartheta_{\mathrm{tail}}^4}{4}\ge1.
\]
The second factor is positive, so taking reciprocals gives the claimed bound. Exponential Markov inequality now yields
\[
\begin{aligned}
\Pr(N_{\mathrm{Pois}}<n\mid\text{latent variables})
&\le
\Pr(\Lambda_{\mathrm{count}}-N_{\mathrm{Pois}}>
x_{\mathrm{tail}}\mid\text{latent variables})\\
&\le
\exp\!\left\{
-\frac{x_{\mathrm{tail}}^2}{\Lambda_{\mathrm{count}}}
+\Lambda_{\mathrm{count}}
(e^{-\vartheta_{\mathrm{tail}}}-1+\vartheta_{\mathrm{tail}})
\right\}\\
&\le
\exp\!\left(-\frac{x_{\mathrm{tail}}^2}
                   {2\Lambda_{\mathrm{count}}}\right)
=e^{-n/8}
\le\frac8n.
\end{aligned}
\]

Uniformly ordering independent Poisson counts produces iid marks with law
\(P_{\sigma,O}\), independently of their total count. Consequently,
conditional on the latent variables, the output of
\(\kappa_n\) on \(N_{\mathrm{Pois}}\ge n\) has law
\(P_{\sigma,O}^{\otimes n}\). Coupling this output with an iid sample and using the fixed fallback on the complementary event gives
\[
\operatorname{TV}\bigl(Q_\sigma^{(n)},\kappa_n\mathsf C_\sigma\bigr)
\le
\mathbb E\bigl[
\Pr(N_{\mathrm{Pois}}<n\mid\text{latent variables})\bigr]
\le\frac8n.
\]
Here \(\kappa_n\mathsf C_\sigma\) denotes the output law of the kernel.
Its independence of the sign and latent variables permits total-variation contraction. The triangle inequality therefore gives
\[
\operatorname{TV}(Q_-^{(n)},Q_+^{(n)})
\le
\operatorname{TV}(\mathsf C_-,\mathsf C_+)+\frac{16}{n}.
\]
Choose integer thresholds
\[
N_{\mathrm{transfer}}\ge2,
\qquad
N_{\mathrm{transfer}}>\frac{32}{\beta},
\qquad
N_{\mathrm{TV}}
=\max\{N_{\mathrm{count}},N_{\mathrm{transfer}}\}.
\]
For every \(n\ge N_{\mathrm{TV}}\), the two terms on the right are each at most \(\beta/2\), so
\[
\operatorname{TV}(Q_-^{(n)},Q_+^{(n)})\le\beta.
\]
% lean: CausalSmith.Stat.MarNearcompleteFrontier.paired_predictive_tv_eventual

\item
We now establish target separation. First, the interpolation norm has the uniform bound
\[
1\le D\le\cosh2.
\]
To see the upper bound, define the transformed nodes, exterior point, and hyperbolic parameter by
\[
y_i=\frac{2x_i-1-h}{1-h}
=\cos\frac{i\pi}{m},
\qquad
y_*=-\frac{1+h}{1-h},
\qquad
t_{\mathrm{hyp}}=\log\frac{K+1}{K-1}.
\]
Use the Chebyshev convention
\[
T_0(t_{\mathrm{Cheb}})=1,\qquad
T_1(t_{\mathrm{Cheb}})=t_{\mathrm{Cheb}},\qquad
T_{\ell_{\mathrm{pol}}+1}(t_{\mathrm{Cheb}})
=2t_{\mathrm{Cheb}}T_{\ell_{\mathrm{pol}}}(t_{\mathrm{Cheb}})
-T_{\ell_{\mathrm{pol}}-1}(t_{\mathrm{Cheb}})
\quad(\ell_{\mathrm{pol}}\ge1,\ t_{\mathrm{Cheb}}\in\mathbb R).
\]
The ordered nodes give
\(\operatorname{sign}(l_i(0))=(-1)^{m+i}\), while
\(T_m(y_i)=(-1)^i\). Interpolating \(T_m\) at \(y_*\) therefore gives
\[
D=|T_m(y_*)|.
\]
Since \(h=K^{-2}\),
\[
y_*=-\cosh(t_{\mathrm{hyp}}),
\qquad
0\le mt_{\mathrm{hyp}}
=(K-1)\log\frac{K+1}{K-1}
\le
(K-1)\left(\frac{K+1}{K-1}-1\right)=2.
\]
The inequality uses \(\log t\le t-1\) for \(t>0\).
The Chebyshev hyperbolic identity then gives
\[
D=\cosh(mt_{\mathrm{hyp}})\le\cosh2.
\]

Define the signed latent mass, its mean, and the fixed hyperbolic constant by
\[
W=\sum_{j=1}^{M}p_jz_j,
\qquad
\mu_W=\frac{MhB_0}{D}
=\frac{M}{1024nKD},
\qquad
C_{\mathrm{hyp}}=\cosh2.
\]
Independence of the latent pairs and the moment identities above imply
\[
\mathbb EW=\mu_W,
\qquad
\operatorname{Var}(W)
=M\left(\mathbb E[(pz)^2]-\left(\frac{hB_0}{D}\right)^2\right)
\le\frac{M}{1024^2n^2}.
\]
Similarly, the deterministic filler mass gives
\[
J=1+2\sum_{j=1}^{M}(p_j-hB_0),
\qquad
\mathbb EJ=1,
\]
and
\[
\operatorname{Var}(J)
=4M\bigl(\mathbb Ep^2-(hB_0)^2\bigr)
\le\frac{4M}{1024^2n^2}.
\]
In both variance computations, the cross terms vanish because independent centered summands have zero product expectation.

Finally, opposite orientations in a pair satisfy the exact identity
\[
\mu_{\sigma,j,1}+\mu_{\sigma,j,-1}
=2\mu_0+\frac{\sigma\delta z_j}{16q}.
\]
This follows by putting the two fractions over a common denominator and using
\[
b^2-\frac{\delta^2}{4}=q.
\]
Since \(Y(0)=0\), summing the filler and pair contributions and dividing by \(J\) yields
\[
\begin{aligned}
\tau(P_\sigma)
&=
\frac{
f\mu_0+\sum_{j=1}^{M}
p_j\left(2\mu_0+\frac{\sigma\delta z_j}{16q}\right)}
{J}\\
&=\mu_0+\frac{\sigma\delta}{16q}\frac WJ.
\end{aligned}
\]
% lean: CausalSmith.Stat.MarNearcompleteFrontier.pairedFullLaw_tau_exact

\item
For \(n,d\ge1\), define the dimension factor
\[
\zeta_{n,d}
=\min\left\{1,\frac d{nL}\right\}\in[0,1],
\qquad
g_{n,d,q}=\delta\zeta_{n,d}.
\]
Under the stated separation regime,
\[
\frac1n\le\delta^2\zeta_{n,d}^2
\le\frac14\zeta_{n,d}^2,
\qquad
\frac4n\le\zeta_{n,d}^2.
\]
In particular, \(\zeta_{n,d}>0\). For \(n\ge2\), the regime also implies \(d\ge3\): if \(d\le2\), then
\[
0\le g_{n,d,q}
\le\frac12\frac2{nL}\le\frac1n,
\qquad
g_{n,d,q}^2\le\frac1{n^2}<\frac1n,
\]
a contradiction.

For \(d\ge3\) and \(nL\ge2\), the integer floors satisfy
\[
\left\lfloor\frac{d-1}{2}\right\rfloor\ge\frac d4,
\qquad
\lfloor nL\rfloor\ge\frac{nL}{2}.
\]
Thus
\[
M\ge\frac14\min\{d,nL\}
=\frac{nL}{4}\zeta_{n,d}>0,
\]
and, using \(K\le9L\),
\[
\frac{M}{nK}\ge\frac1{36}\zeta_{n,d},
\qquad
\mu_W\ge
\frac1{36\cdot1024\,C_{\mathrm{hyp}}}\zeta_{n,d}>0.
\]
Squaring the lower bound for \(M\) and using
\(\zeta_{n,d}^2\ge4/n\) also gives
\[
M^2\ge\frac{nL^2}{4},
\qquad
K^4\le6561L^4,
\qquad
K^4\le\frac{26244L^2}{n}M^2.
\]

The logarithmic factor obeys
\[
\forall\varepsilon>0\ \exists N_\varepsilon\in\mathbb N\
\forall n\ge N_\varepsilon:
\qquad L^2\le\varepsilon n.
\]
Indeed, put \(\gamma_{\log}=e+1\). The standard limit
\((\log \xi)^2/\xi\to0\) as \(\xi\to\infty\) gives, eventually,
\[
L^2\le\frac{\varepsilon}{\gamma_{\log}}(e+n)
\le\varepsilon n,
\]
where the last inequality uses \(n\ge1\).

Define
\[
\eta_W=\frac{\beta}{8C_{\mathrm{hyp}}^2}>0.
\]
Choose \(N_W\ge2\), depending only on \(\beta\), so that
\[
n\ge N_W
\quad\Longrightarrow\quad
L^2\le\frac{\eta_W^2}{26244}n.
\]
For every admissible \(n,d,q\) beyond this threshold,
\[
K^4\le\eta_W^2M^2,
\qquad
\frac{K^2}{M}\le\eta_W.
\]
The second inequality follows by taking nonnegative square roots.
% lean: CausalSmith.Stat.MarNearcompleteFrontier.eventually_priorK_sq_div_pairCount_le

\item
Chebyshev's inequality and the positive mean established above give, for admissible \(n\ge N_W\),
\[
\begin{aligned}
\Pr\left(|W-\mu_W|>\frac{\mu_W}{2}\right)
&\le
\frac{M/(1024^2n^2)}{(\mu_W/2)^2}\\
&=\frac{4K^2D^2}{M}
\le4C_{\mathrm{hyp}}^2\frac{K^2}{M}
\le\frac\beta2.
\end{aligned}
\]
Choose another integer threshold \(N_J\), depending only on \(\beta\), such that
\[
n\ge N_J
\quad\Longrightarrow\quad
L^2\le\frac{\beta\,1024^2}{32}n.
\]
The normalizer variance bound, \(M\le nL\), and \(L\ge1\) then yield
\[
\begin{aligned}
\Pr\left(|J-1|>\frac12\right)
&\le\frac{16M}{1024^2n^2}\\
&\le\frac{16L}{1024^2n}
\le\frac{16L^2}{1024^2n}
\le\frac\beta2.
\end{aligned}
\]

Define the common good event, target threshold, and separation constant by
\[
\mathcal G
=\left\{|W-\mu_W|\le\frac{\mu_W}{2},
          \ |J-1|\le\frac12\right\},
\qquad
N_{\mathrm{target}}=\max\{N_W,N_J\},
\]
\[
c_{\mathrm{large}}
=\frac1{36\cdot1024\,C_{\mathrm{hyp}}\cdot48}>0.
\]
For admissible \(n\ge N_{\mathrm{target}}\), the union bound gives
\[
\Pr(\mathcal G)\ge1-\beta.
\]
On \(\mathcal G\),
\[
W\ge\frac{\mu_W}{2},
\qquad
0<J\le\frac32,
\qquad
\frac WJ\ge\frac{\mu_W}{3}.
\]
Since \(\delta\ge0\) and \(0<q\le1\), it follows that
\[
\frac{\delta}{16q}\frac WJ
\ge\frac{\delta\mu_W}{48}
\ge
\frac{\delta\zeta_{n,d}}
     {36\cdot1024\,C_{\mathrm{hyp}}\cdot48}
=c_{\mathrm{large}}g_{n,d,q}.
\]
Set
\[
m_0=\mu_0.
\]
The exact target identity therefore implies, on \(\mathcal G\),
\[
\tau(P_-)\le m_0-c_{\mathrm{large}}g_{n,d,q},
\qquad
m_0+c_{\mathrm{large}}g_{n,d,q}\le\tau(P_+).
\]
Because \(\Pi_\sigma=\operatorname{Law}(P_\sigma)\), the corresponding prior event masses equal the probabilities of their latent preimages. Hence
\[
\Pi_-\{\tau(P)\le m_0-c_{\mathrm{large}}g_{n,d,q}\}
\ge1-\beta,
\]
\[
\Pi_+\{m_0+c_{\mathrm{large}}g_{n,d,q}\le\tau(P)\}
\ge1-\beta.
\]
% lean: CausalSmith.Stat.MarNearcompleteFrontier.pairedPrior_target_events_eventual

\item
Combine the two large-sample thresholds:
\[
N_\beta=\max\{N_{\mathrm{TV}},N_{\mathrm{target}}\}\ge2.
\]
For every admissible \(n\ge N_\beta\), the paired priors simultaneously satisfy
\[
\operatorname{TV}(Q_-^{(n)},Q_+^{(n)})\le\beta
\]
and both target-event inequalities with separation constant
\(c_{\mathrm{large}}\). The threshold and this positive constant are independent of \(n,d,q\).
% lean: CausalSmith.Stat.MarNearcompleteFrontier.normalized_converse

\item
To cover smaller sample sizes, define
\[
a_\beta
=\min\left\{\frac14,
\frac{\sqrt{\log(1+4\beta^2)}}4\right\}>0,
\qquad
u_n=\frac{a_\beta}{\sqrt n},
\qquad n\ge1.
\]
These choices satisfy
\[
0<u_n\le\frac14,
\qquad
16nu_n^2=16a_\beta^2\le\log(1+4\beta^2).
\]

Choose one baseline label, set both potential surrogates and \(Y(0)\) to zero, assign treatment by an independent fair coin, and make arrival complete in both arms. Let
\(P_{-,\mathrm{small}}\) and \(P_{+,\mathrm{small}}\) be the resulting laws with
\[
Y(1)\sim\operatorname{Bernoulli}(1/2)
\quad\text{and}\quad
Y(1)\sim\operatorname{Bernoulli}(1/2+u_n),
\]
respectively. Both laws belong to \(\mathcal M(d,q)\) for every \(d\ge1\) and \(q\in[1/2,1]\), and
\[
\tau(P_{-,\mathrm{small}})=\frac12,
\qquad
\tau(P_{+,\mathrm{small}})=\frac12+u_n.
\]
In the observed-atom order
\((A,R,RY)=(0,1,0),(1,1,0),(1,1,1)\), their probabilities are
\[
P_{-,\mathrm{small},O}:
\left(\frac12,\frac14,\frac14\right),
\qquad
P_{+,\mathrm{small},O}:
\left(\frac12,\frac14-\frac{u_n}{2},
                 \frac14+\frac{u_n}{2}\right).
\]
All other atoms have zero mass. Thus the plus observed law is absolutely continuous with respect to the minus observed law, and direct summation gives
\[
\begin{aligned}
\chi^2(P_{+,\mathrm{small},O}\Vert P_{-,\mathrm{small},O})
&:=
\sum_{\substack{o\in O\\
P_{-,\mathrm{small},O}(\{o\})>0}}
\frac{
\bigl(P_{+,\mathrm{small},O}(\{o\})
-P_{-,\mathrm{small},O}(\{o\})\bigr)^2}
{P_{-,\mathrm{small},O}(\{o\})}\\
&=2u_n^2\le16u_n^2.
\end{aligned}
\]
For finite probability laws, multiplication of likelihood ratios gives the product identity for \(1+\chi^2\), while Cauchy--Schwarz gives
\(\operatorname{TV}\le\frac12\sqrt{\chi^2}\).
Applying these facts to the displayed pair yields
\[
\begin{aligned}
\operatorname{TV}
\left(P_{-,\mathrm{small},O}^{\otimes n},
      P_{+,\mathrm{small},O}^{\otimes n}\right)
&\le
\frac12\sqrt{
\left(1+
\chi^2(P_{+,\mathrm{small},O}\Vert
       P_{-,\mathrm{small},O})\right)^n-1}\\
&\le\frac12\sqrt{e^{16nu_n^2}-1}
\le\beta.
\end{aligned}
\]
Here the middle inequality follows from \(1+t\le e^t\) for \(t\ge0\).

Define
\[
c_{\mathrm{small}}=\frac{a_\beta}{\sqrt{N_\beta}}>0.
\]
For \(1\le n<N_\beta\), the scale bound
\(0\le g_{n,d,q}\le1/2\) gives
\[
c_{\mathrm{small}}g_{n,d,q}
\le\frac{a_\beta}{2\sqrt{N_\beta}}
\le\frac{a_\beta}{2\sqrt n}
=\frac{u_n}{2}.
\]
Take the point priors and midpoint
\[
\Pi_-^{\mathrm{small}}=\delta_{P_{-,\mathrm{small}}},
\qquad
\Pi_+^{\mathrm{small}}=\delta_{P_{+,\mathrm{small}}},
\qquad
m_0=\frac12+\frac{u_n}{2}.
\]
Their target inequalities hold deterministically:
\[
\tau(P_{-,\mathrm{small}})
\le m_0-c_{\mathrm{small}}g_{n,d,q},
\qquad
m_0+c_{\mathrm{small}}g_{n,d,q}
\le\tau(P_{+,\mathrm{small}}).
\]
Both prior event masses are therefore one, and the preceding product bound supplies the required mixture distance.
% lean: CausalSmith.Stat.MarNearcompleteFrontier.normalized_converse_smallN

\item
Finally, set
\[
c_\beta=\min\{c_{\mathrm{large}},c_{\mathrm{small}}\}>0.
\]
If \(n\ge N_\beta\), select the paired priors and their center. If
\(1\le n<N_\beta\), select the point priors and their midpoint.
In either case, write
\[
c_{\mathrm{branch}}
=
\begin{cases}
c_{\mathrm{large}},&n\ge N_\beta,\\
c_{\mathrm{small}},&1\le n<N_\beta.
\end{cases}
\]
Since \(g_{n,d,q}\ge0\) and \(c_\beta\le c_{\mathrm{branch}}\),
\[
\{P:\tau(P)\le m_0-c_{\mathrm{branch}}g_{n,d,q}\}
\subseteq
\{P:\tau(P)\le m_0-c_\beta g_{n,d,q}\},
\]
\[
\{P:m_0+c_{\mathrm{branch}}g_{n,d,q}\le\tau(P)\}
\subseteq
\{P:m_0+c_\beta g_{n,d,q}\le\tau(P)\}.
\]
Monotonicity of prior probabilities preserves both lower bounds
\(1-\beta\), and the mixture-distance bounds remain valid. These two cases cover every admissible positive sample size and establish all the assertions.
% lean: CausalSmith.Stat.MarNearcompleteFrontier.normalized_converse_of_eventual_certificate
\end{enumerate}
\end{proof}

\bibliographystyle{plainnat}

\bibliography{references}

\end{document}
